% Conception d'un modèle simple de la structure interne de la Terre — SVT 1ère D — Cours signé OnBuch+
% Programme officiel MINESEC. Compiler avec : tectonic ND14-modele-structure-interne-terre.tex
\newcommand{\DOCMATIERE}{SVT}
\newcommand{\DOCNIVEAU}{1ère D}
\newcommand{\DOCMODULE}{Module 3 : L'Éducation à l'Environnement et au Développement Durable}
\newcommand{\DOCLECON}{ND14}
\newcommand{\DOCTITRE}{Conception d'un modèle simple de la structure interne de la Terre}
% OnBuch+ — préambule « cours signés » ENRICHI (compilé avec tectonic / XeLaTeX)
% Système visuel « Pop » d'OnBuch+ : fond crème (jamais blanc), bordures encre
% épaisses, ombres « dures » décalées, titres Archivo Black en capitales.
% Version MATHÉMATIQUES : graphes (pgfplots/tikz), boîtes pédagogiques
% (définition, propriété, méthode, attention, le savais-tu, exercice résolu,
% exercice, TP, bilan), page de garde et signature OnBuch+ ancrée en bas.
%
% Macros attendues dans main.tex AVANT \input{preamble} :
%   \DOCMATIERE  \DOCNIVEAU  \DOCMODULE  \DOCLECON (numéro)  \DOCTITRE
\documentclass[11pt]{article}

\usepackage{fontspec}
\usepackage[french]{babel}
\usepackage[a4paper,margin=2cm,top=3.4cm,bottom=2.8cm,headheight=1.4cm,footskip=1.3cm]{geometry}
\usepackage{amsmath,amssymb,amsfonts}
\usepackage{mathtools} % \xmapsto, \coloneqq... souvent utilisés par les modèles
\usepackage{cancel} % simplifications barrées (bilans de forces)
% \abs{z} (module, valeur absolue) et \norm{u} (norme), utilisés par les modèles
\providecommand{\abs}{}\let\abs\relax\DeclarePairedDelimiter{\abs}{\lvert}{\rvert}
\providecommand{\norm}{}\let\norm\relax\DeclarePairedDelimiter{\norm}{\lVert}{\rVert}
\usepackage[table]{xcolor} % [table] : fournit aussi \rowcolors (alternance de lignes)
\usepackage{pagecolor}
\usepackage{tikz}
\usetikzlibrary{arrows.meta,positioning,calc,shapes.geometric,shapes.misc,shapes.arrows,decorations.pathmorphing,decorations.pathreplacing,decorations.markings,patterns,shadows,mindmap,trees,fit,backgrounds,matrix,intersections,angles,quotes}
\usepackage{pgfplots}
\usepackage[version=4]{mhchem} % équations nucléaires \ce{^{235}_{92}U}
\usepackage[european,siunitx]{circuitikz} % schémas électriques
\pgfplotsset{compat=1.18}
\usepgfplotslibrary{fillbetween,groupplots} % aires entre courbes (calcul intégral)
\usepackage{enumitem}
\usepackage{fancyhdr}
% [most] charge toutes les bibliothèques tcolorbox (listings, minted, raster,
% documentation...) et gonfle inutilement la mémoire TeX sur un document long
% (~300 boîtes) : on ne charge que ce qui est réellement utilisé.
\usepackage[skins,breakable]{tcolorbox}
\usepackage{graphicx}
\usepackage{colortbl}
\usepackage{booktabs}
\usepackage{tabularx}
\usepackage{diagbox} % case d'en-tête diagonale des tableaux à double entrée
% Colonnes extensibles centrée / gauche / droite (C, L, R), souvent utilisées par les modèles
\newcolumntype{C}{>{\centering\arraybackslash}X}
\newcolumntype{L}{>{\raggedright\arraybackslash}X}
\newcolumntype{R}{>{\raggedleft\arraybackslash}X}
\usepackage{array}
\usepackage{multicol}
\usepackage{siunitx}
% parse-numbers=false : accepte aussi \SI{2,5 \times 10^{-3}}{...}
\sisetup{locale=FR,output-decimal-marker={,},group-minimum-digits=4,parse-numbers=false}
\DeclareSIUnit\ohm{\ensuremath{\symup{\Omega}}} % U+2126 absent de la police
\DeclareSIUnit\division{div} % réglages d'oscilloscope (ms/div, V/div)
\DeclareSIUnit\tr{tr} % tours (vitesse de rotation, ex. \SI{1500}{\tr\per\minute})
\DeclareSIUnit\ch{ch} % cheval-vapeur (puissance moteur, unité usuelle hors SI)
\DeclareSIUnit\franccfa{FCFA} % franc CFA (coûts, contexte camerounais)
\DeclareSIUnit\dioptre{\ensuremath{\delta}} % vergence d'une lentille (dioptrie)
\usepackage{qrcode}
% \degree : commande fréquemment utilisée par les modèles pour un angle ou une
% température en dehors de \SI{}{} (non fournie par les paquets chargés).
\providecommand{\degree}{^{\circ}}
% \qed (fin de démonstration), utilisé par les modèles sans amsthm
\providecommand{\qed}{\hfill$\square$}
% \barre : double barre des valeurs interdites dans un tableau de variations (tkz-tab)
\providecommand{\barre}{\ensuremath{\|}}
% environnement proof (amsthm non chargé) utilisé par les modèles
\ifdefined\proof\else\newenvironment{proof}[1][Démonstration]{\par\noindent\textit{#1.}\ }{\qed\par}\fi
\usepackage{lastpage}
\usepackage{hyperref}
\hypersetup{hidelinks}

% ============================================================
% Palette « Pop » OnBuch+ — mêmes hex que la classe P (plus_ui.dart)
% ============================================================
\definecolor{poporange}{HTML}{F59321}
\definecolor{popdark}{HTML}{E07A0C}
\definecolor{popcream}{HTML}{FFF6E8}
\definecolor{popink}{HTML}{1C1714}
\definecolor{poppurple}{HTML}{7A5AE0}
\definecolor{popgreen}{HTML}{2E9E6A}
\definecolor{poppink}{HTML}{E0457B}
\definecolor{popblue}{HTML}{2F7DE1}
\definecolor{popgold}{HTML}{C98A12}
\definecolor{popmuted}{HTML}{8A7F76}
\definecolor{popline}{HTML}{EADFD2}
\definecolor{popgray}{HTML}{8A7F76}
\colorlet{popgrey}{popgray}
% Alias tolérants (le modèle nomme parfois les couleurs différemment)
\colorlet{popwhite}{white}
\colorlet{popblack}{popink}
\colorlet{popred}{poppink}
\colorlet{popyellow}{popgold}
% Teintes claires pour fonds de tableaux / zones
\colorlet{popblueL}{popblue!10!white}
\colorlet{poporangeL}{poporange!14!white}
\colorlet{popgreenL}{popgreen!12!white}
\colorlet{poppurpleL}{poppurple!12!white}
\colorlet{poppinkL}{poppink!12!white}
\colorlet{popgoldL}{popgold!14!white}

\pagecolor{popcream}

% ============================================================
% Typographie
% ============================================================
\defaultfontfeatures{Ligatures=TeX}
\setmainfont{PlusJakartaSans-Medium.ttf}[Path=../../../fonts/,BoldFont=PlusJakartaSans-Bold.ttf]
% Maths en sans-serif (Fira Math) avec chiffres et lettres droites de Plus
% Jakarta, pour que les formules se fondent dans le texte.
\usepackage[math-style=ISO,bold-style=ISO]{unicode-math}
\setmathfont{FiraMath-Regular.otf}[Path=../../../fonts/]
\setmathfont{PlusJakartaSans-Medium.ttf}[Path=../../../fonts/,range={up/num,up/{Latin,latin},\mathup}]
\usepackage{icomma} % virgule décimale sans espace en mode math
% Caractères mathématiques Unicode absents des polices de texte (Plus Jakarta,
% Archivo Black) : rendus par leur équivalent mathématique, en texte comme en
% formule — sinon ils s'affichent en carré vide (ex. « Arithmétique dans ℤ »).
\usepackage{newunicodechar}
\newunicodechar{ℝ}{\ensuremath{\mathbb{R}}}
\newunicodechar{ℤ}{\ensuremath{\mathbb{Z}}}
\newunicodechar{ℂ}{\ensuremath{\mathbb{C}}}
\newunicodechar{ℕ}{\ensuremath{\mathbb{N}}}
\newunicodechar{ℚ}{\ensuremath{\mathbb{Q}}}
\newunicodechar{𝔻}{\ensuremath{\mathbb{D}}}
\newunicodechar{α}{\ensuremath{\alpha}}
\newunicodechar{β}{\ensuremath{\beta}}
\newunicodechar{γ}{\ensuremath{\gamma}}
\newunicodechar{δ}{\ensuremath{\delta}}
\newunicodechar{ε}{\ensuremath{\varepsilon}}
\newunicodechar{θ}{\ensuremath{\theta}}
\newunicodechar{λ}{\ensuremath{\lambda}}
\newunicodechar{μ}{\ensuremath{\mu}}
\newunicodechar{σ}{\ensuremath{\sigma}}
\newunicodechar{τ}{\ensuremath{\tau}}
\newunicodechar{φ}{\ensuremath{\varphi}}
\newunicodechar{ψ}{\ensuremath{\psi}}
\newunicodechar{ω}{\ensuremath{\omega}}
\newunicodechar{Σ}{\ensuremath{\Sigma}}
% majuscules produites par \MakeUppercase dans les titres (α-aminés -> Α)
\newunicodechar{Α}{\ensuremath{\alpha}}
\newunicodechar{Β}{\ensuremath{\beta}}
\newunicodechar{Γ}{\ensuremath{\Gamma}}
\newunicodechar{Δ}{\ensuremath{\Delta}}
\newunicodechar{Ε}{\ensuremath{\varepsilon}}
\newunicodechar{Θ}{\ensuremath{\Theta}}
\newunicodechar{Λ}{\ensuremath{\Lambda}}
\newunicodechar{Μ}{\ensuremath{\mu}}
\newunicodechar{Τ}{\ensuremath{\tau}}
\newunicodechar{Φ}{\ensuremath{\Phi}}
\newunicodechar{Ψ}{\ensuremath{\Psi}}
\newunicodechar{Ω}{\ensuremath{\Omega}}
\newunicodechar{↦}{\ensuremath{\mapsto}}
\newunicodechar{∈}{\ensuremath{\in}}
\newunicodechar{∉}{\ensuremath{\notin}}
\newunicodechar{∧}{\ensuremath{\wedge}}
\newunicodechar{⇔}{\ensuremath{\Leftrightarrow}}
\newunicodechar{⟺}{\ensuremath{\Longleftrightarrow}}
\newunicodechar{⇒}{\ensuremath{\Rightarrow}}
\newunicodechar{⟹}{\ensuremath{\Longrightarrow}}
\newunicodechar{∘}{\ensuremath{\circ}}
\newunicodechar{∪}{\ensuremath{\cup}}
\newunicodechar{∩}{\ensuremath{\cap}}
\newunicodechar{≡}{\ensuremath{\equiv}}
\newunicodechar{⊂}{\ensuremath{\subset}}
\newunicodechar{⊥}{\ensuremath{\perp}}
\newunicodechar{∃}{\ensuremath{\exists}}
\newunicodechar{∀}{\ensuremath{\forall}}
\newunicodechar{∛}{\ensuremath{\sqrt[3]{\,}}}
\newunicodechar{∜}{\ensuremath{\sqrt[4]{\,}}}
\newunicodechar{⃗}{\ensuremath{{}^{\to}}}
\newunicodechar{ⁿ}{\ifmmode^{n}\else\textsuperscript{\ensuremath{n}}\fi}
\newunicodechar{ˣ}{\ifmmode^{x}\else\textsuperscript{\ensuremath{x}}\fi}
\newunicodechar{ᵇ}{\ifmmode^{b}\else\textsuperscript{\ensuremath{b}}\fi}
\newunicodechar{ʳ}{\ifmmode^{r}\else\textsuperscript{\ensuremath{r}}\fi}
\newunicodechar{ᵘ}{\ifmmode^{u}\else\textsuperscript{\ensuremath{u}}\fi}
\newunicodechar{ʸ}{\ifmmode^{y}\else\textsuperscript{\ensuremath{y}}\fi}
\newunicodechar{ᵃ}{\ifmmode^{a}\else\textsuperscript{\ensuremath{a}}\fi}
\newunicodechar{ᵉ}{\ifmmode^{e}\else\textsuperscript{\ensuremath{e}}\fi}
\newunicodechar{ᵏ}{\ifmmode^{k}\else\textsuperscript{\ensuremath{k}}\fi}
\newunicodechar{ᵖ}{\ifmmode^{p}\else\textsuperscript{\ensuremath{p}}\fi}
\newunicodechar{ᴺ}{\ifmmode^{N}\else\textsuperscript{\ensuremath{N}}\fi}
\newunicodechar{ᶜ}{\ifmmode^{c}\else\textsuperscript{\ensuremath{c}}\fi}
\newunicodechar{ʰ}{\ifmmode^{h}\else\textsuperscript{\ensuremath{h}}\fi}
\newunicodechar{ᵠ}{\ifmmode^{q}\else\textsuperscript{\ensuremath{q}}\fi}
\newunicodechar{ₙ}{\ifmmode_{n}\else\textsubscript{\ensuremath{n}}\fi}
\newunicodechar{ₐ}{\ifmmode_{a}\else\textsubscript{\ensuremath{a}}\fi}
\newunicodechar{ᵢ}{\ifmmode_{i}\else\textsubscript{\ensuremath{i}}\fi}
\newunicodechar{ᵣ}{\ifmmode_{r}\else\textsubscript{\ensuremath{r}}\fi}
\newunicodechar{ₖ}{\ifmmode_{k}\else\textsubscript{\ensuremath{k}}\fi}
\newunicodechar{ᵦ}{\ifmmode_{\beta}\else\textsubscript{\ensuremath{\beta}}\fi}
\newunicodechar{ₓ}{\ifmmode_{x}\else\textsubscript{\ensuremath{x}}\fi}
\newfontfamily\popdisplay{ArchivoBlack-Regular.ttf}[Path=../../../fonts/]
\newfontfamily\poplogo{SpaceGrotesk-Bold.ttf}[Path=../../../fonts/]
\newfontfamily\popsemi{PlusJakartaSans-SemiBold.ttf}[Path=../../../fonts/]
\newfontfamily\popxbold{PlusJakartaSans-ExtraBold.ttf}[Path=../../../fonts/]
\newfontfamily\popxboldls{PlusJakartaSans-ExtraBold.ttf}[Path=../../../fonts/,LetterSpace=3.0]

\setlist[enumerate]{leftmargin=1.7em,itemsep=5pt,topsep=4pt}
\setlist[itemize]{leftmargin=1.7em,itemsep=4pt,topsep=4pt,label={\color{poporange}\textbullet}}
\setlength{\parindent}{0pt}
\setlength{\parskip}{5pt}
\renewcommand{\arraystretch}{1.25}

% pgfplots : style Pop par défaut
\pgfplotsset{
  every axis/.append style={axis line style={popink,line width=1pt},tick style={popink},
    grid style={popline,line width=0.5pt},label style={font=\small},tick label style={font=\footnotesize}},
  popcurve/.style={poporange,line width=1.8pt,no markers,smooth},
}
% Même style disponible dans un \draw TikZ simple (hors pgfplots)
\tikzset{popcurve/.style={poporange,line width=1.8pt}}
\tikzset{popgrid/.style={popline,very thin}}
\colorlet{popcurve}{poporange} % le nom du style est parfois utilisé comme couleur

% ============================================================
% Pill / squiggle
% ============================================================
\newcommand{\poppill}[3]{%
  \tikz[baseline=-0.55ex]\node[draw=popink,line width=1.2pt,rounded corners=2.6pt,%
    fill=#2,inner xsep=6.5pt,inner ysep=3pt]{%
    \popxboldls\fontsize{7}{7}\selectfont\color{#3}\MakeUppercase{#1}};%
}
\newcommand{\popsquiggle}[1]{%
  \tikz[baseline=-0.4ex]{
    \draw[#1,line width=2.6pt,line cap=round]
      plot [smooth] coordinates {(0,0.09) (0.34,0.01) (0.68,0.21) (1.02,0.01) (1.36,0.10)};
  }%
}
% Mot-clé surligné (marqueur orange sous le texte)
\newcommand{\cle}[1]{{\popxbold\color{popdark}#1}}

% ============================================================
% En-tête / pied
% ============================================================
\pagestyle{fancy}
\fancyhf{}
\renewcommand{\headrulewidth}{0pt}
\renewcommand{\footrulewidth}{0pt}
\lhead{\raisebox{-0.15ex}{\poplogo\fontsize{13}{13}\selectfont\color{popink}OnBuch\color{poporange}+}}
\rhead{\poppill{\DOCMATIERE\ \textbullet\ \DOCNIVEAU\ \textbullet\ Leçon \DOCLECON}{white}{popink}}
\lfoot{\popxbold\fontsize{7.6}{7.6}\selectfont\color{popmuted}\MakeUppercase{Cours signé OnBuch+}\ \textbullet\ {\popsemi\DOCTITRE}}
\rfoot{\tikz[baseline=-0.6ex]\node[rounded corners=8.5pt,fill=popink,minimum height=17pt,minimum width=17pt,inner xsep=5pt,inner sep=0]{\popxbold\fontsize{8}{8}\selectfont\color{popcream}\thepage\,/\,\pageref*{LastPage}};}

% ============================================================
% Titres
% ============================================================
\newcommand{\chaptitre}[2]{%
  \begin{tcolorbox}[enhanced,colback=poporange,colframe=popink,boxrule=2.6pt,arc=11pt,
      drop shadow=popink,left=15pt,right=15pt,top=12pt,bottom=11pt,before skip=3pt,after skip=18pt]
    \poppill{#2}{popink}{popcream}\par\vspace{9pt}
    {\raggedright\hyphenpenalty=10000\exhyphenpenalty=10000\popdisplay\fontsize{22}{24}\selectfont\color{popcream}\MakeUppercase{#1}\par}
  \end{tcolorbox}
}
% Section numérotée : pastille orange avec le numéro + titre Archivo
\newcounter{popsec}
\newcommand{\coursec}[1]{%
  \stepcounter{popsec}\setcounter{popsub}{0}%
  \par\vspace{16pt}\noindent
  \begin{minipage}{\linewidth}
  \tikz[baseline=-0.7ex]\node[draw=popink,line width=1.6pt,fill=poporange,rounded corners=4pt,minimum size=22pt,inner sep=0]{\popdisplay\fontsize{12}{12}\selectfont\color{popcream}\thepopsec};%
  \hspace{8pt}{\popdisplay\fontsize{15}{17}\selectfont\color{popink}\MakeUppercase{#1}}\par
  \vspace{4pt}\noindent{\color{poporange}\rule{36pt}{3.6pt}}
  \end{minipage}\par\nopagebreak\vspace{8pt}
}
\newcounter{popsub}
\newcommand{\courssub}[1]{%
  \stepcounter{popsub}%
  \par\vspace{9pt}\noindent{\popxbold\fontsize{11.5}{13}\selectfont\color{popink}{\color{poporange}\thepopsec.\thepopsub}\hspace{6pt}#1}\par\nopagebreak\vspace{4pt}
}

% ============================================================
% Boîtes pédagogiques — même recette : fond blanc, bordure épaisse colorée,
% ombre dure encre, badge plein. Titre optionnel [#1].
% ============================================================
\tcbset{popbase/.style={breakable,enhanced,colback=white,boxrule=2.3pt,arc=9pt,
  shadow={3pt}{-3pt}{0pt}{popink},left=14pt,right=14pt,top=11pt,bottom=11pt,before skip=11pt,after skip=15pt,
  fontupper=\popsemi\fontsize{10.6}{14.5}\selectfont\color{popink},
  fontlower=\popsemi\fontsize{10.6}{14.5}\selectfont\color{popink}}}
% Coche dessinée (le glyphe ✓ n'existe pas dans les polices du document)
\newcommand{\popcheck}[1]{\tikz[baseline=-0.5ex]\draw[#1,line width=1.6pt,line cap=round,line join=round](-0.13,0.0)--(-0.04,-0.09)--(0.13,0.1);}
\providecommand{\coche}{\popcheck{popgreen}}
\providecommand{\resultat}[1]{\par\smallskip\noindent\fcolorbox{popgreen}{popgreenL}{\parbox{\dimexpr\linewidth-2\fboxsep-2\fboxrule\relax}{\textbf{Résultat.} #1}}\par\smallskip}
\newcommand{\popboxhead}[3]{\poppill{#1}{#2}{white}\ifx\relax#3\relax\else\hspace{6pt}{\popxbold\fontsize{10.6}{12}\selectfont\color{popink}#3}\fi\par\vspace{7pt}}

\newtcolorbox{aretenir}[1][]{popbase,colframe=popgreen,before upper={\popboxhead{À retenir}{popgreen}{#1}}}
\newtcolorbox{exemplebox}[1][]{popbase,colframe=poporange,before upper={\popboxhead{Exemple}{poporange}{#1}}}
\newtcolorbox{definition}[1][]{popbase,colframe=popblue,before upper={\popboxhead{Définition}{popblue}{#1}}}
\newtcolorbox{propriete}[1][]{popbase,colframe=poppurple,before upper={\popboxhead{Propriété}{poppurple}{#1}}}
\newtcolorbox{methode}[1][]{popbase,colframe=popgold,before upper={\popboxhead{Méthode}{popgold}{#1}}}
\newtcolorbox{attention}[1][]{popbase,colframe=poppink,before upper={\popboxhead{Attention}{poppink}{#1}}}
\newtcolorbox{experience}[1][]{popbase,colframe=popblue,colback=popblueL,before upper={\popboxhead{Expérience}{popblue}{#1}}}
% « Le savais-tu ? » : carte violette pleine (contexte, culture, Cameroun)
\newtcolorbox{savaistu}[1][]{popbase,colback=poppurple,colframe=popink,
  fontupper=\popsemi\fontsize{10.4}{14.2}\selectfont\color{white},
  before upper={\poppill{Le savais-tu\,?}{white}{poppurple}\ifx\relax#1\relax\else\hspace{6pt}{\popxbold\color{white}#1}\fi\par\vspace{7pt}}}
% Exercice résolu : énoncé en haut, \tcblower puis la solution sur fond crème
\newcounter{popexo}\newcounter{popexr}
\newtcolorbox{exoresolu}[1][]{popbase,colframe=poporange,colbacklower=popcream,
  segmentation style={popink,line width=1.2pt,dashed},
  before upper={\stepcounter{popexr}\popboxhead{Exercice résolu \thepopexr}{poporange}{#1}},
  before lower={\poppill{Solution}{popink}{popcream}\par\vspace{6pt}}}
% Exercice (non résolu en place) — la correction est dans la partie Corrigés
\newtcolorbox{exercice}[1][]{popbase,colframe=popink,
  before upper={\stepcounter{popexo}\popboxhead{Exercice \thepopexo}{popink}{#1}}}
% Corrigé
\newtcolorbox{corrige}[1][]{popbase,colframe=popgreen,colback=popgreenL,
  before upper={\popboxhead{Corrigé}{popgreen}{#1}}}
% Objectifs / prérequis / bilan
\newtcolorbox{objectifs}{popbase,colframe=popink,colback=poporangeL,
  before upper={\popboxhead{Objectifs de la leçon}{popink}{}}}
\newtcolorbox{prerequis}{popbase,colframe=popink,colback=popblueL,
  before upper={\popboxhead{Prérequis}{popblue}{}}}
\newtcolorbox{bilan}{popbase,colframe=popink,colback=white,boxrule=2.8pt,
  before upper={\popboxhead{Fiche bilan}{poporange}{L'essentiel à connaître pour l'examen}}}
% Figure encadrée avec légende
\newtcolorbox{popfigure}[1][]{enhanced,breakable=false,colback=white,colframe=popline,boxrule=1.4pt,arc=8pt,
  left=10pt,right=10pt,top=10pt,bottom=8pt,before skip=10pt,after skip=12pt,halign=center,
  lower separated=false,halign lower=center,
  fontlower=\popsemi\fontsize{9}{11.5}\selectfont\color{popmuted},
  #1}
\newcommand{\legende}[1]{\par\vspace{6pt}{\popsemi\fontsize{9}{11.5}\selectfont\color{popmuted}#1}}

% ============================================================
% Signature OnBuch+
% ============================================================
\newcommand{\poplogoinline}[1]{{\poplogo\fontsize{#1}{#1}\selectfont\color{popink}OnBuch\color{poporange}+}}
\newcommand{\signatureonbuch}{%
  \begin{tcolorbox}[enhanced,colback=popink,colframe=popink,boxrule=2.2pt,arc=10pt,
      drop shadow=poporange,left=14pt,right=14pt,top=10pt,bottom=10pt,before skip=0pt,after skip=0pt]
    \tikz[baseline=-0.7ex]\node[circle,fill=popgreen,draw=popcream,line width=1.3pt,minimum size=22pt,inner sep=0]{\popcheck{white}};%
    \hspace{9pt}\begin{minipage}[c]{0.62\linewidth}
      {\popxbold\fontsize{12}{13}\selectfont\color{popcream}Signé\ \textbullet\ {\poplogo OnBuch\color{poporange}+}}\par\vspace{2pt}
      {\raggedright\popsemi\fontsize{8}{10.5}\selectfont\color{popline}\MakeUppercase{Équipe pédagogique OnBuch+}\\\MakeUppercase{Conforme au programme officiel MINESEC}\par}
    \end{minipage}\hfill
    {\poplogo\fontsize{22}{22}\selectfont\color{popcream}OnBuch\color{poporange}+}
  \end{tcolorbox}%
}

% ============================================================
% Page de garde
% #1 titre · #2 matière · #3 niveau · #4 module · #5 n° leçon · #6 durée ·
% #7 description / objectifs (texte ou itemize)
% ============================================================
\newcommand{\pagegarde}[7]{%
  \thispagestyle{empty}
  \newgeometry{a4paper,margin=1.7cm,top=1.6cm,bottom=1.4cm}%
  \begin{tikzpicture}[remember picture,overlay]
    \fill[poporange!18] ([xshift=-3.2cm,yshift=-3.6cm]current page.north east) circle (4.6cm);
    \fill[poppurple!14] ([xshift=2.4cm,yshift=7.6cm]current page.south west) circle (3.8cm);
  \end{tikzpicture}%
  {\poplogo\fontsize{30}{30}\selectfont\color{popink}OnBuch\color{poporange}+}\hfill
  \raisebox{0.6ex}{\poppill{Cours signé \textbullet\ Premium}{popink}{popcream}}\par
  \vspace{1pt}\popsquiggle{poppurple}\par
  \vspace{8pt}
  \begin{tcolorbox}[enhanced,colback=poporange,colframe=popink,boxrule=2.8pt,arc=13pt,
      drop shadow=popink,left=18pt,right=18pt,top=12pt,bottom=13pt,after skip=10pt]
    \poppill{#2 \textbullet\ #3}{popink}{popcream}\hspace{5pt}\poppill{#4}{white}{popink}\par\vspace{8pt}
    {\popdisplay\fontsize{14}{14}\selectfont\color{popink}LEÇON #5}\par\vspace{4pt}
    {\raggedright\hyphenpenalty=10000\exhyphenpenalty=10000\popdisplay\fontsize{25}{27}\selectfont\color{popcream}\MakeUppercase{#1}\par}
    \vspace{8pt}
    {\popxbold\fontsize{9.5}{9.5}\selectfont\color{popink}\MakeUppercase{Durée conseillée : #6}}
  \end{tcolorbox}
  \begin{tcolorbox}[enhanced,colback=white,colframe=popink,boxrule=2.2pt,arc=10pt,
      drop shadow=popink,left=16pt,right=16pt,top=10pt,bottom=10pt,after skip=8pt]
    \poppill{Dans cette leçon}{poppurple}{white}\par\vspace{5pt}
    \popsemi\fontsize{9.3}{12.6}\selectfont\color{popink}#7
  \end{tcolorbox}
  \noindent\begin{minipage}[t]{0.487\linewidth}
    \begin{tcolorbox}[enhanced,colback=popgreen,colframe=popink,boxrule=2.2pt,arc=10pt,
        drop shadow=popink,left=11pt,right=11pt,top=10pt,bottom=10pt,height=3.1cm,valign=center]
      \tcbox[colback=white,colframe=popink,boxrule=1.3pt,arc=5pt,boxsep=0pt,nobeforeafter,
        left=3pt,right=3pt,top=3pt,bottom=3pt,valign=center]{\qrcode[height=1.9cm]{https://play.google.com/store/apps/details?id=com.luvvix.onbuch&hl=fr}}%
      \hspace{8pt}\begin{minipage}[c]{0.5\linewidth}
        {\popdisplay\fontsize{11}{12.5}\selectfont\color{white}\MakeUppercase{Télécharge OnBuch}}\par\vspace{4pt}
        {\raggedright\popsemi\fontsize{8.4}{11}\selectfont\color{white}Gratuit \textbullet\ Google Play\\Tuteur IA, annales, concours, résultats\par}
      \end{minipage}
    \end{tcolorbox}
  \end{minipage}\hfill
  \begin{minipage}[t]{0.487\linewidth}
    \begin{tcolorbox}[enhanced,colback=poppurple,colframe=popink,boxrule=2.2pt,arc=10pt,
        drop shadow=popink,left=11pt,right=11pt,top=10pt,bottom=10pt,height=3.1cm,valign=center]
      \tikz[baseline=-0.7ex]\node[circle,fill=white,draw=popink,line width=1.3pt,minimum size=1.6cm,inner sep=0]{\popdisplay\fontsize{24}{24}\selectfont\color{poppurple}+};%
      \hspace{8pt}\begin{minipage}[c]{0.6\linewidth}
        {\popdisplay\fontsize{11}{12.5}\selectfont\color{white}\MakeUppercase{Passe à OnBuch+}}\par\vspace{4pt}
        {\raggedright\popsemi\fontsize{8.4}{11}\selectfont\color{white}Tous les cours signés, Tuteur IA illimité, fiches PDF — onglet OnBuch+ de l'app\par}
      \end{minipage}
    \end{tcolorbox}
  \end{minipage}
  \vspace{8pt}
  \popcontenu
  \vfill
  \signatureonbuch
  \restoregeometry
  \newpage
}

% Rangée « Ce cours contient » (page de garde)
\newcommand{\poptile}[3]{% #1 couleur #2 gros texte #3 légende
  \begin{tcolorbox}[enhanced,colback=#1,colframe=popink,boxrule=2pt,arc=9pt,shadow={2.5pt}{-2.5pt}{0pt}{popink},
      left=6pt,right=6pt,top=9pt,bottom=9pt,halign=center,valign=center,height=1.75cm,width=\linewidth,nobeforeafter]
    {\popdisplay\fontsize{12.5}{13}\selectfont\color{white}\MakeUppercase{#2}}\par\vspace{3pt}
    {\centering\popsemi\fontsize{7.8}{9.5}\selectfont\color{white}#3\par}
  \end{tcolorbox}}
\newcommand{\popcontenu}{%
  \par\noindent{\popxboldls\fontsize{7.5}{7.5}\selectfont\color{popmuted}CE COURS SIGNÉ CONTIENT}\par\vspace{7pt}
  \noindent\begin{minipage}[t]{0.235\linewidth}\poptile{popblue}{Cours}{complet, illustré, pas à pas}\end{minipage}\hfill
  \begin{minipage}[t]{0.235\linewidth}\poptile{poporange}{Méthodes}{et exercices résolus}\end{minipage}\hfill
  \begin{minipage}[t]{0.235\linewidth}\poptile{poppink}{Exercices}{type examen + corrigés}\end{minipage}\hfill
  \begin{minipage}[t]{0.235\linewidth}\poptile{popgreen}{Fiche bilan}{l'essentiel à retenir}\end{minipage}\par}

% Sommaire
\newtcolorbox{sommaire}{enhanced,colback=white,colframe=popink,boxrule=2.2pt,arc=10pt,
  drop shadow=popink,left=18pt,right=18pt,top=13pt,bottom=13pt,before skip=6pt,after skip=20pt,
  before upper={\poppill{Sommaire}{poppurple}{white}\par\vspace{9pt}\popsemi\fontsize{10.6}{14.5}\selectfont\color{popink}}}

% Signature de fin de cours — poussée tout en bas de la dernière page
\newcommand{\findecours}{%
  \par\vfill\nopagebreak
  \begin{center}\popsquiggle{poporange}\end{center}\vspace{4pt}
  \signatureonbuch
}
\AtBeginDocument{\def\llbracket{\mathopen{[\mkern-3.5mu[}}\def\rrbracket{\mathclose{]\mkern-3.5mu]}}} % intervalles d'entiers (glyphe ⟦ absent de la police)
% Glyphes absents de Fira Math : redéfinis à partir de glyphes disponibles
\AtBeginDocument{%
  \def\setminus{\mathbin{\backslash}}\let\smallsetminus\setminus
  \def\vdots{\mathord{\vbox{\baselineskip3.2pt\lineskiplimit0pt\kern4pt\hbox{.}\hbox{.}\hbox{.}}}}%
  \def\ddots{\mathinner{\mkern1mu\raise6.5pt\hbox{.}\mkern2mu\raise3.6pt\hbox{.}\mkern2mu\raise0.7pt\hbox{.}\mkern1mu}}%
  \def\triangle{\mathord{\scalebox{0.9}{$\symup{\Delta}$}}}%
  \def\nabla{\mathord{\raisebox{\depth}{\scalebox{1}[-1]{$\symup{\Delta}$}}}}%
  \def\checkmark{\popcheck{popdark}}%
  \def\star{\mathbin{\ast}}\def\bigstar{\mathord{\ast}}%
  \def\diamond{\mathbin{\scalebox{0.6}{$\Diamond$}}}%
  \def\bigcup{\mathop{\mathchoice{\raisebox{-0.3em}{\scalebox{1.7}{$\cup$}}}{\scalebox{1.25}{$\cup$}}{\cup}{\cup}}\displaylimits}%
  \def\bigcap{\mathop{\mathchoice{\raisebox{-0.3em}{\scalebox{1.7}{$\cap$}}}{\scalebox{1.25}{$\cap$}}{\cap}{\cap}}\displaylimits}%
  \def\lesseqqgtr{\mathrel{\lessgtr}}\def\bigoplus{\mathop{\oplus}\displaylimits}%
}
\providecommand{\ssi}{si et seulement si} % abréviation fréquente
\AtBeginDocument{\providecommand{\wideparen}{\overparen}} % arc de cercle

%% FIGFIX-BEGIN v5 — correctifs globaux des figures (docs/onbuch-generation/tools/figfix)
\usepackage{adjustbox}\usepackage{etoolbox}\tcbuselibrary{raster}
% toute figure TikZ/pgfplots plus large que la ligne est réduite à la largeur disponible
\BeforeBeginEnvironment{tikzpicture}{\begin{adjustbox}{max width=\linewidth}}
\AfterEndEnvironment{tikzpicture}{\end{adjustbox}}
% légendes pgfplots lisibles par-dessus les courbes
\pgfplotsset{every axis/.append style={legend style={fill=white,fill opacity=0.92,text opacity=1,draw=black!15,inner sep=3pt}}}
% étiquettes d'arêtes (syntaxe edge["..."]) avec fond blanc
\tikzset{every edge quotes/.append style={fill=white,inner sep=1.2pt}}
%% FIGFIX-END
\begin{document}

\pagegarde{Conception d'un modèle simple de la structure interne de la Terre}{SVT}{1ère D}{Module 3 : L'Éducation à l'Environnement et au Développement Durable}{ND14}{14 h}{Cette leçon montre comment l'étude de la propagation des ondes sismiques (P et S) permet de déterminer la structure interne du globe terrestre. L'élève apprend à exploiter des temps de parcours d'ondes pour localiser les discontinuités et à construire un modèle simplifié en couches : croûte, manteau et noyau.
\vspace{4pt}
\begin{multicols}{2}\small
\begin{itemize}
\item À la fin de cette leçon, je sais définir les notions de foyer, d'épicentre et d'onde sismique.
\item À la fin de cette leçon, je sais distinguer les ondes P et les ondes S par leur mode de propagation, leur vitesse et les milieux qu'elles traversent.
\item À la fin de cette leçon, je sais expliquer le comportement des ondes (réflexion, réfraction, arrêt) au niveau d'une discontinuité.
\item À la fin de cette leçon, je sais exploiter des données de temps de parcours d'ondes sismiques pour situer une discontinuité.
\item À la fin de cette leçon, je sais interpréter l'existence des zones d'ombre pour démontrer la présence d'un noyau liquide.
\item À la fin de cette leçon, je sais construire un modèle simplifié en couches de la structure interne du globe (croûte, manteau, noyau).
\end{itemize}\end{multicols}}

\begin{sommaire}
\begin{multicols}{2}
\begin{enumerate}
\item Séismes et ondes sismiques
\item Comportement des ondes aux discontinuités
\item Exploitation des temps de parcours : la démarche de la sismologie
\item Les grandes discontinuités et les enveloppes de la Terre
\item Construction du modèle en couches et limites du modèle
\item Activité d'intégration
\item Méthodes et exercices résolus
\item Exercices
\item Corrigés des exercices
\item Fiche bilan
\end{enumerate}
\end{multicols}
\end{sommaire}

\begin{objectifs}
\begin{itemize}
\item À la fin de cette leçon, je sais définir les notions de foyer, d'épicentre et d'onde sismique.
\item À la fin de cette leçon, je sais distinguer les ondes P et les ondes S par leur mode de propagation, leur vitesse et les milieux qu'elles traversent.
\item À la fin de cette leçon, je sais expliquer le comportement des ondes (réflexion, réfraction, arrêt) au niveau d'une discontinuité.
\item À la fin de cette leçon, je sais exploiter des données de temps de parcours d'ondes sismiques pour situer une discontinuité.
\item À la fin de cette leçon, je sais interpréter l'existence des zones d'ombre pour démontrer la présence d'un noyau liquide.
\item À la fin de cette leçon, je sais construire un modèle simplifié en couches de la structure interne du globe (croûte, manteau, noyau).
\item À la fin de cette leçon, je sais relier les propriétés des ondes sismiques aux propriétés physiques (état solide ou liquide, densité) des couches traversées.
\item À la fin de cette leçon, je sais citer les caractéristiques (état, épaisseur approximative, composition) des trois enveloppes de la Terre.
\end{itemize}
\end{objectifs}

\begin{prerequis}
\begin{itemize}
\item Notion d'onde mécanique progressive et de vitesse de propagation (Physique, 2nde/1ère)
\item Relation fondamentale v = d/t et calculs de durées et de distances
\item Notion de réfraction et de changement de milieu de propagation (Physique)
\item Structure externe du globe : lithosphère, continents, océans (géographie du secondaire)
\item Lecture et exploitation de graphiques et de tableaux de données
\end{itemize}
\end{prerequis}

\newpage

\coursec{Séismes et ondes sismiques}

En février 2023, un violent séisme frappe la Turquie et la Syrie, faisant des dizaines de milliers de victimes. Quelques minutes après le choc, des stations sismiques situées jusqu'en Afrique enregistrent un léger tremblement : la Terre entière a « vibré ». Comment des vibrations nées à des kilomètres de profondeur peuvent-elles traverser toute la planète et être détectées à Yaoundé ? Mieux : comment les scientifiques ont-ils pu découvrir que la Terre possède un noyau métallique, alors qu'aucun forage n'a jamais dépassé \SI{12}{km} de profondeur (le forage de Kola, en Russie) ? La réponse tient en une image : c'est en « écoutant » les ondes des séismes que l'Homme a appris à voir l'intérieur du globe, comme un médecin qui fait une échographie du corps humain. Cette leçon va nous apprendre, pas à pas, à construire ce modèle en couches de la Terre. Mais avant de sonder la planète, il faut comprendre ce qu'est un séisme et quelles sont les ondes qu'il émet : c'est l'objet de cette première section.

\courssub{Qu'est-ce qu'un séisme ?}

\textbf{Une observation de la vie courante.}\par

Pose une règle en plastique à moitié sur le bord d'une table, maintiens-la fermement, puis plie l'extrémité libre et lâche-la d'un coup : la règle vibre brutalement en libérant l'énergie accumulée pendant la flexion. La croûte terrestre se comporte de la même façon : les roches, soumises à des contraintes lentes mais continues (mouvements des plaques lithosphériques), se déforment élastiquement et accumulent de l'énergie. Lorsque la contrainte dépasse la résistance des roches, celles-ci \emph{rompent} brusquement le long d'une faille : l'énergie accumulée est libérée en quelques secondes sous forme de vibrations qui se propagent dans toutes les directions. C'est un séisme.

\begin{definition}[Séisme]
Un \cle{séisme} (ou tremblement de terre) est une vibration brutale et brève du sol, provoquée par la libération soudaine, en profondeur, de l'énergie élastique accumulée par les roches soumises à des contraintes. Cette énergie se propage dans le globe sous forme d'\cle{ondes sismiques}.
\end{definition}

\begin{definition}[Foyer et épicentre]
Le \cle{foyer} (ou \cle{hypocentre}) est le point situé \emph{en profondeur} où débute la rupture des roches : c'est le point de départ du séisme. L'\cle{épicentre} est le point de la \emph{surface} du globe situé à la verticale du foyer. C'est généralement à l'épicentre que les dégâts sont les plus importants.
\end{definition}

La profondeur du foyer varie de quelques kilomètres (séismes superficiels, les plus destructeurs) à près de \SI{700}{km} (séismes profonds, dans les zones de subduction). La distance entre le foyer et l'épicentre s'appelle la \emph{profondeur focale}.

\begin{attention}[Ne confonds jamais foyer et épicentre !]
Erreur très fréquente au Probatoire : le \textbf{foyer} est \emph{en profondeur} (c'est là que naît le séisme), l'\textbf{épicentre} est \emph{en surface} (à la verticale du foyer). Une phrase comme « le séisme a son épicentre à \SI{10}{km} de profondeur » est fausse : c'est le \emph{foyer} qui est à \SI{10}{km} de profondeur. Retiens : \textbf{F}oyer = \textbf{F}ond ; \textbf{É}picentre = \textbf{É}piderme (la surface).
\end{attention}

\begin{popfigure}
\begin{center}
\begin{tikzpicture}[scale=0.9]
% Sol
\fill[popgreenL] (-6,0) rectangle (6,0.5);
\draw[thick, popdark] (-6,0.5) -- (6,0.5);
\node[popdark] at (-4.8,0.25) {\small Surface du sol};
% Faille
\draw[very thick, poporange] (0.4,-4.6) -- (-0.4,-1.2);
\node[poporange, rotate=68] at (1.35,-2.9) {\small faille};
% Foyer
\fill[popink] (0,-3.5) circle (0.12);
\node[popink, right] at (0.2,-3.7) {\small \textbf{Foyer} (hypocentre)};
% Epicentre
\fill[poppurple] (0,0.5) circle (0.12);
\node[poppurple, above] at (0,0.62) {\small \textbf{Épicentre}};
% Verticale foyer-épicentre
\draw[dashed, popmuted] (0,-3.5) -- (0,0.5);
\draw[<->, popmuted] (0.55,-3.5) -- (0.55,0.5) node[midway, right]{\scriptsize profondeur focale};
% Ondes rayonnantes
\draw[popblue, thick] (0,-3.5) circle (0.7);
\draw[popblue] (0,-3.5) circle (1.3);
\draw[popblue!60] (0,-3.5) circle (1.9);
\draw[popblue!35] (0,-3.5) circle (2.5);
\node[popblue] at (-2.6,-2.1) {\small ondes sismiques};
\draw[->, popblue] (-2.2,-2.35) -- (-1.5,-2.75);
% Maisons
\draw[popdark] (2.6,0.5) rectangle (3.4,1.1);
\draw[popdark] (2.5,1.1) -- (3.0,1.5) -- (3.5,1.1);
\draw[popdark] (-3.6,0.5) rectangle (-2.8,1.1);
\draw[popdark] (-3.7,1.1) -- (-3.2,1.5) -- (-2.7,1.1);
\end{tikzpicture}
\end{center}
\legende{Coupe schématique d'un séisme : la rupture naît au foyer, en profondeur, le long d'une faille ; l'épicentre est le point de la surface à la verticale du foyer ; les ondes sismiques rayonnent depuis le foyer dans toutes les directions.}
\end{popfigure}

\begin{savaistu}[Le Cameroun, une terre sismiquement active]
Le Cameroun possède des zones sismiques actives le long de la \emph{ligne volcanique du Cameroun}, qui s'étend du golfe de Guinée (île de Bioko) au mont Cameroun (\SI{4095}{m}) et jusqu'à la région de l'Ouest (monts Bamboutos, lac Manengouba). Le mont Cameroun, volcan encore actif (éruptions en 1999 et 2000), est surveillé car l'activité magmatique s'accompagne de séismes. La région de l'Ouest a connu des séismes destructeurs, et les catastrophes des lacs Nyos (1986) et Monoun (1984), bien que d'origine gazeuse, rappellent que cette zone est géologiquement vivante. Des stations sismologiques, notamment à l'IRGM (Institut de recherches géologiques et minières), enregistrent l'activité sismique du pays.
\end{savaistu}

\courssub{Les ondes sismiques : des messagères de l'intérieur du globe}

\textbf{Intuition.}\par

Quand tu jettes un caillou dans l'eau d'une mare, des rides circulaires s'éloignent du point d'impact : l'énergie du choc se propage sans que l'eau elle-même ne se déplace durablement. De même, l'énergie libérée au foyer d'un séisme se propage à travers les roches sous forme d'ondes : ce sont les \cle{ondes sismiques}.

\begin{definition}[Onde sismique]
Une \cle{onde sismique} est une onde élastique qui transporte, à travers les roches du globe, l'énergie libérée au foyer d'un séisme. Elle correspond à une déformation temporaire du milieu traversé : les roches vibrent puis retrouvent leur état initial après le passage de l'onde.
\end{definition}

On distingue deux grandes familles d'ondes sismiques :

\begin{itemize}
\item les \cle{ondes de volume}, qui traversent l'intérieur du globe dans toutes les directions : ce sont les ondes \cle{P} et les ondes \cle{S}. Ce sont elles qui nous renseignent sur la structure interne de la Terre ;
\item les \cle{ondes de surface} (ondes \cle{L} de Love et \cle{R} de Rayleigh), qui ne se propagent qu'à la surface du globe, comme les rides à la surface de l'eau. Plus lentes, ce sont elles qui provoquent l'essentiel des dégâts des constructions. Nous ne les étudierons pas en détail.
\end{itemize}

\courssub{Les ondes P : compression-dilatation}

\textbf{Observation.}\par

Prends un ressort, pose-le sur une table et donne une petite poussée brusque à une extrémité : tu vois une zone de compression qui se déplace le long du ressort, suivie d'une zone de dilatation. Chaque spire vibre \emph{dans le sens du déplacement} de la compression. C'est exactement le mécanisme d'une onde P.

\begin{definition}[Ondes P]
Les ondes \cle{P} (ondes \emph{primaires}) sont des ondes de \cle{compression-dilatation} : elles compriment puis dilatent successivement les roches traversées. La vibration des particules de roche est \emph{parallèle} à la direction de propagation de l'onde (ondes longitudinales).
\end{definition}

\begin{propriete}[Propriétés des ondes P]
\begin{itemize}
\item Ce sont les ondes sismiques \emph{les plus rapides} : leur vitesse varie d'environ \SI{6}{km/s} dans la croûte à environ \SI{13}{km/s} dans le manteau inférieur ;
\item elles se propagent dans \emph{tous} les milieux : solides \textbf{et} liquides (les liquides se laissent comprimer) ;
\item elles sont donc toujours les \emph{premières} enregistrées par une station sismique — d'où leur nom d'ondes \emph{primaires}.
\end{itemize}
\end{propriete}

\courssub{Les ondes S : cisaillement}

\textbf{Observation.}\par

Tiens une corde tendue par une extrémité et secoue-la latéralement d'un coup sec : une ondulation se déplace le long de la corde. Chaque point de la corde vibre \emph{perpendiculairement} à la direction dans laquelle l'ondulation avance. C'est le mécanisme d'une onde S.

\begin{definition}[Ondes S]
Les ondes \cle{S} (ondes \emph{secondaires}) sont des ondes de \cle{cisaillement} : elles déforment les roches par un mouvement de glissement latéral. La vibration des particules de roche est \emph{perpendiculaire} à la direction de propagation de l'onde (ondes transversales).
\end{definition}

\begin{propriete}[Propriétés des ondes S]
\begin{itemize}
\item Elles sont \emph{plus lentes} que les ondes P : leur vitesse varie d'environ \SI{3,5}{km/s} à \SI{7}{km/s} ;
\item elles ne se propagent \emph{que dans les solides}. Un liquide ne peut pas être cisaillé : il s'écoule au lieu de transmettre la déformation. Les ondes S sont donc \emph{arrêtées} par tout milieu liquide ;
\item elles arrivent en \emph{second} sur les sismogrammes — d'où leur nom d'ondes \emph{secondaires}.
\end{itemize}
\end{propriete}

\begin{aretenir}[L'essentiel sur les ondes P et S]
\begin{center}
\begin{tabularx}{\linewidth}{|l|X|X|}
\hline
\rowcolor{popblueL} \textbf{Critère} & \textbf{Ondes P} & \textbf{Ondes S} \\
\hline
Type de déformation & Compression-dilatation & Cisaillement \\
\hline
Vibration des particules & Parallèle à la propagation & Perpendiculaire à la propagation \\
\hline
Vitesse & \SI{6}{km/s} à \SI{13}{km/s} (les plus rapides) & \SI{3,5}{km/s} à \SI{7}{km/s} \\
\hline
Milieux traversés & Solides \textbf{et} liquides & Solides \textbf{uniquement} \\
\hline
Ordre d'arrivée & Premières & Secondes \\
\hline
\end{tabularx}
\end{center}
Le fait décisif pour la suite de la leçon : \textbf{les ondes S ne traversent pas les liquides}. C'est cette propriété qui révélera l'existence du noyau externe liquide.
\end{aretenir}

\begin{popfigure}
\begin{center}
\begin{tikzpicture}[scale=0.95]
% ---- Onde P : ressort ----
\node[popdark, font=\small\bfseries] at (-3.5,2.6) {Onde P (compression-dilatation)};
% ressort avec spires resserrées puis espacées
\draw[popblue, thick, decorate, decoration={coil, aspect=0.4, segment length=3.2pt, amplitude=4pt}] (-6.5,1.6) -- (-4.2,1.6);
\draw[popblue, thick, decorate, decoration={coil, aspect=0.4, segment length=1.6pt, amplitude=4pt}] (-4.2,1.6) -- (-3.4,1.6);
\draw[popblue, thick, decorate, decoration={coil, aspect=0.4, segment length=3.2pt, amplitude=4pt}] (-3.4,1.6) -- (-1.1,1.6);
\draw[popblue, thick, decorate, decoration={coil, aspect=0.4, segment length=1.6pt, amplitude=4pt}] (-1.1,1.6) -- (-0.3,1.6);
\draw[popblue, thick, decorate, decoration={coil, aspect=0.4, segment length=3.2pt, amplitude=4pt}] (-0.3,1.6) -- (1.6,1.6);
\draw[->, very thick, popink] (-6.5,0.9) -- (1.6,0.9) node[midway, below]{\scriptsize sens de propagation};
\draw[<->, poppurple] (-4.0,2.15) -- (-3.6,2.15);
\node[poppurple] at (-3.8,2.4) {\scriptsize vibration};
% ---- Onde S : corde ----
\node[popdark, font=\small\bfseries] at (-3.5,-0.6) {Onde S (cisaillement)};
\draw[poporange, very thick] plot[smooth, domain=-6.5:1.6, samples=120] (\x, {0.55*sin(240*(\x+6.5)/8.1) - 1.6});
\draw[->, very thick, popink] (-6.5,-2.6) -- (1.6,-2.6) node[midway, below]{\scriptsize sens de propagation};
\draw[<->, poppurple] (-3.9,-2.15) -- (-3.9,-1.05);
\node[poppurple] at (-3.9,-0.85) {\scriptsize vibration};
\end{tikzpicture}
\end{center}
\legende{Comparaison des déformations : pour l'onde P (ressort), les zones de compression (spires resserrées) et de dilatation se déplacent et la vibration est parallèle à la propagation ; pour l'onde S (corde), la vibration est perpendiculaire à la propagation.}
\end{popfigure}

\courssub{L'enregistrement des ondes : sismographe et sismogramme}

\textbf{Comment « écouter » la Terre ?}\par

Pour détecter les vibrations du sol, on utilise un appareil dont le principe est simple : une masse lourde suspendue à un ressort reste quasiment immobile par inertie pendant que le sol (et le bâti de l'appareil) vibre ; un stylet solidaire du bâti trace sur un cylindre tournant (ou un capteur électronique enregistre) l'écart entre le sol et la masse.

\begin{definition}[Sismographe et sismogramme]
Le \cle{sismographe} est l'appareil qui détecte et enregistre les vibrations du sol. L'enregistrement obtenu s'appelle un \cle{sismogramme} : c'est une courbe donnant l'amplitude des vibrations du sol en fonction du temps.
\end{definition}

\textbf{Lecture d'un sismogramme.} Sur un sismogramme, on observe d'abord une ligne presque plate (le calme), puis une première arrivée d'ondes d'amplitude modérée : ce sont les \textbf{ondes P}. Après un intervalle de temps, arrive un second train d'ondes, d'amplitude généralement plus forte : ce sont les \textbf{ondes S}. L'intervalle de temps entre l'arrivée des P et celle des S est appelé \cle{écart S-P} (noté $t_S - t_P$).

\begin{propriete}[Les P arrivent toujours avant les S — propriété admise]
Sur tout sismogramme enregistré à distance d'un séisme, les ondes P sont \emph{toujours} enregistrées avant les ondes S. Cette propriété, admise en Première, se justifie par l'observation systématique des sismogrammes et s'explique par le fait que les ondes P sont plus rapides que les ondes S dans tous les matériaux de la croûte et du manteau.
\end{propriete}

\begin{popfigure}
\begin{center}
\begin{tikzpicture}
\begin{axis}[
  width=\linewidth, height=5.2cm,
  xlabel={Temps (s)}, ylabel={Amplitude},
  xmin=0, xmax=60, ymin=-1.6, ymax=1.6,
  xtick={0,10,20,30,40,50,60}, ytick=\empty,
  axis lines=left, clip=false]
\addplot[popblue, thick, domain=0:12, samples=60] {0.04*sin(400*x)};
\addplot[popblue, thick, domain=12:26, samples=200] {0.45*sin(600*(x-12))*exp(-(x-12)/6)};
\addplot[popblue, thick, domain=26:60, samples=300] {1.15*sin(500*(x-26))*exp(-(x-26)/10)};
\draw[<-, popink, thick] (axis cs:14,0.5) -- (axis cs:14,1.25) node[above]{\small arrivée des P};
\draw[<-, poporange, thick] (axis cs:28,1.0) -- (axis cs:28,1.45) node[above]{\small arrivée des S};
\draw[<->, poppurple, thick] (axis cs:12,-1.25) -- (axis cs:26,-1.25) node[midway, below]{\small écart S-P};
\end{axis}
\end{tikzpicture}
\end{center}
\legende{Sismogramme simplifié : les ondes P, plus rapides, arrivent en premier (amplitude modérée), puis les ondes S (amplitude plus forte). L'écart temporel S-P augmente avec la distance de la station à l'épicentre.}
\end{popfigure}

\courssub{L'écart S-P : une mesure de la distance au séisme}

\textbf{Démarche d'investigation.}\par

\emph{Observation} : après un séisme, on compare les sismogrammes de plusieurs stations. À la station la plus proche de l'épicentre, les ondes P et S arrivent presque en même temps ; plus la station est éloignée, plus l'écart S-P est grand.

\emph{Hypothèse} : l'écart S-P est lié à la distance parcourue, car les deux ondes voyagent à des vitesses différentes.

\emph{Raisonnement sur les données} : les ondes P et S partent du foyer au même instant. Si $d$ est la distance à parcourir, $v_P$ et $v_S$ les vitesses moyennes des ondes P et S, alors les temps de parcours sont $t_P = \dfrac{d}{v_P}$ et $t_S = \dfrac{d}{v_S}$. L'écart vaut donc :
\[ t_S - t_P = d\left(\frac{1}{v_S} - \frac{1}{v_P}\right). \]
Comme $v_S < v_P$, cet écart est \emph{proportionnel} à la distance $d$ : il croît avec la distance à l'épicentre.

\emph{Conclusion} : la mesure de l'écart S-P sur un sismogramme permet de calculer la distance entre la station et le séisme. C'est le principe de la \emph{localisation} d'un séisme : avec trois stations, on trace trois cercles centrés sur chaque station, de rayon égal à la distance calculée ; leur intersection donne la position de l'épicentre.

\begin{exoresolu}[Calculer la distance d'un séisme à une station]
Énoncé. Une station sismique de Garoua enregistre un séisme : les ondes P arrivent à 14 h 02 min 10 s et les ondes S à 14 h 02 min 34 s. On admet que dans la croûte, $v_P = \SI{6,0}{km/s}$ et $v_S = \SI{3,5}{km/s}$. Calculer la distance $d$ entre la station et le foyer (supposé proche de la surface).
\tcblower
Solution détaillée.
\begin{enumerate}
\item \textbf{Écart S-P} : $t_S - t_P = 34 - 10 = \SI{24}{s}$.
\item \textbf{Relation} : $t_S - t_P = d\left(\dfrac{1}{v_S} - \dfrac{1}{v_P}\right)$.
\item \textbf{Calcul du facteur} :
\begin{align*}
\frac{1}{v_S} - \frac{1}{v_P} &= \frac{1}{3.5} - \frac{1}{6.0} = 0.2857 - 0.1667 = 0.1190 \ \text{s/km}.
\end{align*}
\item \textbf{Distance} :
\[ d = \frac{t_S - t_P}{0.1190} = \frac{24}{0.1190} \approx \SI{202}{km}. \]
\end{enumerate}
Le séisme s'est produit à environ \SI{200}{km} de la station. \emph{Vérification} : $t_P = 202/6.0 \approx \SI{34}{s}$ et $t_S = 202/3.5 \approx \SI{58}{s}$ ; l'écart est bien $58 - 34 = \SI{24}{s}$.
\end{exoresolu}

\begin{attention}[Pièges fréquents sur l'écart S-P]
\begin{itemize}
\item L'écart S-P donne la \emph{distance} au séisme, \textbf{pas sa direction} : une seule station ne suffit pas à localiser l'épicentre, il en faut au moins trois.
\item Ne conclus jamais « l'écart S-P est grand donc le séisme est puissant » : l'écart S-P renseigne sur la \emph{distance}, pas sur la \emph{magnitude}. C'est l'\emph{amplitude} des vibrations qui est liée à l'énergie libérée.
\item Dans les calculs, respecte la cohérence des unités : vitesses en km/s, temps en secondes, distance en kilomètres.
\end{itemize}
\end{attention}

\begin{exemplebox}[Pourquoi ressent-on deux secousses ?]
Lors d'un séisme ressenti à Bafoussam, un témoin décrit souvent « deux secousses successives » : une première vibration rapide, verticale, puis quelques secondes plus tard une secousse horizontale plus forte. La première correspond à l'arrivée des ondes P (compression, mouvement dans le sens de propagation), la seconde aux ondes S (cisaillement, mouvement perpendiculaire), plus lentes mais souvent plus destructrices. L'écart de quelques secondes entre les deux secousses est tout simplement l'écart S-P ressenti par le corps humain !
\end{exemplebox}

\begin{aretenir}[Bilan de la section]
\begin{itemize}
\item Un séisme est une vibration brutale du sol due à la libération d'énergie au \textbf{foyer} (en profondeur) ; l'\textbf{épicentre} est le point de surface à sa verticale.
\item L'énergie se propage par des \textbf{ondes sismiques} : ondes de volume (P et S) et ondes de surface (L et R).
\item Les ondes \textbf{P} (compression-dilatation, vibration parallèle) sont les plus rapides (\SI{6}{} à \SI{13}{km/s}) et traversent solides et liquides ; les ondes \textbf{S} (cisaillement, vibration perpendiculaire) sont plus lentes (\SI{3,5}{} à \SI{7}{km/s}) et ne traversent que les solides.
\item Le \textbf{sismographe} enregistre un \textbf{sismogramme} sur lequel les P arrivent avant les S ; l'\textbf{écart S-P} croît avec la distance à l'épicentre, ce qui permet de localiser le séisme.
\end{itemize}
Dans la section suivante, nous verrons comment ces ondes se comportent lorsqu'elles rencontrent une discontinuité — et comment ce comportement trahit l'existence des couches de la Terre.
\end{aretenir}

\coursec{Comportement des ondes aux discontinuités}

\courssub{Transition depuis la section précédente}
Dans la section « Séismes et ondes sismiques », nous avons vu qu'un séisme génère deux types d'ondes de volume : les ondes \cle{P} (primaires, longitudinales) et les ondes \cle{S} (secondaires, transversales). Leur vitesse de propagation n'est pas constante dans la Terre ; elle dépend des caractéristiques physiques du milieu traversé. C'est précisément cette variation de vitesse qui permet aux sismologues de « radiographier » l'intérieur du globe. Nous allons maintenant analyser comment les ondes se comportent lorsqu'elles rencontrent une surface de séparation entre deux milieux de propriétés différentes : une \cle{discontinuité sismique}.

\courssub{La vitesse des ondes dépend des propriétés physiques du milieu}
\textbf{Observation.}\par
Si l'on mesure la vitesse des ondes P et S dans différentes roches en laboratoire, on constate que le granite, le basalte, le péridotite ou le fer liquide ne transmettent pas les ondes à la même vitesse.

\textbf{Hypothèse.}\par
La vitesse $v$ d'une onde élastique dans un milieu isotrope dépend de sa \cle{densité} $\rho$ et de ses \cle{modules d'élasticité} : le module de compressibilité $K$ (résistance au changement de volume) et le module de rigidité $\mu$ (résistance au cisaillement).

\textbf{Formules de référence.}\par
Pour les ondes P (compressionnelles) :
\[
v_P = \sqrt{\frac{K + \frac{4}{3}\mu}{\rho}}
\]
Pour les ondes S (de cisaillement) :
\[
v_S = \sqrt{\frac{\mu}{\rho}}
\]

\textbf{Interprétation.}\par
\begin{itemize}
  \item Une augmentation de la rigidité $\mu$ accélère les deux types d'ondes.
  \item Une augmentation de la densité $\rho$ tend à ralentir les ondes, mais l'effet de la rigidité domine généralement dans les roches du manteau.
  \item Dans un \cle{liquide}, le module de rigidité $\mu$ est nul : les ondes S ne peuvent pas s'y propager ($v_S = 0$), tandis que les ondes P y persistent avec une vitesse $v_P = \sqrt{K/\rho}$.
\end{itemize}

\textbf{Conclusion.}\par
La vitesse des ondes sismiques est un indicateur direct de l'état physique (solide/liquide), de la composition minéralogique et de la température/pression du milieu traversé.

\begin{definition}[Discontinuité sismique]
Une \textbf{discontinuité sismique} est une surface de séparation entre deux couches de la Terre dont les propriétés physiques (densité, rigidité, état de la matière) changent brutalement. Elle se manifeste sur les sismogrammes par une variation soudaine de la vitesse des ondes P et/ou S, et par l'apparition d'ondes réfléchies et réfractées.
\end{definition}

\courssub{Réflexion et réfraction à une discontinuité : analogie avec l'optique}
Lorsqu'une onde sismique rencontre une discontinuité, deux phénomènes se produisent simultanément :
\begin{enumerate}
  \item \textbf{Réflexion} : une partie de l'énergie rebondit dans le milieu d'incidence. L'angle de réflexion $r$ est égal à l'angle d'incidence $i$ (loi de Descartes).
  \item \textbf{Réfraction (transmission)} : l'autre partie franchit la discontinuité en changeant de direction et de vitesse. La loi de Snell‑Descartes s'applique :
  \[
  \frac{\sin i}{v_1} = \frac{\sin t}{v_2}
  \]
  où $v_1$ et $v_2$ sont les vitesses dans le milieu incident et le milieu transmis, $i$ l'angle d'incidence et $t$ l'angle de transmission (réfraction).
\end{enumerate}
Si $v_2 > v_1$, le rayon transmis s'éloigne de la normale ($t > i$) ; si $v_2 < v_1$, il se rapproche de la normale ($t < i$). Au-delà d'un \cle{angle critique} $i_c = \arcsin(v_1/v_2)$, la transmission disparaît et on observe une réflexion totale.

\begin{popfigure}
\begin{tikzpicture}[scale=1.2, >=Stealth]
% Couche supérieure
\draw[popblue!30, fill=popblue!10] (-3,0) rectangle (3,1.5);
\node[popblue] at (0,0.75) {Milieu 1 ($v_1$)};
% Couche inférieure
\draw[poporange!30, fill=poporange!10] (-3,-2) rectangle (3,0);
\node[poporange] at (0,-1) {Milieu 2 ($v_2$)};
% Discontinuité
\draw[thick, popdark] (-3,0) -- (3,0) node[midway, above=2pt] {Discontinuité};
% Normale
\draw[dashed, thin] (0,-2) -- (0,1.5) node[midway, right] {Normale};
% Rayon incident
\draw[->, thick, poppink] (-1.5,1.5) -- (0,0) node[midway, left=2pt] {Incident};
% Angle incidence
\draw[poppink] (0,0) ++(150:0.6) arc (150:180:0.6) node[midway, left=2pt] {$i$};
% Rayon réfléchi
\draw[->, thick, popgreen] (0,0) -- (1.5,1.5) node[midway, right=2pt] {Réfléchi};
% Angle réflexion
\draw[popgreen] (0,0) ++(30:0.6) arc (30:0:0.6) node[midway, right=2pt] {$r=i$};
% Rayon réfracté (transmis)
\draw[->, thick, poppurple] (0,0) -- (1.2,-1.8) node[midway, right=2pt] {Réfracté};
% Angle réfraction
\draw[poppurple] (0,0) ++(-30:0.6) arc (-30:0:0.6) node[midway, right=2pt] {$t$};
% Légendes vitesses
\node[popblue] at (-2.5,1.2) {$v_1$};
\node[poporange] at (-2.5,-1.2) {$v_2$};
\end{tikzpicture}
\legende{Schéma d'un rayon sismique P ou S rencontrant une discontinuité plane. L'onde incidente (rose) se partage en une onde réfléchie (verte) et une onde réfractée (violette). Les angles sont mesurés par rapport à la normale.}
\end{popfigure}

\courssub{Comportement particulier des ondes S et P dans les liquides}
\begin{propriete}[Propagation des ondes dans un milieu liquide]
Les \textbf{ondes S ne se propagent pas dans les liquides} (module de rigidité $\mu = 0 \Rightarrow v_S = 0$). Les \textbf{ondes P y sont ralenties} par rapport au solide environnant (seul le module de compressibilité $K$ intervient) et subissent une forte réfraction. Cette propriété est admise au Probatoire comme conséquence de la mécanique des milieux continus.
\end{propriete}

\textbf{Conséquence majeure : la découverte du noyau liquide.}\par
Lorsque les sismologues ont observé qu'au-delà d'une distance épicentrale d'environ $103^\circ$ les ondes S directes disparaissaient complètement, tandis que les ondes P présentaient une \cle{zone d'ombre} entre $103^\circ$ et $143^\circ$, ils en ont déduit l'existence d'un noyau externe liquide (arrêt des ondes S, ralentissement et déviation des ondes P). C'est une illustration parfaite de la démarche d'investigation : \emph{observation} (zones d'ombre) $\rightarrow$ \emph{hypothèse} (noyau liquide) $\rightarrow$ \emph{modélisation} (calcul des trajectoires) $\rightarrow$ \emph{confirmation} (accord avec les temps de parcours).

\begin{popfigure}
\begin{tikzpicture}[scale=0.9]
\begin{axis}[
    width=\linewidth,
    height=8cm,
    xlabel={Profondeur (km)},
    ylabel={Vitesse (km/s)},
    xmin=0, xmax=6371,
    ymin=0, ymax=14,
    xtick={0, 30, 2891, 5150, 6371},
    xticklabels={Surface, 30 (Moho), 2891 (Gutenberg), 5150 (Lehmann), Centre},
    ytick={0,2,4,6,8,10,12,14},
    grid=major,
    grid style={popline},
    legend style={at={(0.02,0.98)}, anchor=north west, fill=popcream},
    tick label style={font=\small},
    label style={font=\small},
]
% Ondes P
\addplot[popcurve, popblue, thick, domain=0:30, samples=50] {5.8 + 0.02*x};
\addplot[popcurve, popblue, thick, domain=30:2891, samples=200] {6.4 + 0.0005*(x-30)};
\addplot[popcurve, popblue, thick, domain=2891:5150, samples=200] {8.0 + 0.0003*(x-2891)};
\addplot[popcurve, popblue, thick, domain=5150:6371, samples=100] {10.5 + 0.0001*(x-5150)};
% Sauts aux discontinuités (traits verticaux)
\draw[popblue, dashed] (axis cs:30,5.8) -- (axis cs:30,6.4);
\draw[popblue, dashed] (axis cs:2891,8.0) -- (axis cs:2891,8.0);
\draw[popblue, dashed] (axis cs:5150,10.5) -- (axis cs:5150,10.5);
% Ondes S
\addplot[popcurve, poporange, thick, domain=0:30, samples=50] {3.2 + 0.015*x};
\addplot[popcurve, poporange, thick, domain=30:2891, samples=200] {3.6 + 0.0004*(x-30)};
% Arrêt des ondes S dans le noyau externe
\draw[poporange, thick, dashed] (axis cs:2891,4.6) -- (axis cs:5150,4.6);
\addplot[popcurve, poporange, thick, domain=5150:6371, samples=100] {5.0 + 0.0001*(x-5150)};
\legend{Ondes P ($v_P$), Ondes S ($v_S$)};
\end{axis}
\end{tikzpicture}
\legende{Variation de la vitesse des ondes P (bleu) et S (orange) en fonction de la profondeur. Les sauts brutaux correspondent aux grandes discontinuités : Moho (30~km), Gutenberg (2891~km, limite manteau/noyau externe), Lehmann (5150~km, limite noyau externe/noyau interne). Les ondes S s'annulent dans le noyau externe liquide (zone grisée).}
\end{popfigure}

\courssub{Notion de zone d'ombre}
Une \cle{zone d'ombre} est une région à la surface de la Terre où aucune onde directe (P ou S) n'est reçue après un séisme, en raison de la réfraction (et parfois de l'arrêt) des ondes à l'intérieur du globe.
\begin{itemize}
  \item \textbf{Zone d'ombre des ondes S :} au-delà de $\approx 103^\circ$ de distance épicentrale, les ondes S directes ne sont pas enregistrées. Elles sont bloquées par le noyau externe liquide.
  \item \textbf{Zone d'ombre des ondes P :} entre $\approx 103^\circ$ et $143^\circ$, les ondes P directes manquent. Elles sont fortement déviées (réfractées) en traversant le noyau externe à vitesse plus faible, puis reviennent en surface au-delà de $143^\circ$.
\end{itemize}
L'existence de ces zones d'ombre est une preuve directe de la structure en couches et de l'état liquide du noyau externe.

\begin{methode}[Repérer une discontinuité sur un graphique vitesse–profondeur]
\begin{enumerate}
  \item Tracer les courbes $v_P(z)$ et $v_S(z)$ à partir des données de temps de parcours (ou utiliser une courbe de référence comme la figure ci‑dessus).
  \item Rechercher les \textbf{ruptures de pente} (changement brutal de la dérivée) ou les \textbf{sauts de vitesse} (discontinuité de la fonction).
  \item Identifier la profondeur correspondante : c'est la profondeur de la discontinuité.
  \item Associer chaque discontinuité à une interface connue : Moho (croûte/manteau), Gutenberg (manteau/noyau externe), Lehmann (noyau externe/noyau interne).
  \item Vérifier la cohérence : une discontinuité où $v_S$ tombe à zéro signale un passage au liquide (noyau externe).
\end{enumerate}
\end{methode}

\courssub{Exemple concret : localisation du Moho sous le Cameroun}
\begin{exoresolu}[Détermination de l'épaisseur de la croûte sous Yaoundé]
\textbf{Énoncé.} Un sismogramme enregistré à Yaoundé (station sismique) montre une onde P directe arrivant à $t_1 = 12,3$~s et une onde P réfléchie sur le Moho (notée PmP) arrivant à $t_2 = 15,8$~s pour un séisme local à distance épicentrale $\Delta = 50$~km. La vitesse moyenne dans la croûte est $v_P = 6,2$~km/s. Calculer l'épaisseur de la croûte $h$ sous la station.
\textbf{Solution détaillée.}
L'onde PmP effectue un aller‑retour dans la croûte : elle descend verticalement (incidence quasi normale pour un séisme proche), se réfléchit sur le Moho et remonte. Le temps supplémentaire $\Delta t = t_2 - t_1 = 3,5$~s correspond au parcours aller‑retour dans la croûte : $2h = v_P \times \Delta t$.
\[
h = \frac{v_P \times \Delta t}{2} = \frac{6{,}2 \times 3{,}5}{2} = 10{,}85 \text{ km}
\]
L'épaisseur de la croûte sous Yaoundé est donc d'environ \textbf{11~km}, ce qui est cohérent avec une croûte continentale amincie dans la zone du rift camerounais (proche du Mont Cameroun).
\end{exoresolu}

\begin{attention}[Pièges fréquents]
\begin{itemize}
  \item \textbf{Confondre discontinuité et changement d'état.} Une discontinuité sismique marque un changement brutal de vitesse, mais pas nécessairement un changement de phase (solide/liquide). Par exemple, le Moho sépare deux milieux solides (croûte et manteau) de composition différente.
  \item \textbf{Oublier que les ondes S ne traversent pas les liquides.} Au Probatoire, on doit pouvoir justifier l'absence d'ondes S dans le noyau externe par $\mu = 0$, et non par « les ondes S sont absorbées ».
  \item \textbf{Mésinterpréter la zone d'ombre.} La zone d'ombre des ondes P ($103^\circ$–$143^\circ$) n'est pas due à l'arrêt des ondes P, mais à leur forte déviation dans le noyau externe. Les ondes P réapparaissent au‑delà de $143^\circ$ après avoir traversé le noyau.
  \item \textbf{Utiliser des angles en degrés dans les formules trigonométriques sans conversion.} La loi de Snell s'écrit avec le sinus de l'angle ; assurez‑vous que votre calculatrice est en mode degrés (ou convertissez en radians).
  \item \textbf{Écrire les nombres décimaux avec une virgule en mode mathématique.} En mode math ($...$), écrivez $6{,}2$ et non $6,2$ pour éviter l'espace parasite après la virgule.
\end{itemize}
\end{attention}

\begin{savaistu}[Le sismomètre de l'IRGM à Yaoundé]
L'Institut de Recherche Géologique et Minière (IRGM) de Yaoundé exploite un réseau de stations sismiques qui enregistrent en temps réel les séismes locaux (liens avec l'activité volcanique du Mont Cameroun) et télé‑séismes. Les données de ces stations alimentent les modèles de vitesse utilisés pour la prévention des risques sismiques au Cameroun et contribuent au réseau sismologique mondial (ISC, USGS). C'est un exemple concret de la manière dont la sismologie fondamentale sert la sécurité des populations.
\end{savaistu}

\courssub{Synthèse de la section}
\begin{aretenir}[Points clés à retenir]
\begin{itemize}
  \item La vitesse des ondes sismiques dépend de la densité, de la rigidité et de l'état du milieu ($v_P = \sqrt{(K+4\mu/3)/\rho}$, $v_S = \sqrt{\mu/\rho}$).
  \item Une \cle{discontinuité sismique} est une surface où la vitesse change brutalement ; elle génère réflexion et réfraction (loi de Snell‑Descartes).
  \item Les ondes S ne se propagent pas dans les liquides ($\mu=0$) ; les ondes P y sont ralenties et déviées.
  \item Les \cle{zones d'ombre} (absence d'ondes directes) révèlent la structure en couches et l'existence du noyau externe liquide.
  \item Sur un graphique $v(z)$, une discontinuité se repère par une rupture de pente ou un saut de vitesse.
\end{itemize}
\end{aretenir}

\coursec{Exploitation des temps de parcours : la démarche de la sismologie}

Dans la section précédente, nous avons vu que les ondes sismiques ne traversent pas le globe de façon uniforme : les ondes P se réfractent et se réfléchissent aux discontinuités, tandis que les ondes S disparaissent dans les milieux liquides. Ces comportements, loin d'être de simples curiosités, constituent de véritables \cle{outils d'investigation}. Tout comme le médecin du Centre Hospitalier d'Essos à Yaoundé utilise l'échographie — des ondes ultrasonores réfléchies par les organes — pour « voir » à l'intérieur du corps sans opérer, le sismologue utilise les ondes sismiques pour « voir » à l'intérieur de la Terre sans y creuser. Mais comment passe-t-on concrètement d'un enregistrement sismique à une profondeur en kilomètres ? C'est toute la démarche de la \cle{sismologie} que nous allons maintenant construire, pas à pas.

\courssub{Le principe général : écouter la Terre en plusieurs stations}

\begin{definition}[Sismographe, sismogramme, station sismique]
Un \cle{sismographe} est un appareil qui enregistre les vibrations du sol. Le document obtenu est un \cle{sismogramme} : il indique, pour une station donnée, les \cle{heures d'arrivée} des ondes P (premières arrivées, car les plus rapides) puis des ondes S. Un ensemble de stations réparties à la surface du globe constitue un \cle{réseau sismologique}.
\end{definition}

Imaginons un séisme qui se produit dans la région de Foumban. Une station à Bafoussam enregistre les ondes P à 14 h 03 min 12 s, une station à Douala les enregistre plus tard, une station à N'Djamena encore plus tard. Chaque station fournit une donnée essentielle : le \cle{temps de parcours} $t$, c'est-à-dire la durée écoulée entre l'instant du séisme (l'origine, déterminée par comparaison entre stations) et l'arrivée de l'onde.

L'idée fondatrice de la sismologie est alors d'une simplicité remarquable : si l'on connaît la \cle{vitesse} $v$ de l'onde dans le milieu traversé et le \cle{temps de parcours} $t$, on remonte à la \cle{distance} parcourue $d$ grâce à la relation fondamentale de la cinématique :

\[ v = \frac{d}{t} \quad \Longleftrightarrow \quad d = v \times t \quad \Longleftrightarrow \quad t = \frac{d}{v} \]

\begin{aretenir}[La relation fondamentale]
Pour une onde se propageant à vitesse constante $v$ : la distance parcourue est $d = v \times t$, avec $v$ en \si{km/s}, $t$ en \si{s} et $d$ en \si{km}. Toute la sismologie repose sur cette relation, appliquée à des milliers d'enregistrements.
\end{aretenir}

En pratique, le raisonnement s'inverse : on \emph{mesure} $t$ sur les sismogrammes, on \emph{connaît} $d$ (distance entre l'épicentre et la station, mesurable à la surface), et on en \emph{déduit} la vitesse moyenne $v = d/t$ des roches traversées. En comparant les vitesses obtenues pour des trajets de plus en plus profonds, on détecte les changements de milieu : les discontinuités.

\courssub{Les hodochrones : la courbe qui trahit les discontinuités}

\begin{definition}[Hodochrone]
Une \cle{hodochrone} est la courbe représentant le \cle{temps de parcours} $t$ d'une onde sismique en fonction de la \cle{distance angulaire} $\Delta$ entre l'épicentre et la station réceptrice (mesurée en degrés, de 0° à 180°, l'angle au centre de la Terre entre épicentre et station).
\end{definition}

Si l'intérieur de la Terre était parfaitement homogène, les ondes iraient en ligne droite à vitesse constante : la hodochrone serait une simple droite passant par l'origine (plus la station est loin, plus le trajet est long, proportionnellement). Or les hodochrones réelles présentent deux particularités :

\begin{itemize}
\item une \cle{courbure} progressive : la vitesse des ondes augmente avec la profondeur (les roches sont plus denses et plus rigides sous pression), donc les ondes profondes arrivent « plus tôt » que prévu ;
\item des \cle{brisures} (changements brusques de pente) et des \cle{retards} localisés : ce sont les signatures des discontinuités, où la vitesse change brutalement.
\end{itemize}

\begin{popfigure}
\begin{center}
\begin{tikzpicture}
\begin{axis}[
  width=0.85\linewidth, height=7cm,
  xlabel={Distance angulaire $\Delta$ (degrés)},
  ylabel={Temps de parcours $t$ (min)},
  xmin=0, xmax=160, ymin=0, ymax=26,
  xtick={0,20,40,60,80,100,120,140},
  ytick={0,5,10,15,20,25},
  grid=both, grid style={popline!40},
  axis lines=left,
  legend style={at={(0.03,0.97)}, anchor=north west, font=\small}
]
% Branche 1 : ondes directes croûte/manteau supérieur (quasi droite)
\addplot[popcurve, domain=0:60, samples=60] {0.16*x};
\addlegendentry{Ondes directes (manteau supérieur)}
% Branche 2 : ondes ayant plongé plus profondément, courbure + brisure
\addplot[popcurve, poporange, domain=60:150, samples=80] {9.6 + 0.115*(x-60) - 0.0004*(x-60)^2};
\addlegendentry{Ondes profondes (réfractées)}
% Point de brisure
\addplot[only marks, mark=*, popdark, mark size=2pt] coordinates {(60,9.6)};
\node[popdark, font=\small, anchor=south west] at (axis cs:62,9.6) {Brisure : discontinuité};
\end{axis}
\end{tikzpicture}
\end{center}
\legende{Hodochrone simplifiée d'ondes P : le temps de parcours croît avec la distance angulaire, mais la courbe présente une \textbf{brisure} vers 60°. Ce changement brutal de pente signale que les ondes ont rencontré une discontinuité où leur vitesse change brutalement.}
\end{popfigure}

\begin{attention}[Interpréter une brisure]
Une brisure sur une hodochrone n'est jamais un « hasard » ou une erreur de mesure : elle indique qu'à partir d'une certaine distance, les ondes empruntent un trajet différent (plus profond, réfléchi ou réfracté) traversant un milieu de propriétés différentes. C'est l'\emph{anomalie} qui renseigne, pas la régularité.
\end{attention}

\courssub{Le raisonnement de l'écho : mesurer une profondeur sans creuser}

Revenons à l'analogie de l'échographie. Quand on crie devant la falaise du Mont Cameroun, l'écho revient après un certain temps : plus la falaise est loin, plus l'écho est tardif. Le son a fait un \cle{aller-retour}. La sismologie de réflexion applique exactement ce principe : une onde P descendante se \emph{réfléchit} sur une discontinuité, puis remonte jusqu'à la station. Le temps mesuré $t$ correspond donc à un trajet \emph{double}.

\begin{propriete}[Profondeur d'une discontinuité par écho sismique — démonstration guidée exigible]
Si une onde sismique de vitesse moyenne $v$ se réfléchit sur une discontinuité située à la profondeur $h$ et revient à la station après un temps aller-retour $t$, alors :
\[ h = \frac{v \times t}{2} \]
\textbf{Démonstration guidée (à établir en classe) :}
\begin{enumerate}
\item L'onde parcourt la profondeur $h$ à la descente, puis la même profondeur $h$ à la remontée. Quelle est la distance totale parcourue ? \emph{Réponse attendue :} $d = 2h$.
\item Appliquons la relation fondamentale $d = v \times t$ à ce trajet total : $2h = v \times t$.
\item Isolons $h$ en divisant les deux membres par 2 : $h = \dfrac{v \times t}{2}$. \hfill $\blacksquare$
\end{enumerate}
\end{propriete}

\begin{methode}[Exploiter des temps de parcours pour situer une discontinuité]
\begin{enumerate}
\item \textbf{Identifier l'onde et sa vitesse} : s'agit-il d'une onde P ou S ? Quelle vitesse moyenne $v$ est donnée (ou calculable par $v = d/t$ sur un trajet connu) ?
\item \textbf{Mesurer le temps} : relever sur le sismogramme le temps $t$ entre l'émission (ou l'arrivée de l'onde directe) et l'arrivée de l'onde réfléchie.
\item \textbf{Vérifier la nature du trajet} : l'onde a-t-elle fait un aller simple ou un aller-retour ? Si c'est un aller-retour (écho), appliquer $h = \dfrac{v \times t}{2}$ ; si c'est un aller simple, $d = v \times t$.
\item \textbf{Calculer} en veillant à la cohérence des unités ($v$ en \si{km/s}, $t$ en \si{s}, $h$ en \si{km}).
\item \textbf{Vérifier la cohérence du résultat} : la profondeur obtenue est-elle compatible avec les enveloppes connues (croûte : quelques km à $\sim 70$~km ; manteau : jusqu'à $\sim 2900$~km) ? Un résultat de $0{,}5$~km ou de $50\,000$~km doit alerter.
\end{enumerate}
\end{methode}

\begin{exemplebox}[Un écho sous une chaîne de montagnes]
Sous une chaîne de montagnes, les géophysiciens enregistrent une onde P réfléchie qui revient à la station $16$~s après son émission, avec une vitesse moyenne $v = 8$~km/s dans la croûte. La profondeur de la discontinuité réfléchissante est :
\[ h = \frac{v \times t}{2} = \frac{8 \times 16}{2} = \frac{128}{2} = 64~\si{km} \]
La discontinuité se situe à $64$~km de profondeur : il s'agit d'une racine crustale profonde, typique des chaînes de montagnes dont la croûte est épaissie.
\end{exemplebox}

\begin{attention}[L'erreur classique de l'aller-retour]
L'erreur la plus fréquente au Probatoire est d'\textbf{oublier de diviser par 2} le temps d'un trajet aller-retour. Dans l'exemple précédent, écrire $h = 8 \times 16 = 128$~km revient à confondre la distance totale parcourue ($2h$) avec la profondeur cherchée ($h$). Avant tout calcul, demande-toi toujours : \emph{l'onde a-t-elle fait un aller simple ou un aller-retour ?}
\end{attention}

\begin{exoresolu}[Temps de parcours et profondeur du Moho]
Lors d'une campagne de sismique réflexion dans la plaine du Logone, une onde P de vitesse moyenne $v = 6{,}5$~km/s se réfléchit sur la discontinuité de Mohorovičić (Moho) et revient à la station au bout de $t = 10$~s.
\begin{enumerate}
\item Calculer la profondeur $h$ du Moho sous cette région.
\item Ce résultat est-il cohérent avec une croûte continentale ? Justifier.
\end{enumerate}
\tcblower
\textbf{Solution détaillée.}
\begin{enumerate}
\item Le trajet est un aller-retour, donc la distance totale parcourue est $2h = v \times t$, d'où :
\[ h = \frac{v \times t}{2} = \frac{6{,}5 \times 10}{2} = \frac{65}{2} = 32{,}5~\si{km} \]
\item La croûte continentale a une épaisseur moyenne de $30$ à $40$~km (jusqu'à $70$~km sous les montagnes, $5$ à $10$~km sous les océans). Une profondeur de $32{,}5$~km est donc parfaitement cohérente avec le Moho d'une croûte continentale de plaine.
\end{enumerate}
\end{exoresolu}

\courssub{Les zones d'ombre : quand l'absence d'onde devient une preuve}

La démarche la plus spectaculaire de la sismologie ne repose pas sur ce que l'on enregistre, mais sur ce que l'on \emph{n'enregistre pas}. En compilant les sismogrammes de stations réparties sur tout le globe après de grands séismes, les sismologues ont constaté deux faits troublants :

\begin{itemize}
\item \textbf{Au-delà d'environ 103°} de distance angulaire, les \cle{ondes S disparaissent totalement} : aucune station située entre 103° et 180° ne reçoit d'ondes S directes. C'est la \cle{zone d'ombre des ondes S}.
\item \textbf{Entre environ 103° et 142°}, les \cle{ondes P directes disparaissent également}, puis \emph{réapparaissent au-delà de 142°}, affaiblies et déviées. C'est la \cle{zone d'ombre des ondes P}.
\end{itemize}

\begin{experience}[Démarche d'investigation : que cache la zone d'ombre ?]
\textbf{Observation :} après un grand séisme, aucune onde S n'est enregistrée au-delà de 103°, et les ondes P disparaissent entre 103° et 142° puis réapparaissent au-delà.
\begin{enumerate}
\item \textbf{Hypothèse 1 :} il existe un obstacle central qui bloque les ondes S. Or les ondes S ne se propagent pas dans les liquides. $\Rightarrow$ L'obstacle contiendrait un milieu \emph{liquide}.
\item \textbf{Hypothèse 2 :} les ondes P, qui traversent les liquides, sont \emph{déviées} (réfractées) en entrant dans cet obstacle, ce qui explique leur absence entre 103° et 142° et leur réapparition au-delà.
\item \textbf{Données complémentaires :} la vitesse des ondes P chute brutalement vers $2\,900$~km de profondeur (de $\sim 13$~km/s à $\sim 8$~km/s), signature d'un changement d'état.
\item \textbf{Interprétation :} la chute de vitesse des P et la disparition des S à la même profondeur indiquent une discontinuité majeure vers $2\,900$~km, séparant le manteau solide d'un milieu liquide.
\item \textbf{Conclusion :} la Terre possède un \cle{noyau externe liquide}, d'environ $2\,900$~km à $5\,100$~km de profondeur. Les ondes P réfractées à l'entrée et à la sortie de ce noyau sont déviées, créant la zone d'ombre entre 103° et 142° ; les ondes S, incapables de traverser le liquide, ne réapparaissent jamais.
\end{enumerate}
\end{experience}

\begin{popfigure}
\begin{center}
\begin{tikzpicture}[scale=1.05]
% Noyau externe (liquide)
\fill[poporangeL] (0,0) circle (1.5);
% Manteau
\fill[popgreenL, even odd rule] (0,0) circle (1.5) (0,0) circle (3);
% Croûte (ligne épaisse)
\draw[line width=1.6pt, popdark] (0,0) circle (3);
\draw[line width=0.8pt, poporange] (0,0) circle (1.5);
% Foyer (séisme) en haut à gauche
\coordinate (F) at (-2.4,1.8);
\fill[popink] (F) circle (0.09);
\node[popink, font=\small\bfseries, anchor=east] at ($(F)+(-0.08,0.15)$) {Foyer};
% Ondes S : bloquées par le noyau
\draw[->, poppurple, line width=1pt] (F) -- (1.05,-1.05);
\node[poppurple, font=\small, anchor=west] at (0.4,-1.6) {Onde S arrêtée};
\draw[->, poppurple, line width=1pt] (F) -- (0.2,-1.49);
% Onde P directe (reste dans le manteau)
\draw[->, popblue, line width=1pt] (F) to[bend left=15] (2.4,1.8);
\node[popblue, font=\small, anchor=south] at (1.1,2.35) {Onde P directe};
% Onde P réfractée par le noyau
\draw[->, popblue, line width=1pt] (F) -- (-0.5,-1.42);
\draw[popblue, line width=1pt] (-0.5,-1.42) -- (1.15,-0.96);
\draw[->, popblue, line width=1pt] (1.15,-0.96) -- (2.85,-0.85);
\node[popblue, font=\small, anchor=west] at (2.6,-1.3) {P réfractée};
% Zone d'ombre (arc entre 103° et 142°) côté droit-haut
\fill[poppink!45] (3,0) arc[start angle=0, end angle=38, radius=3] -- (2.366,1.847) arc[start angle=38, end angle=0, radius=2.366] -- cycle;
\node[popdark, font=\small, align=center] at (3.6,1.4) {Zone d'ombre\\(103°--142°)};
% Labels couches
\node[font=\small\bfseries, popdark] at (0,0) {Noyau};
\node[font=\small, popdark] at (0,0.55) {(externe : liquide)};
\node[font=\small, popdark!70] at (-1.7,-2.35) {Manteau (solide)};
\node[font=\small, popdark!70, anchor=west] at (-3.35,2.9) {Croûte};
\end{tikzpicture}
\end{center}
\legende{Coupe de la Terre montrant les trajectoires des ondes depuis un foyer. Les ondes S (violettes) sont stoppées par le noyau externe liquide. Les ondes P (bleues) qui pénètrent le noyau sont réfractées et déviées : elles réapparaissent au-delà de 142°, laissant une zone d'ombre entre 103° et 142° où aucune onde directe n'arrive.}
\end{popfigure}

\begin{aretenir}[Les deux preuves du noyau]
\begin{itemize}
\item La \cle{zone d'ombre des ondes S} (au-delà de 103°) prouve l'existence d'un \cle{noyau externe liquide}, car les ondes S ne se propagent pas dans les liquides.
\item La \cle{réapparition des ondes P au-delà de 142°} prouve que ces ondes sont \cle{réfractées} (déviées) par le noyau, dont les propriétés physiques diffèrent nettement de celles du manteau.
\end{itemize}
\end{aretenir}

\begin{savaistu}[Inge Lehmann et les ondes « impossibles »]
En 1936, la sismologue danoise \cle{Inge Lehmann} étudiait des sismogrammes de grands séismes. Elle y repéra des ondes P arrivant \emph{dans la zone d'ombre}, là où aucune onde P directe n'aurait dû parvenir — des ondes « impossibles » selon le modèle de l'époque (un noyau entièrement liquide). En appliquant rigoureusement la démarche des temps de parcours, elle montra que ces ondes avaient été \emph{réfléchies et réfractées} par une structure encore plus profonde : une \cle{graine}, c'est-à-dire un \cle{noyau interne solide} d'environ $1\,220$~km de rayon, au cœur du noyau externe liquide. Sa découverte, obtenue uniquement par le calcul à partir de temps d'arrivée, illustre la puissance extraordinaire de la sismologie : on peut découvrir un « continent » de fer solide à $5\,100$~km sous nos pieds sans jamais y descendre.
\end{savaistu}

\begin{exercice}[À toi de jouer]
Une onde P se propage dans le manteau à la vitesse moyenne $v = 10$~km/s. Elle se réfléchit sur la discontinuité située à la base du manteau (vers $2\,900$~km de profondeur) et revient à la station.
\begin{enumerate}
\item Calculer le temps aller-retour théorique $t$ de cette onde.
\item En réalité, la station enregistre l'écho au bout de $560$~s. Proposer une explication à cet écart.
\end{enumerate}
\end{exercice}

\begin{corrige}[Exercice : temps aller-retour dans le manteau]
\begin{enumerate}
\item Le trajet aller-retour vaut $2h = 2 \times 2\,900 = 5\,800$~km. Donc :
\[ t = \frac{2h}{v} = \frac{5\,800}{10} = 580~\si{s} \quad \text{soit environ } 9~\text{min}~40~\text{s}. \]
\item L'écho arrive à $560$~s, soit $20$~s \emph{plus tôt} que prévu. Cela signifie que la vitesse réelle moyenne est légèrement supérieure à $10$~km/s : en effet, la vitesse des ondes P \emph{augmente avec la profondeur} dans le manteau (les roches deviennent plus denses et plus rigides). La valeur $v = 10$~km/s n'était qu'une moyenne simplifiée. Cet écart illustre précisément la courbure des hodochrones.
\end{enumerate}
\end{corrige}

En résumé, la démarche de la sismologie est un remarquable exemple de raisonnement scientifique : à partir de simples \emph{temps d'arrivée} mesurés en surface, et de la seule relation $d = v \times t$ correctement appliquée (en n'oubliant jamais le facteur 2 des échos), on déduit les profondeurs des discontinuités, l'état physique des couches et même l'existence de structures insoupçonnées. Il ne nous reste plus qu'à assembler ces résultats pour construire, dans la prochaine section, le modèle complet en couches du globe terrestre — et à en discuter les limites.

\coursec{Les grandes discontinuités et les enveloppes de la Terre}

\courssub{Transition depuis l'exploitation des temps de parcours}
Dans la section précédente, nous avons vu comment les sismologues mesurent les temps de parcours des ondes P et S pour repérer les profondeurs où leurs vitesses changent brutalement. Ces ruptures de vitesse correspondent à des \textbf{discontinuités sismiques} : ce sont les frontières entre les grandes enveloppes internes de la Terre. Nous allons maintenant identifier les trois discontinuités majeures, décrire les enveloppes qu'elles délimitent et construire le modèle en couches qui en résulte.

\courssub{La discontinuité de Mohorovičić (Moho) : limite croûte / manteau}
\begin{definition}[Discontinuité de Mohorovičić (Moho)]
La \cle{Moho} est la surface de discontinuité sismique qui sépare la croûte terrestre du manteau supérieur. Elle se manifeste par une augmentation soudaine de la vitesse des ondes P (de $\approx 6,5\ \text{km/s}$ dans la croûte inférieure à $\approx 8,0\ \text{km/s}$ dans le manteau) et des ondes S, due au passage de roches felsiques (granitoïdes, gabbros) à des roches ultramafiques (péridotites).
\end{definition}

\begin{experience}[Mise en évidence de la Moho par réfraction sismique]
\textbf{Matériel :} sismogrammes d'un séisme lointain enregistrés par un réseau de stations alignées.\\
\textbf{Protocole :} on trace les temps de parcours $t$ en fonction de la distance épicentrale $\Delta$. On observe deux branches linéaires : la première (ondes directes dans la croûte) de pente $1/v_{\text{croûte}}$, la seconde (ondes réfractées le long de la Moho) de pente $1/v_{\text{manteau}}$ plus faible.\\
\textbf{Observation :} le point de rupture des deux droites donne le temps critique $t_c$ et la distance critique $\Delta_c$.\\
\textbf{Interprétation :} en supposant une croûte plane d'épaisseur $h$ et des vitesses constantes $v_1$ (croûte) et $v_2$ (manteau), la géométrie du rayon réfracté donne $t_c = \frac{2h}{v_1}\sqrt{1-(v_1/v_2)^2}$. On en déduit $h$.
\end{experience}

\begin{exemplebox}[Profondeur de la Moho sous le Cameroun]
Au Cameroun, les études de sismologie de réfraction (ex. campagne \emph{CRUSTAL} 2015) montrent une Moho à $\approx 35\ \text{km}$ sous le plateau de l'Adamaoua (croûte continentale épaisse) et à $\approx 12\ \text{km}$ au large de Kribi (croûte océanique). Cela illustre la variabilité latérale de la discontinuité.
\end{exemplebox}

\courssub{La discontinuité de Gutenberg : limite manteau / noyau externe}
\begin{definition}[Discontinuité de Gutenberg]
La \cle{discontinuité de Gutenberg} se situe à une profondeur moyenne de $2\,900\ \text{km}$. Elle marque la transition brutale du manteau solide (péridotites) au noyau externe liquide (alliage Fe-Ni). Les ondes P voient leur vitesse chuter de $\approx 13,7$ à $8,1\ \text{km/s}$ ; les ondes S \textbf{disparaissent} (aucune onde S ne traverse cette interface).
\end{definition}

\begin{propriete}[Preuve de la liquidité du noyau externe]
Les ondes de cisaillement (ondes S) ne se propagent pas dans un milieu fluide (module de cisaillement $\mu = 0$). L'absence totale d'ondes S au-delà de $2\,900\ \text{km}$ (zone d'ombre des ondes S) démontre que le noyau externe se comporte comme un liquide à l'échelle des temps sismiques.
\end{propriete}

\begin{attention}[Piège fréquent : confondre « liquide » et « magma »]
Le noyau externe est un \textbf{liquide métallique} (fer-nickel) à très haute température ($4\,000$–$5\,000\ ^\circ\text{C}$) et pression ($130$–$330\ \text{GPa}$). Ce n'est \emph{pas} un magma silicate comme dans les chambres magmatiques volcaniques. Le manteau, lui, est \textbf{solide} (les ondes S le traversent) bien qu'il puisse s'écouler à très long terme (fluage).
\end{attention}

\courssub{La discontinuité de Lehmann : limite noyau externe / noyau interne}
\begin{definition}[Discontinuité de Lehmann]
À $\approx 5\,150\ \text{km}$ de profondeur, la vitesse des ondes P augmente légèrement (de $\approx 8,1$ à $11,2\ \text{km/s}$) : c'est la \cle{discontinuité de Lehmann}. Elle correspond au passage du noyau externe liquide au noyau interne solide (appelé \emph{graine}). Les ondes S réapparaissent dans le noyau interne (ondes PKJKP observées).
\end{definition}

\begin{savaistu}[Inge Lehmann, pionnière de la sismologie]
La sismologue danoise Inge Lehmann (1888–1993) a découvert cette discontinuité en 1936 en analysant les arrivées précoces d'ondes P dans la zone d'ombre. Elle a ainsi prouvé l'existence d'un noyau interne solide, révolutionnant le modèle de la Terre.
\end{savaistu}

\courssub{Les enveloppes de la Terre : description physique et chimique}
\begin{popfigure}
\begin{tikzpicture}[scale=0.45]
% Earth radius = 6370 km -> scaled to 6.37 cm
% Radii in cm (scale 1 cm = 1000 km)
% Crust: 35 km -> 0.035 cm (too thin to see). We exaggerate crust thickness for visibility.
% Let's use a scale where 1 cm = 100 km for the upper part? No, keep whole earth scale but exaggerate crust in drawing.
% Better: Draw whole earth radius 6.37. Crust 0-35km is invisible.
% Standard pedagogical diagram: Exaggerate crust and inner core sizes.
% We'll draw schematic circles with exaggerated thicknesses for clarity, but label true values.
\def\R{6.37}       % Earth radius 6370 km
\def\rcrust{6.02}   % Radius at Moho (6370-35=6335 -> 6.335). Let's use 6.02 for visual gap? No, use real scaled: 6.335.
% Let's use real scaled radii for boundaries, but draw crust as a thick ring.
% Real scaled:
% Surface: 6.37
% Moho (35km): 6.335
% Gutenberg (2900km): 3.47
% Lehmann (5150km): 1.22
% Center: 0

% Draw from center outwards so labels on top
% Inner Core (Pink)
\fill[poppinkL] (0,0) circle (1.22);
% Outer Core (Orange)
\fill[poporangeL] (0,0) circle (3.47);
% Mantle (Green)
\fill[popgreenL] (0,0) circle (6.335);
% Crust (Blue) - draw as thick ring? fill circle 6.37 then overwrite? No, fill blue last? 
% Order: Blue (whole earth), Green (mantle+core), Orange (core), Pink (inner core).
% Result: Blue ring (crust), Green ring (mantle), Orange ring (outer core), Pink disk (inner core).

\fill[popblueL] (0,0) circle (\R);
\fill[popgreenL] (0,0) circle (6.335);
\fill[poporangeL] (0,0) circle (3.47);
\fill[poppinkL] (0,0) circle (1.22);

% Labels (placed manually for readability)
\node[align=center, font=\small] at (0, 5.8) {Croûte continentale $\sim 35$ km};
\node[align=center, font=\small] at (0, 4.5) {Croûte océanique $\sim 7$ km};
\node[align=center, font=\small] at (0, 3.0) {Manteau (péridotites)};
\node[align=center, font=\small] at (0, 0.5) {Noyau externe (Fe-Ni liquide)};
\node[align=center, font=\small] at (0, -1.5) {Noyau interne (Fe-Ni solide)};

% Discontinuity arrows and labels
\draw[<->, popdark, thick] (0, 6.335) -- ++(2.8, 0) node[right, font=\small] {Moho ($\approx 35$ km)};
\draw[<->, popdark, thick] (0, 3.47) -- ++(2.8, 0) node[right, font=\small] {Gutenberg ($2\,900$ km)};
\draw[<->, popdark, thick] (0, 1.22) -- ++(2.8, 0) node[right, font=\small] {Lehmann ($5\,150$ km)};

% Scale note
\node[font=\small, popmuted] at (0, -6.8) {Échelle non respectée : épaisseurs exagérées pour la lisibilité (croûte $\times 100$)};
\end{tikzpicture}
\legende{Coupe schématique du globe terrestre montrant les trois grandes discontinuités (Moho, Gutenberg, Lehmann) et les enveloppes résultantes. Les épaisseurs ne sont pas à l'échelle (la croûte est grossie $\times 100$).}
\end{popfigure}

\textbf{La croûte terrestre}\par
\begin{itemize}
\item \textbf{Croûte continentale} : épaisseur $30$–$70\ \text{km}$ (moyenne $\approx 35\ \text{km}$), composition granitique à granodioritique (riche en Si, Al, K, Na), densité moyenne $2,7\ \text{g/cm}^3$, vitesse P $\approx 6,0$–$6,5\ \text{km/s}$ (croûte supérieure) à $6,8$–$7,2\ \text{km/s}$ (croûte inférieure).
\item \textbf{Croûte océanique} : épaisseur $5$–$10\ \text{km}$ (moyenne $\approx 7\ \text{km}$), composition basaltique (riche en Fe, Mg, Ca), densité moyenne $3,0\ \text{g/cm}^3$, vitesse P $\approx 6,7$–$7,2\ \text{km/s}$.
\end{itemize}

\begin{popfigure}
\begin{center}
\begin{tabularx}{\linewidth}{|l|X|X|}\hline
\rowcolor{popblueL}
\textbf{Caractère} & \textbf{Croûte continentale} & \textbf{Croûte océanique} \\ \hline
Épaisseur moyenne & $35$–$40\ \text{km}$ (jusqu'à $70\ \text{km}$ sous les chaînes) & $6$–$7\ \text{km}$ \\ \hline
Composition dominante & Granitoïdes (quartz, feldspaths, micas) & Basaltes (pyroxènes, plagioclases) \\ \hline
Densité & $\approx 2,7\ \text{g/cm}^3$ & $\approx 3,0\ \text{g/cm}^3$ \\ \hline
Âge moyen & Jusqu'à $4\ \text{Ga}$ (cratons) & $< 200\ \text{Ma}$ (rénovée aux dorsales) \\ \hline
Vitesse des ondes P & $6,0$–$6,5\ \text{km/s}$ (sup.) ; $6,8$–$7,2\ \text{km/s}$ (inf.) & $6,7$–$7,2\ \text{km/s}$ \\ \hline
\end{tabularx}
\end{center}
\legende{Comparatif des deux types de croûte terrestre.}
\end{popfigure}

\textbf{Le manteau}\par
Épaisseur $\approx 2\,900\ \text{km}$ (de la Moho à $2\,900\ \text{km}$ de profondeur). Roches ultramafiques (péridotites : olivine + pyroxènes). État \textbf{solide} (les ondes S s'y propagent). Toutefois, entre $\approx 100$ et $250\ \text{km}$ de profondeur, l'\cle{asthénosphère} présente une fusion partielle ($< 1\ \%$ de fond) qui abaisse la viscosité et la vitesse des ondes S (zone de basse vitesse), permettant la dérive des plaques lithosphériques.

\textbf{Le noyau externe}\par
Épaisseur $2\,900$–$5\,150\ \text{km}$ (rayon $3\,470$–$1\,220\ \text{km}$). Alliage liquide Fe-Ni ($\approx 85\ \%$ Fe, $5\ \%$ Ni, éléments légers S, O, Si). Convection vigoureuse $\rightarrow$ géodynamo $\rightarrow$ champ magnétique terrestre. Pression $130$–$330\ \text{GPa}$, température $4\,000$–$5\,000\ ^\circ\text{C}$.

\textbf{Le noyau interne (graine)}\par
Rayon $\approx 1\,220\ \text{km}$ (profondeur $5\,150$–$6\,370\ \text{km}$). Solide malgré $T \approx 5\,500\ ^\circ\text{C}$ car la pression ($330$–$360\ \text{GPa}$) stabilise la phase $\varepsilon$-Fe (compacte hexagonale). Croissance lente ($\approx 1\ \text{mm/an}$) par solidification du noyau externe, libérant de la chaleur latente et des éléments légers, ce qui alimente la convection du noyau externe.

\courssub{Tableau récapitulatif et construction du modèle en couches}
\begin{exoresolu}[Calcul du volume relatif du noyau]
\textbf{Énoncé :} Le rayon terrestre $R_T = 6\,370\ \text{km}$. Le noyau (externe + interne) a un rayon $R_N = 3\,470\ \text{km}$. Calculer le pourcentage du volume terrestre occupé par le noyau.\\
\textbf{Solution :}
\[
V_{\text{Terre}} = \frac{4}{3}\pi R_T^3,\qquad
V_{\text{noyau}} = \frac{4}{3}\pi R_N^3
\]
\[
\frac{V_{\text{noyau}}}{V_{\text{Terre}}} = \left(\frac{R_N}{R_T}\right)^3
= \left(\frac{3\,470}{6\,370}\right)^3
\approx (0,545)^3 \approx 0,162 \; (\text{soit } 16,2\%)
\]
Le noyau représente donc environ \textbf{16 \% du volume} de la Terre, alors qu'il en occupe 55 \% du rayon.
\end{exoresolu}

\begin{center}
\begin{tabular}{|l|c|c|c|c|c|}\hline
\rowcolor{popblueL}
\textbf{Enveloppe} & \textbf{Limites (km)} & \textbf{Épaisseur (km)} & \textbf{État} & \textbf{Composition principale} & \textbf{Ondes traversées} \\ \hline
Croûte continentale & $0$ – $30$–$70$ & $30$–$70$ & Solide & Granitoïdes (Si, Al, K) & P, S \\ \hline
Croûte océanique & $0$ – $5$–$10$ & $5$–$10$ & Solide & Basaltes (Fe, Mg, Ca) & P, S \\ \hline
Manteau & $30$–$70$ – $2\,900$ & $\approx 2\,830$ & Solide (asthénosphère ductile, fusion partielle $<1\%$) & Péridotites (Ol, Px) & P, S \\ \hline
Noyau externe & $2\,900$ – $5\,150$ & $2\,250$ & Liquide & Fe-Ni + éléments légers & P (ralenties), \textbf{pas de S} \\ \hline
Noyau interne & $5\,150$ – $6\,370$ & $1\,220$ & Solide & Fe-Ni (phase $\varepsilon$) & P (accélérées), S (réapparaissent) \\ \hline
\end{tabular}
\end{center}
\legende{Modèle simplifié en couches de la structure interne de la Terre (valeurs moyennes de référence PREM).}

\begin{aretenir}[Modèle de référence pour le Probatoire]
\begin{itemize}
\item Trois discontinuités majeures : Moho (croûte/manteau $\approx 35\ \text{km}$), Gutenberg (manteau/noyau externe $2\,900\ \text{km}$), Lehmann (noyau externe/noyau interne $5\,150\ \text{km}$).
\item La disparition des ondes S à $2\,900\ \text{km}$ $\Rightarrow$ noyau externe liquide.
\item La réapparition des ondes S (PKJKP) au-delà de $5\,150\ \text{km}$ $\Rightarrow$ noyau interne solide.
\item Croûte : deux types (continentale épaisse granitique, océanique mince basaltique, plus dense).
\item Manteau : solide, péridotitique, asthénosphère ductile (zone de basse vitesse pour les ondes S).
\item Noyau : externe liquide Fe-Ni (géodynamo), interne solide Fe-Ni (graine, phase $\varepsilon$).
\end{itemize}
\end{aretenir}

\begin{methode}[Construire le modèle en couches à partir des temps de parcours]
\begin{enumerate}
\item Tracer les courbes $t(\Delta)$ pour les ondes P et S observées sur un réseau de stations.
\item Repérer les changements de pente (ruptures de vitesse) et les zones d'ombre (absence d'ondes).
\item Associer chaque rupture à une discontinuité de profondeur connue (Moho $\approx 35\ \text{km}$, Gutenberg $2\,900\ \text{km}$, Lehmann $5\,150\ \text{km}$).
\item Déduire l'état physique de chaque couche : propagation des S $\Rightarrow$ solide ; absence des S $\Rightarrow$ liquide.
\item Remplir le tableau récapitulatif (limites, épaisseur, état, composition, ondes).
\end{enumerate}
\end{methode}

\begin{attention}[Erreurs à ne pas commettre au Probatoire]
\begin{itemize}
\item Dire que le manteau est « liquide » ou « en fusion » : \textbf{faux}, les ondes S le traversent (sauf zone de faible vitesse de l'asthénosphère due à une fusion partielle $<1\%$).
\item Confondre la Moho (croûte/manteau) avec la discontinuité de Gutenberg (manteau/noyau).
\item Oublier que la croûte océanique est plus dense que la croûte continentale (basalte vs granite).
\item Écrire les pressions du noyau en MPa : ce sont des \textbf{GPa} ($130$–$360\ \text{GPa}$).
\item Dire que le noyau interne est de la « phase $\alpha$-Fe cubique centrée » : c'est la phase $\varepsilon$-Fe (compacte hexagonale) stable à haute pression.
\end{itemize}
\end{attention}

\coursec{Construction du modèle en couches et limites du modèle}

Dans la section précédente, nous avons identifié les grandes discontinuités du globe terrestre — Moho, Gutenberg, Lehmann — et nous avons associé à chacune d'elles une enveloppe : la croûte, le manteau, le noyau. Mais ces connaissances ne sont pas tombées du ciel : elles résultent d'un travail de \cle{modélisation}, c'est-à-dire de la construction d'une représentation simplifiée de la réalité, bâtie à partir d'observations. Cette dernière section te montre \emph{comment} on passe des sismogrammes au modèle en couches concentriques, et surtout \emph{pourquoi} ce modèle, si utile soit-il, reste perfectible. C'est exactement la démarche que l'on te demandera de restituer au Probatoire.

\courssub{De l'observation au modèle : la démarche de modélisation}

\textbf{Une situation concrète pour comprendre.} Imagine que tu te trouves devant un avocatier du quartier, et que tu veux savoir si un avocat est mûr sans l'ouvrir. Tu le tapotes : un son sourd indique une chair molle, un son clair une chair ferme. Tu viens, sans le savoir, de faire de la \cle{sismologie} à petite échelle : tu as envoyé une onde (le tapotement) dans un objet opaque, enregistré sa réponse, et déduit une propriété interne. Les sismologues font exactement cela avec la Terre : les séismes sont les « tapotements », les stations sismiques sont les « oreilles », et le modèle en couches est la conclusion.

\begin{definition}[Modèle scientifique]
Un \cle{modèle} est une représentation simplifiée d'un objet ou d'un phénomène réel, construite à partir d'observations et de mesures, qui permet de décrire, d'expliquer et de prévoir. Un modèle n'est pas la réalité : il en est une image volontairement simplifiée, valable tant qu'il reste en accord avec les observations.
\end{definition}

Le problème scientifique posé est le suivant : \emph{la Terre est inaccessible à l'observation directe} (le forage le plus profond, celui de Kola en Russie, n'atteint que \SI{12,3}{km}, soit moins de 0,2 \% du rayon terrestre). Comment alors connaître sa structure interne ? La réponse : en exploitant les seules « messagères » capables de la traverser, les ondes sismiques.

\begin{methode}[Construire un modèle simplifié en couches à partir de données sismologiques]
La construction du modèle suit quatre étapes rigoureuses, à connaître et à savoir rédiger :
\begin{enumerate}
\item \textbf{Recueillir les temps de parcours} des ondes P et S enregistrées par un réseau de stations sismiques réparties sur le globe, pour de nombreux séismes. On construit les \cle{courbes hodochrones} (temps de parcours en fonction de la distance angulaire à l'épicentre).
\item \textbf{Repérer les discontinuités} : toute brusque modification des hodochrones (changement de pente, zone d'ombre, arrivée d'ondes réfléchies ou réfractées) traduit un changement brutal de milieu à une certaine profondeur.
\item \textbf{Déduire l'état physique des couches} à partir du comportement des ondes : les ondes S ne se propagent que dans les solides ; un ralentissement brutal des ondes P signale un changement d'état ou de composition ; une zone d'ombre des S révèle un milieu liquide.
\item \textbf{Dessiner le modèle en couches} : on représente la Terre comme un empilement de couches concentriques séparées par les discontinuités identifiées, en annotant chaque couche (nom, épaisseur, état physique) et en justifiant chaque annotation par une observation.
\end{enumerate}
\end{methode}

\courssub{Chaque propriété du modèle repose sur une observation}

\begin{propriete}[Justification obligatoire des affirmations]
Toute affirmation sur la nature ou l'état physique d'une couche terrestre doit être \cle{justifiée} par le comportement des ondes sismiques qui la traversent. Au Probatoire, une affirmation du type « le noyau externe est liquide » sans justification (« car les ondes S ne le traversent pas ») est considérée comme non démontrée et perd des points.
\end{propriete}

Le tableau ci-dessous établit la correspondance systématique entre les propriétés du modèle et les observations sismologiques qui les fondent. Apprends-le : c'est l'outil central des exercices d'exploitation de données.

\begin{center}
\begin{tabularx}{\linewidth}{|l|X|X|}
\hline
\rowcolor{popblueL} \textbf{Propriété du modèle} & \textbf{Observation sismologique} & \textbf{Interprétation} \\
\hline
Existence de la croûte (limite : Moho, $\approx$ 35 km sous les continents) & Accélération brutale des ondes P (de $\approx$ 6 à $\approx$ \SI{8}{km/s}) à faible profondeur & Changement de composition minéralogique : roches plus denses en dessous \\
\hline
Existence du manteau solide & Les ondes P \emph{et} les ondes S traversent toute la couche située entre 35 et \SI{2900}{km} & Milieu solide (les S ne se propagent que dans les solides) \\
\hline
Discontinuité de Gutenberg (\SI{2900}{km}) & Ralentissement brutal des ondes P ; disparition totale des ondes S au-delà & Passage d'un solide à un liquide \\
\hline
Noyau externe liquide (2 900 à \SI{5100}{km}) & Zone d'ombre des ondes S pour toute station située au-delà de 103\textdegree{} de l'épicentre & Les S, ondes de cisaillement, ne se propagent pas dans un liquide \\
\hline
Noyau interne (graine) solide, à partir de \SI{5100}{km} & Réaccélération des ondes P et apparition d'ondes P réfractées dans la graine (discontinuité de Lehmann) & Milieu de nouveau solide, malgré la température élevée, sous l'effet de la pression immense \\
\hline
\end{tabularx}
\end{center}

\begin{attention}[Un modèle n'est pas la réalité]
Le modèle en couches concentriques est une \emph{représentation}. Il ne prétend pas que la Terre est faite de sphères parfaitement lisses et homogènes : la surface du Moho est irrégulière, le manteau est brassé par des mouvements de convection, le noyau est en turbulence. Le modèle est \cle{valable} parce qu'il est cohérent avec l'ensemble des observations sismologiques ; il sera \cle{affiné} ou remplacé si de nouvelles observations le contredisent. Écrire « la Terre est constituée de trois couches » sans nuance est une erreur ; écrire « le modèle le plus simple de la structure terrestre comporte trois couches » est exact.
\end{attention}

\courssub{Le modèle en couches concentriques}

En appliquant la méthode en quatre étapes à l'ensemble des données sismologiques mondiales, on aboutit au modèle suivant, que tu dois savoir dessiner et légender entièrement.

\begin{popfigure}
\begin{center}
\begin{tikzpicture}[scale=0.9]
% Noyau interne
\fill[poporange!85] (0,0) circle (0.9);
% Noyau externe
\fill[popgold!70,even odd rule] (0,0) circle (1.7) (0,0) circle (0.9);
% Manteau
\fill[popgreen!45,even odd rule] (0,0) circle (2.9) (0,0) circle (1.7);
% Croûte (trait épais pour la rendre visible à l'échelle)
\draw[popdark,line width=2.2pt] (0,0) circle (2.9);
% Discontinuités
\draw[popdark,line width=0.8pt] (0,0) circle (1.7);
\draw[popdark,line width=0.8pt] (0,0) circle (0.9);
% Légendes pointées
\draw[thin,popmuted] (2.9,0) -- (4.6,0.2) node[right,align=left,font=\small]{\textbf{Croûte} (solide)\\ 5 à 70 km d'épaisseur\\ Moho à sa base};
\draw[thin,popmuted] (2.05,2.05) -- (4.0,3.1) node[right,align=left,font=\small]{\textbf{Manteau} (solide)\\ $\approx$ 2 865 km\\ jusqu'à Gutenberg (2 900 km)};
\draw[thin,popmuted] (-1.2,-1.2) -- (-3.4,-2.4) node[left,align=right,font=\small]{\textbf{Noyau externe} (liquide)\\ $\approx$ 2 200 km\\ de 2 900 à 5 100 km};
\draw[thin,popmuted] (0.3,-0.85) -- (1.6,-2.6) node[right,align=left,font=\small]{\textbf{Graine} (solide)\\ rayon $\approx$ 1 270 km\\ (Lehmann : 5 100 km)};
% Centre
\fill[popdark] (0,0) circle (0.05);
\end{tikzpicture}
\end{center}
\legende{Modèle simplifié de la structure interne de la Terre en couches concentriques. Chaque limite correspond à une discontinuité repérée par les ondes sismiques ; l'état physique de chaque couche est déduit du comportement des ondes P et S. L'échelle n'est pas respectée (la croûte est en réalité beaucoup plus fine).}
\end{popfigure}

\begin{exoresolu}[Calculer l'épaisseur des couches à partir des discontinuités]
\textbf{Énoncé.} Les discontinuités sont situées aux profondeurs suivantes : Moho à \SI{35}{km} (sous un continent), Gutenberg à \SI{2900}{km}, Lehmann à \SI{5100}{km}. Le rayon terrestre vaut $R = \SI{6370}{km}$. Calculer l'épaisseur du manteau, du noyau externe et le rayon de la graine.
\tcblower
\textbf{Solution détaillée.} L'épaisseur d'une couche est la différence entre les profondeurs des discontinuités qui l'encadrent.
\begin{align*}
\text{Manteau} &= 2900 - 35 = \SI{2865}{km} \\
\text{Noyau externe} &= 5100 - 2900 = \SI{2200}{km} \\
\text{Rayon de la graine} &= 6370 - 5100 = \SI{1270}{km}
\end{align*}
Vérification : la croûte continentale a une épaisseur de \SI{35}{km} (par définition, du sol au Moho). Le manteau représente à lui seul environ 84 \% du volume terrestre : c'est de loin la couche la plus volumineuse, alors que la croûte, seule accessible directement, est une pellicule quasi négligeable à l'échelle du globe.
\end{exoresolu}

\begin{attention}[Pièges de calcul fréquents]
Deux erreurs reviennent chaque année. Premièrement, confondre \cle{profondeur} (comptée depuis la surface) et \cle{distance au centre} : la graine commence à \SI{5100}{km} de profondeur, donc son rayon est $6370 - 5100 = \SI{1270}{km}$, et non \SI{5100}{km}. Deuxièmement, oublier de soustraire l'épaisseur de la croûte pour le manteau : le manteau ne fait pas \SI{2900}{km} d'épaisseur mais $2900 - 35 = \SI{2865}{km}$.
\end{attention}

\courssub{Complément utile au Probatoire : vers des modèles plus fins}

\textbf{(Complément — utile pour les exercices à données, mais non exigible comme restitution de cours.)} En accumulant les enregistrements et en affinant le traitement des hodochrones, les sismologues ont subdivisé les trois grandes enveloppes :

\begin{itemize}
\item le \cle{manteau supérieur} et le \cle{manteau inférieur}, séparés vers \SI{660}{km} de profondeur par une zone de transition où les ondes accélèrent progressivement (changements de structure cristalline des minéraux sous l'effet de la pression) ;
\item la \cle{lithosphère} (croûte + partie supérieure rigide du manteau, environ \SI{100}{km} d'épaisseur), rigide et fragmentée en plaques ;
\item l'\cle{asthénosphère}, zone du manteau supérieur (entre environ 100 et \SI{350}{km}) où les ondes S sont nettement ralenties : les roches y sont proches de leur point de fusion, donc \emph{ductiles} (déformables), ce qui permet le mouvement des plaques lithosphériques.
\end{itemize}

\begin{popfigure}
\begin{center}
\begin{tikzpicture}[scale=1]
% Blocs
\fill[popgreen!35] (0,0) rectangle (8,1.2); % Croûte océanique
\fill[popgreen!60] (0,1.2) rectangle (8,1.9); % Manteau sup. rigide
\fill[popblue!45] (0,1.9) rectangle (3.6,2.25); % Croûte continentale (partie sup)
\fill[popgold!70] (3.6,1.9) rectangle (8,2.25); % Croûte continentale (partie inf) - simplification visuelle
\fill[poporange!55] (0,-1.6) rectangle (8,0); % Asthénosphère
\fill[poporange!25] (0,-2.4) rectangle (8,-1.6); % Manteau inférieur
% Contours
\draw[popdark] (0,-2.4) rectangle (8,2.25);
\draw[popdark,dashed] (0,1.9) -- (8,1.9); % Base croûte cont / Moho
\draw[popdark,dashed] (0,1.2) -- (8,1.2); % Base croûte océ / Moho
\draw[popdark] (0,0) -- (8,0); % Base lithosphère
\draw[popdark,dashed] (0,-1.6) -- (8,-1.6); % Base asthénosphère
% Accolades simulées (sans bibliothèque decorations.pathreplacing)
% Lithosphère
\draw[popdark,thick] (8.15,2.25) -- (8.15,0);
\draw[popdark,thick] (8.15,2.25) -- (8.25,2.25);
\draw[popdark,thick] (8.15,0) -- (8.25,0);
\node[right=7pt,align=left,font=\small] at (8.25,1.125) {\textbf{Lithosphère}\\ (rigide, $\approx$ 100 km)};
% Asthénosphère
\draw[popdark,thick] (8.15,0) -- (8.15,-1.6);
\draw[popdark,thick] (8.15,-1.6) -- (8.25,-1.6);
\node[right=7pt,align=left,font=\small] at (8.25,-0.8) {\textbf{Asthénosphère}\\ (ductile, S ralenties)};
% Labels internes
\node[font=\small] at (1.8,2.07) {Croûte continentale};
\node[font=\small] at (5.8,2.07) {Croûte continentale};
\node[font=\small] at (1.8,0.6) {Croûte océanique};
\node[font=\small] at (5.8,0.6) {Manteau sup. rigide};
\node[font=\small] at (4,1.55) {Manteau supérieur rigide};
\node[font=\small] at (4,-0.8) {Manteau sup. partiellement ramolli};
\node[font=\small] at (4,-2.0) {Manteau inférieur (solide)};
% Flèches de mouvement des plaques
\draw[-{Stealth},very thick,popink] (1.2,2.5) -- (2.6,2.5);
\draw[-{Stealth},very thick,popink] (6.8,2.5) -- (5.4,2.5);
\node[font=\footnotesize,popink] at (4,2.75) {Mouvement des plaques lithosphériques};
\end{tikzpicture}
\end{center}
\legende{Modèle affiné lithosphère/asthénosphère (complément Probatoire). La lithosphère rigide (croûte + manteau supérieur rigide) « flotte » et coulisse sur l'asthénosphère ductile, où le ralentissement des ondes S trahit des roches proches de la fusion. Ce découpage mécanique explique la tectonique des plaques, que le modèle simple en trois couches ne décrit pas.}
\end{popfigure}

\begin{savaistu}[Le Cameroun, terrain d'illustration des limites du modèle]
La \cle{Ligne du Cameroun}, alignement volcanique qui court du golfe de Guinée (île de Bioko) jusqu'aux monts Bambouto en passant par le Mont Cameroun (\SI{4095}{m}), est l'un des rares exemples mondiaux de volcanisme intraplaque océanique et continental continu. Le Mont Cameroun, volcan actif (dernières éruptions en 1999 et 2000), n'est pas situé à une frontière de plaques : il traduirait la remontée de matériaux chauds depuis le manteau profond. Un modèle en couches parfaitement concentriques et homogènes ne peut pas expliquer ce volcanisme : il faut y ajouter des hétérogénéités internes, ce que font les modèles modernes de tomographie sismique.
\end{savaistu}

\courssub{Limites du modèle simple et statut d'un modèle}

Le modèle en couches concentriques suppose que chaque couche est \emph{homogène} et que ses propriétés ne varient qu'avec la profondeur. Or les observations montrent des \cle{hétérogénéités latérales} que le modèle simple ne rend pas compte :

\begin{itemize}
\item la croûte n'a pas la même épaisseur partout : environ 5 à \SI{10}{km} sous les océans, 30 à \SI{70}{km} sous les continents (jusqu'à \SI{70}{km} sous les chaînes de montagnes) ;
\item la lithosphère est fragmentée en \cle{plaques} en mouvement, avec des zones de subduction où des matériaux froids plongent dans le manteau ;
\item des \cle{points chauds} (comme la Ligne du Cameroun ou Hawaï) témoignent de remontées localisées de matériaux mantelliques ;
\item la tomographie sismique révèle d'immenses anomalies de vitesse des ondes à la base du manteau, incompatibles avec des couches parfaitement régulières.
\end{itemize}

\begin{aretenir}[Ce qu'il faut retenir de la démarche de modélisation]
\begin{itemize}
\item Le modèle en couches concentriques (croûte – manteau – noyau externe liquide – graine solide) se \emph{construit} en quatre étapes : temps de parcours $\rightarrow$ discontinuités $\rightarrow$ états physiques $\rightarrow$ schéma.
\item Chaque propriété du modèle est \emph{justifiée} par une observation sismologique précise (P ralenties, S arrêtées, zones d'ombre).
\item Un modèle est une représentation \emph{simplifiée et perfectible} : il est valide tant qu'il est en accord avec les observations, et il est affiné (lithosphère/asthénosphère, tomographie) quand de nouvelles données apparaissent.
\end{itemize}
\end{aretenir}

\begin{exercice}[Construire et critiquer un modèle]
Un séisme est enregistré par des stations du monde entier. On observe : (a) les ondes P et S arrivent à toutes les stations situées à moins de 103\textdegree{} de l'épicentre ; (b) au-delà de 103\textdegree{}, plus aucune onde S n'arrive et les ondes P directes disparaissent entre 103\textdegree{} et 142\textdegree{} ; (c) certaines stations au-delà de 142\textdegree{} reçoivent des ondes P légèrement accélérées ayant traversé le centre.
\begin{enumerate}
\item Construis, en suivant les quatre étapes de la méthode, le modèle de structure interne compatible avec ces données (schéma légendé exigé).
\item Justifie l'état physique attribué à chaque couche.
\item Ce modèle peut-il expliquer l'existence du volcanisme du Mont Cameroun ? Conclus sur le statut du modèle.
\end{enumerate}
\end{exercice}

\begin{corrige}[Exercice : construire et critiquer un modèle]
\textbf{1.} Étape 1 : les temps de parcours montrent des arrivées normales jusqu'à 103\textdegree{}, puis des anomalies. Étape 2 : la limite à 103\textdegree{} repère une discontinuité majeure (Gutenberg, \SI{2900}{km}) ; les P accélérées au-delà de 142\textdegree{} révèlent une seconde discontinuité (Lehmann, \SI{5100}{km}). Étape 3 : les S traversent l'enveloppe externe $\rightarrow$ manteau solide ; les S disparaissent au-delà de 103\textdegree{} $\rightarrow$ noyau externe liquide ; les P réaccélèrent au centre $\rightarrow$ graine solide. Étape 4 : schéma de trois couches concentriques (croûte, manteau, noyau externe + graine), identique à la figure de cours.
\textbf{2.} Manteau solide : P \emph{et} S s'y propagent. Noyau externe liquide : les S, ondes de cisaillement, ne s'y propagent pas (zone d'ombre des S) et les P y sont ralenties. Graine solide : les P y sont accélérées, ce qui indique un milieu de rigidité élevée.
\textbf{3.} Non : ce modèle suppose des couches homogènes et ne contient ni plaques ni remontées localisées de matériaux chauds. Il ne peut donc pas expliquer un volcanisme intraplaque comme celui du Mont Cameroun. Cela illustre le statut du modèle : représentation simplifiée, valable pour décrire la structure globale, mais perfectible et à compléter (modèles de convection mantellique, tomographie sismique) dès que les observations l'exigent.
\end{corrige}

\newpage

\coursec{Activité d'intégration}

\coursec{ND14 : Conception d'un modèle simple de la structure interne de la Terre}

\courssub{Objectifs de l'activité}
Cette activité d'intégration vise à mobiliser l'ensemble des savoirs et savoir-faire de la leçon pour :
\begin{itemize}
    \item \textbf{Exploiter} des données de temps de parcours d'ondes sismiques (ondes P directes et réfléchies) afin de \textbf{localiser une discontinuité majeure} : le Moho (discontinuité de Mohorovičić) sous le volcan actif du mont Cameroun.
    \item \textbf{Comparer} l'épaisseur de la croûte continentale et océanique pour en déduire la structure hétérogène du globe.
    \item \textbf{Analyser} l'existence d'une zone d'ombre des ondes S pour un séisme lointain et \textbf{en déduire l'état physique du noyau externe}.
    \item \textbf{Construire}, à l'échelle, un \textbf{modèle simplifié en couches} de la structure interne de la Terre, en justifiant chaque enveloppe par des observations sismologiques.
\end{itemize}

\courssub{Situation : À la recherche du Moho sous le mont Cameroun}
Le mont Cameroun (\emph{Mongo ma Ndemi}), stratovolcan actif culminant à 4\,095\,m d'altitude près de Buéa (région du Sud-Ouest), est le siège d'une sismicité volcanique et tectonique régulière. Le 14 mars 2023, un séisme de faible magnitude ($M_L = 3,2$) se produit à une profondeur focale de 10\,km, à quelques kilomètres au large des côtes, au sud-ouest du volcan. 

Trois stations sismologiques du réseau national (Buéa \textsc{bua}, Douala \textsc{dla}, Yaoundé \textsc{yau}) enregistrent les ondes issues de ce séisme. L'épicentre est situé à une distance épicentrale $\Delta \approx 15\,\text{km}$ de la station \textsc{bua}, $\approx 75\,\text{km}$ de \textsc{dla} et $\approx 300\,\text{km}$ de \textsc{yau}.

Les sismogrammes (composante verticale) montrent clairement, pour la station \textsc{bua} la plus proche :
\begin{enumerate}
    \item Une première arrivée nette : l'\textbf{onde P directe} (trajet dans la croûte).
    \item Une arrivée secondaire, environ 11,7\,s plus tard : l'\textbf{onde P réfléchie sur le Moho} (notée \emph{PmP}), qui a traversé la croûte en aller-retour selon une incidence quasi verticale (géométrie source-récepteur proche de la verticale).
\end{enumerate}

Par ailleurs, l'analyse d'un séisme lointain (télésisme, $\Delta > 100^\circ$) enregistré par le même réseau montre une \textbf{disparition brutale des ondes S} au-delà d'une distance angulaire de $103^\circ$ par rapport à l'épicentre, tandis que les ondes P restent détectées (avec une zone d'ombre P entre $103^\circ$ et $143^\circ$).

\textbf{Données numériques fournies :}
\begin{center}
\begin{tabularx}{\linewidth}{|l|X|}\hline
\textbf{Paramètre} & \textbf{Valeur} \\ \hline
Vitesse moyenne des ondes P dans la croûte continentale ($V_{P,\text{croûte}}$) & $6,0\,\text{km}\cdot\text{s}^{-1}$ \\ \hline
Temps aller-retour de l'onde PmP à la station \textsc{bua} ($t_{\text{aller-retour}}$) & $11,7\,\text{s}$ \\ \hline
Temps aller-retour de l'onde PmP pour un profil océanique de référence ($t_{\text{océan}}$) & $1,7\,\text{s}$ \\ \hline
Rayon moyen de la Terre ($R_T$) & $6\,371\,\text{km}$ \\ \hline
Profondeur de la discontinuité de Gutenberg (mantelle/noyau) & $2\,890\,\text{km}$ \\ \hline
Profondeur de la discontinuité de Lehmann (noyau externe/interne) & $5\,150\,\text{km}$ \\ \hline
Échelle de construction du modèle & $1\,\text{cm} = 500\,\text{km}$ \\ \hline
\end{tabularx}
\end{center}

\courssub{Consignes}
\emph{Travail en binôme. Réponses attendues sur papier millimétré pour la question 4.}

\begin{enumerate}
    \item \textbf{Profondeur du Moho sous le mont Cameroun.} 
    En utilisant la vitesse des ondes P dans la croûte et le temps aller-retour de l'onde réfléchie \emph{PmP} enregistré à Buéa, calculer la profondeur $h$ de la discontinuité de Mohorovičić (Moho) sous le volcan. Exprimer le résultat en kilomètres avec une décimale.
    
    \item \textbf{Comparaison croûte continentale / croûte océanique.} 
    En reprenant le même raisonnement avec le temps aller-retour de référence pour un domaine océanique ($t = 1,7\,\text{s}$), calculer l'épaisseur de la croûte océanique. Comparer les deux valeurs et \textbf{en déduire la nature de la croûte} sous le mont Cameroun.
    
    \item \textbf{État du noyau externe : analyse de la zone d'ombre des ondes S.} 
    Un télésisme est enregistré par le réseau. Les ondes S ne sont pas détectées au-delà de $103^\circ$ de distance épicentrale.
    \begin{enumerate}
        \item Rappeler la propriété fondamentale des ondes S (ondes de cisaillement) vis-à-vis des milieux traversés.
        \item Expliquer pourquoi l'existence d'une \textbf{zone d'ombre des ondes S} (absence de réception au-delà de $103^\circ$) implique que le noyau externe est à l'état \textbf{liquide}.
        \item Préciser le nom de la discontinuité qui marque la limite mantelle / noyau externe.
    \end{enumerate}
    
    \item \textbf{Construction du modèle en couches de la Terre à l'échelle.} 
    Sur une feuille de papier millimétré (format A4 paysage recommandé), tracer, à l'échelle $1\,\text{cm} = 500\,\text{km}$, une coupe méridienne du globe terrestre (cercle de rayon $R_T$). 
    \begin{enumerate}
        \item Calculer le rayon du cercle à tracer (en cm).
        \item Calculer les rayons (en cm) des trois cercles concentriques matérialisant : le Moho (profondeur trouvée en Q1), la discontinuité de Gutenberg, la discontinuité de Lehmann.
        \item Tracer les quatre enveloppes (croûte, mantelle, noyau externe, noyau interne) en les coloriant différemment. Légender chaque couche (nom, état physique, épaisseur approximative en km) et chaque discontinuité (nom, profondeur).
        \item Justifier brièvement, en regard du schéma, l'attribution de l'état solide/liquide à chaque enveloppe en s'appuyant sur les observations sismologiques de la leçon (vitesses, réflexion, réfraction, zones d'ombre).
    \end{enumerate}
\end{enumerate}

\courssub{Ressources à mobiliser}
\begin{itemize}
    \item \textbf{Relation fondamentale :} $v = \dfrac{d}{t}$ $\Rightarrow$ pour un aller-retour vertical : $2h = v \times t_{\text{aller-retour}}$ $\Rightarrow$ $h = \dfrac{v \times t}{2}$.
    \item \textbf{Propriétés des ondes sismiques :}
    \begin{itemize}
        \item Ondes P (primaires, de compression) : se propagent dans les solides \textbf{et} les liquides.
        \item Ondes S (secondaires, de cisaillement) : se propagent \textbf{uniquement dans les solides} (module de cisaillement nul dans les fluides).
    \end{itemize}
    \item \textbf{Comportement aux discontinuités :} Réflexion (retour dans le même milieu) et Réfraction (changement de direction/vitesse selon la loi de Snell-Descartes $\frac{\sin i}{v} = \text{cste}$) si contraste d'impédance acoustique.
    \item \textbf{Zones d'ombre :} Conséquence de la réfraction aux interfaces entre milieux de vitesses différentes. Zone d'ombre S $\rightarrow$ milieu liquide (noyau externe). Zone d'ombre P ($103^\circ-143^\circ$) $\rightarrow$ forte réfraction à l'entrée dans le noyau liquide (ralentissement brutal des ondes P).
    \item \textbf{Échelle graphique :} $1\,\text{cm} \leftrightarrow 500\,\text{km}$. Rayon Terre $R_T = 6\,371\,\text{km}$.
\end{itemize}

\courssub{Démarche et corrigé commenté}
\begin{exoresolu}[Corrigé détaillé de l'activité d'intégration]
\tcblower
\textbf{1. Profondeur du Moho sous le mont Cameroun (station \textsc{bua})}
\par
L'onde \emph{PmP} effectue un trajet aller-retour vertical dans la croûte (incidence quasi normale car source et récepteur sont très proches). La distance parcourue est $2h$.
\begin{align*}
    h &= \frac{V_{P,\text{croûte}} \times t_{\text{aller-retour}}}{2} \\
      &= \frac{6,0\,\text{km}\cdot\text{s}^{-1} \times 11,7\,\text{s}}{2} \\
      &= \frac{70,2\,\text{km}}{2} \\
      &= 35,1\,\text{km}
\end{align*}
\par
\cle{Résultat} : Le Moho se situe à une profondeur de \textbf{35,1\,km} sous le mont Cameroun. C'est une valeur typique de \textbf{croûte continentale épaisse}, cohérente avec un contexte de rift volcanique (ligne du Cameroun) où la croûte est amincie par rapport aux cratons anciens (40-45\,km) mais reste bien plus épaisse que la croûte océanique.

\textbf{2. Comparaison avec la croûte océanique}
\par
Même formule, vitesse $V_P$ supposée identique en moyenne dans la croûte supérieure (gabbro/basalte $\approx 6,0-6,7\,\text{km/s}$, on conserve $6,0\,\text{km/s}$ pour la comparaison) :
\begin{align*}
    h_{\text{océan}} &= \frac{6,0 \times 1,7}{2} = \frac{10,2}{2} = 5,1\,\text{km}
\end{align*}
\par
\cle{Comparaison} : 
\begin{itemize}
    \item Croûte sous le mont Cameroun (continentale) : $\approx 35\,\text{km}$.
    \item Croûte océanique (référence) : $\approx 5\,\text{km}$.
\end{itemize}
\par
\cle{Interprétation} : La croûte sous le mont Cameroun est \textbf{7 fois plus épaisse} que la croûte océanique. Elle est composée de roches felsiques à intermédiaires (granites, gneiss, roches volcaniques différenciées) moins denses ($\rho \approx 2,7-2,9$), flottant plus haut sur la mantelle (isostasie), contrairement à la croûte océanique basaltique dense ($\rho \approx 3,0$) et mince. Cela confirme que le mont Cameroun est édifié sur \textbf{croûte continentale} (ligne du Cameroun = chaîne volcanique intraplaque sur croûte continentale).

\textbf{3. État du noyau externe : zone d'ombre des ondes S}
\begin{enumerate}
    \item \textbf{Propriété des ondes S :} Ce sont des ondes de \textbf{cisaillement} (déformation transversale). Elles nécessitent un milieu possédant une \textbf{rigidité} (module de cisaillement $\mu > 0$). Les fluides (liquides, gaz) ont $\mu = 0$ : ils ne résistent pas au cisaillement. \textbf{Les ondes S ne se propagent pas dans les liquides.}
    \item \textbf{Raisonnement :} L'absence d'ondes S au-delà de $103^\circ$ (zone d'ombre S) signifie que les ondes S émises par le foyer n'ont pas pu traverser la région centrale du globe pour atteindre ces stations. Le trajet direct des ondes S vers ces distances passe nécessairement par le \textbf{noyau}. Comme les ondes S sont bloquées, le noyau (ou une partie majeure) doit être \textbf{liquide}. C'est la preuve historique (Oldham 1906, Gutenberg 1914) que le \textbf{noyau externe est liquide}.
    \item \textbf{Discontinuité :} La limite mantelle / noyau externe est la \textbf{discontinuité de Gutenberg} (à $\approx 2\,890\,\text{km}$ de profondeur). On y observe une chute brutale de la vitesse des ondes P (de $\approx 13,7$ à $\approx 8,1\,\text{km/s}$) et la disparition des ondes S.
\end{enumerate}

\textbf{4. Construction du modèle en couches à l'échelle $1\,\text{cm} = 500\,\text{km}$}
\par
\textbf{Calculs préliminaires (rayons en cm) :}
\begin{center}
\begin{tabular}{|l|c|c|c|}\hline
\textbf{Enveloppe / Discontinuité} & \textbf{Profondeur (km)} & \textbf{Rayon depuis le centre (km)} & \textbf{Rayon à l'échelle (cm)} \\ \hline
Surface (Rayon Terre $R_T$) & 0 & 6\,371 & $6371 / 500 = \mathbf{12,74}$ \\ \hline
\textbf{Moho} (Q1) & 35,1 & $6371 - 35,1 = 6\,335,9$ & $6335,9 / 500 = \mathbf{12,67}$ \\ \hline
\textbf{Gutenberg} (Mantelle/Noyau externe) & 2\,890 & $6371 - 2890 = 3\,481$ & $3481 / 500 = \mathbf{6,96}$ \\ \hline
\textbf{Lehmann} (Noyau externe/Interne) & 5\,150 & $6371 - 5150 = 1\,221$ & $1221 / 500 = \mathbf{2,44}$ \\ \hline
Centre & 6\,371 & 0 & 0 \\ \hline
\end{tabular}
\end{center}

\par
\textbf{Observation importante :} L'épaisseur de la croûte à l'échelle est $12,74 - 12,67 = 0,07\,\text{cm} = 0,7\,\text{mm}$. Elle est \textbf{quasi invisible} sur un dessin à cette échelle (trait de crayon fin $\approx 0,3-0,5\,\text{mm}$). Il est conseillé de la représenter par un trait très fin ou de la mentionner en légende avec une flèche d'appel (« Croûte continentale $\approx 35\,\text{km}$ »).

\par
\textbf{Réalisation graphique (code TikZ pour illustration du corrigé) :}
\begin{popfigure}
\begin{tikzpicture}[scale=1.0]
    % Couleurs douces
    \definecolor{crustcol}{RGB}{210,180,140}    % sable/brun clair
    \definecolor{mantlecol}{RGB}{255,180,100}   % orange mantelle
    \definecolor{outercorecol}{RGB}{255,100,100}% rouge liquide
    \definecolor{innercorecol}{RGB}{200,0,0}    % rouge foncé solide
    
    % Remplissage des couches (du centre vers l'extérieur pour superposition correcte)
    % Rayons calculés : Surface=12.742, Moho=12.672, Gutenberg=6.962, Lehmann=2.442
    % Noyau interne
    \fill[innercorecol] (0,0) circle (2.442);
    % Noyau externe
    \fill[outercorecol] (0,0) circle (6.962);
    % Mantelle
    \fill[mantlecol] (0,0) circle (12.672);
    % Croûte (très fine, on la dessine en dernier comme un anneau fin ou juste le cercle surface)
    \fill[crustcol] (0,0) circle (12.742);
    
    % Cercles de contours
    \draw[very thick, popdark] (0,0) circle (12.742); % Surface
    \draw[thick, dashed, popdark] (0,0) circle (12.672); % Moho
    \draw[thick, popdark] (0,0) circle (6.962); % Gutenberg
    \draw[thick, popdark] (0,0) circle (2.442); % Lehmann
    
    % Légendes avec lignes de rappel (callouts)
    % Croûte (flèche vers l'extérieur)
    \draw[<-, thick, popdark] (12.7, 0.2) -- ++(1.5, 0.8) node[right, align=left, font=\small] {Croûte continentale\\ $\approx 35$ km (0,7 mm)\\ Solide, $V_P \approx 6$ km/s};
    % Moho
    \draw[<-, thick, popdark] (12.67, -0.5) -- ++(1.5, -1.0) node[right, align=left, font=\small] {Moho (35 km)\\ Discontinuité $V_P$ : 6 $\to$ 8 km/s};
    % Mantelle
    \draw[<-, thick, popdark] (9, 0) -- ++(2.0, 1.5) node[right, align=left, font=\small] {Mantelle\\ $\approx 2\,855$ km d'épaisseur\\ Solide (rhéologie ductile)\\ $V_P$ : 8 $\to$ 13,7 km/s};
    % Gutenberg
    \draw[<-, thick, popdark] (6.96, -1.5) -- ++(1.5, -1.5) node[right, align=left, font=\small] {Gutenberg (2\,890 km)\\ Mantelle $\to$ Noyau externe\\ $V_P$ chute : 13,7 $\to$ 8,1 km/s\\ Ondes S : STOP};
    % Noyau externe
    \draw[<-, thick, popdark] (4, -3) -- ++(1.5, -1.0) node[right, align=left, font=\small] {Noyau externe liquide\\ $\approx 2\,260$ km\\ Fe-Ni fondu\\ $V_P \approx 8,1 \to 10,3$ km/s};
    % Lehmann
    \draw[<-, thick, popdark] (2.44, 1.5) -- ++(-2.5, 1.5) node[left, align=right, font=\small] {Lehmann (5\,150 km)\\ Noyau ext. $\to$ int.\\ $V_P$ saut : 10,3 $\to$ 11,2 km/s};
    % Noyau interne
    \draw[<-, thick, popdark] (0, 2.44) -- ++(-2.5, 1.5) node[left, align=right, font=\small] {Noyau interne solide\\ Rayon $\approx 1\,220$ km\\ Fe-Ni solide};
    
    % Échelle graphique
    \draw[|-|, thick, popdark] (-14.242, -13.542) -- ++(2, 0) node[midway, below, font=\small] {500 km = 1 cm};
    \node[below, font=\small] at (0, -14.242) {Échelle : 1 cm = 500 km};
\end{tikzpicture}
\legende{Modèle simplifié de la structure interne de la Terre à l'échelle 1 cm = 500 km. La croûte continentale (35 km) est quasi invisible (0,7 mm). Les trois discontinuités majeures (Moho, Gutenberg, Lehmann) séparent quatre enveloppes aux états physiques et vitesses sismiques distincts.}
\end{popfigure}

\par
\textbf{Justification des enveloppes par les observations sismologiques :}
\begin{center}
\begin{tabularx}{\linewidth}{|l|X|X|}\hline
\textbf{Enveloppe} & \textbf{Observation sismologique clé} & \textbf{Déduction (État / Composition)} \\ \hline
\textbf{Croûte} & Réflexion nette PmP (ondes P) ; $V_P \approx 5,5-6,5\,\text{km/s}$ ; $V_S \approx 3,2-3,7\,\text{km/s}$ & \textbf{Solide}. Roches silicatées différenciées (granitoïdes, basaltes). Discontinuité Moho = saut de $V_P$ (contraste croûte/mantelle). \\ \hline
\textbf{Mantelle} & Ondes P et S propagées ; $V_P$ croît avec la profondeur ($8 \to 13,7\,\text{km/s}$) ; pas de zone d'ombre S dans la mantelle. & \textbf{Solide} (mais ductile à long terme / haute T,P). Péridotites (olivine, pyroxènes). Transitions de phase (410, 660 km) = sauts de vitesse. \\ \hline
\textbf{Noyau externe} & \textbf{ZONE D'OMBRE S} ($>103^\circ$) ; $V_P$ chute brutalement ($13,7 \to 8,1\,\text{km/s}$) à Gutenberg ; $V_P$ recroît ensuite. & \textbf{Liquide}. Alliage Fe-Ni (+ éléments légers S, O, Si). Arrêt des ondes S ($\mu=0$). Réfraction forte des ondes P (ralentissement). \\ \hline
\textbf{Noyau interne} & Réapparition des ondes P (PKIKP) ; saut de $V_P$ à Lehmann ($10,3 \to 11,2\,\text{km/s}$) ; anisotropie (vitesse plus rapide axe polaire). & \textbf{Solide}. Fe-Ni cristallisé (pression extrême $> 360$ GPa). Croissance par solidification du noyau externe (libération chaleur latente $\to$ géodynamo). \\ \hline
\end{tabularx}
\end{center}
\end{exoresolu}

\courssub{Bilan de l'activité}
Cette activité d'intégration a permis de \textbf{valider la démarche scientifique de la sismologie globale} : partir de l'observation de signaux (temps de parcours, zones d'ombre) pour construire un modèle physique cohérent de l'intérieur inaccessible de la Terre.

\begin{aretenir}[Acquis majeurs de l'activité]
\begin{itemize}
    \item \textbf{Méthode du temps de parcours (réflexion verticale) :} $h = \frac{v \cdot t}{2}$ permet d'accéder à la profondeur des discontinuités peu profondes (Moho). Application concrète : \textbf{35\,km sous le mont Cameroun} (croûte continentale) vs \textbf{5\,km en domaine océanique}.
    \item \textbf{Preuve de la liquidité du noyau externe :} L'\textbf{absence d'ondes S} au-delà de $103^\circ$ (zone d'ombre S) est la preuve directe que le noyau externe ne transmet pas le cisaillement $\Rightarrow$ \textbf{état liquide}.
    \item \textbf{Preuve de la solidité du noyau interne :} La détection d'ondes P ayant traversé le centre (phase PKIKP) et le saut de vitesse à la discontinuité de Lehmann ($5\,150\,\text{km}$) imposent un milieu solide au centre.
    \item \textbf{Modèle en 4 enveloppes :} Croûte (solide, hétérogène), Mantelle (solide, homogène chimiquement mais transitions de phase), Noyau externe (liquide, Fe-Ni), Noyau interne (solide, Fe-Ni). Les trois discontinuités majeures (Moho, Gutenberg, Lehmann) correspondent à des contrastes chimiques (Moho, Lehmann) ou chimico-physiques (Gutenberg) majeurs.
    \item \textbf{Limite du modèle :} L'échelle $1/50\,000\,000$ (1 cm = 500 km) rend la croûte invisible (sub-mm). Le modèle est \emph{simplifié} : il ne montre pas les transitions de phase de la mantelle (410, 660 km), ni l'hétérogénéité latérale (panaches, plaques subductées), ni l'anisotropie du noyau interne.
\end{itemize}
\end{aretenir}

\begin{attention}[Pièges fréquents à éviter au Probatoire]
\begin{itemize}
    \item \textbf{Oublier le facteur 1/2} dans le calcul de la profondeur par réflexion ($h = v t / 2$ et non $v t$).
    \item \textbf{Confondre discontinuité de Gutenberg} (mantelle/noyau externe, 2\,890 km, \textbf{arrêt des ondes S}) \textbf{et discontinuité de Lehmann} (noyau externe/interne, 5\,150 km, \textbf{réapparition ondes P / saut vitesse}).
    \item \textbf{Affirmer que « les ondes P ne traversent pas le noyau »} : FAUX. Elles le traversent mais sont réfractées (ralenties dans le liquide), créant une zone d'ombre P ($103^\circ-143^\circ$).
    \item \textbf{Négliger l'échelle graphique :} Dessiner la croûte avec une épaisseur visible (ex: 1 cm) au lieu de 0,7 mm. Il faut le mentionner explicitement : « Croûte non visible à cette échelle, épaisseur réelle 35 km ».
    \item \textbf{Confondre état physique et composition :} La mantelle est \textbf{solide} (ondes S passent) mais se déforme plastiquement (ductile) à l'échelle géologique. Le noyau externe est \textbf{liquide} (ondes S bloquées).
\end{itemize}
\end{attention}

\begin{savaistu}[Le mont Cameroun, fenêtre sur la structure profonde]
Le mont Cameroun est l'un des rares volcans africains où l'on peut observer directement la transition croûte-manteau grâce aux \textbf{xénolites} (enclaves de roches mantelliques remontées par les laves basaltiques). Les basaltes du mont Cameroun contiennent fréquemment des \textbf{péridotites} (lherzolites, harzburgites) qui sont des fragments de la mantelle supérieure arrachés lors de la remontée du magma. Leur étude pétrologique (thermobarométrie) confirme des températures de $900-1100^\circ\text{C}$ et des pressions correspondant à $30-40\,\text{km}$ de profondeur, en \textbf{parfait accord avec la profondeur sismique du Moho (35 km)} calculée dans cette activité. C'est une magnifique convergence entre \textbf{sismologie} (imagerie indirecte) et \textbf{pétrologie} (échantillon direct) !
\end{savaistu}

\newpage

\coursec{Méthodes et exercices résolus}

\coursec{Leçon ND14 : Conception d'un modèle simple de la structure interne de la Terre}

\courssub{Séismes et ondes sismiques : origine et propagation}

\begin{definition}[Séisme, foyer, épicentre]
Un \textbf{séisme} (ou tremblement de terre) est une secousse brutale du sol due à une rupture brutale de roches dans la lithosphère (faillesse, subduction, volcanisme).
\begin{itemize}
    \item \textbf{Foyer (hypocentre) :} Point de rupture initial à l'intérieur de la Terre (profondeur $h$).
    \item \textbf{Épicentre :} Projection verticale du foyer à la surface de la Terre.
    \item \textbf{Distance épicentrale $d$ :} Distance à la surface entre l'épicentre et une station sismique (en km).
    \item \textbf{Distance angulaire $\Delta$ :} Angle au centre de la Terre $\widehat{ECS}$ ($E$ épicentre, $C$ centre, $S$ station), en degrés. $d = R_T \times \Delta_{\text{rad}}$ ($R_T = 6371$ km).
\end{itemize}
\end{definition}

\begin{propriete}[Ondes de volume : ondes P et ondes S]
À partir du foyer, se propagent deux types d'ondes de volume (corps) dans le milieu continu :
\begin{itemize}
    \item \textbf{Ondes P (Primaires, de compression) :} Ondes longitudinales. Le mouvement des particules est \textbf{parallèle} à la direction de propagation. Elles compriment et dilatent le milieu (module d'incompressibilité $K$). Vitesse $v_P = \sqrt{\frac{K + \frac{4}{3}\mu}{\rho}}$. Elles traversent \textbf{solides et liquides}. Ce sont les plus rapides ($v_P \approx 6$ à $13,7$ km/s dans le manteau).
    \item \textbf{Ondes S (Secondaires, de cisaillement) :} Ondes transversales. Le mouvement des particules est \textbf{perpendiculaire} à la direction de propagation (polarisation verticale $SV$ ou horizontale $SH$). Elles déforment le milieu sans changement de volume (module de rigidité $\mu$). Vitesse $v_S = \sqrt{\frac{\mu}{\rho}}$. Elles ne traversent \textbf{que les solides} ($\mu=0$ dans un fluide $\Rightarrow v_S=0$). Plus lentes ($v_S \approx 3,5$ à $7,3$ km/s dans le manteau).
\end{itemize}
\cle{Rapport $v_P/v_S > \sqrt{4/3} \approx 1,15$ (souvent $\approx 1,73$ dans la croûte/manteau).}
\end{propriete}

\begin{experience}[Mise en évidence des ondes P et S sur un sismogramme]
\textbf{Matériel :} Sismomètre à 3 composantes (Z, N, E), chronométrage GPS. \textbf{Protocole :} Enregistrer le signal d'un séisme local (ex. : tir de carrière ou séisme naturel au Cameroun). \textbf{Observations :}
\begin{enumerate}
    \item Première arrivée nette sur la composante verticale (Z) : \textbf{onde P} (mouvement vertical de compression).
    \item Arrivée suivante, plus ample, sur les composantes horizontales (N, E) et Z : \textbf{onde S} (cisaillement).
    \item Arrivées tardives, très amples, basse fréquence : ondes de surface (Rayleigh, Love).
\end{enumerate}
\textbf{Interprétation :} L'intervalle de temps $\Delta T = T_S - T_P$ (temps S moins temps P) croît avec la distance épicentrale $d$. C'est la base de la localisation.
\end{experience}

\begin{popfigure}
\begin{tikzpicture}[scale=0.8]
% Onde P
\draw[popblue, thick, ->] (0,0) -- (6,0) node[midway, above] {Direction de propagation};
\draw[popblue, thick, decorate, decoration={coil, aspect=0.3, segment length=4pt, amplitude=3pt}] (0,0) -- (6,0);
\foreach \x in {0.5,1.5,2.5,3.5,4.5,5.5} {
    \draw[popblue, <->] (\x,-0.5) -- (\x+0.3,-0.5) node[midway, below, font=\tiny] {Compression};
    \draw[popblue, <->] (\x+0.3,-0.5) -- (\x+0.6,-0.5) node[midway, below, font=\tiny] {Dilatation};
}
\node[popblue, font=\bfseries] at (3,-1.5) {Onde P (Longitudinale)};

% Onde S
\begin{scope}[yshift=-4cm]
\draw[poporange, thick, ->] (0,0) -- (6,0) node[midway, above] {Direction de propagation};
\draw[poporange, thick] (0,0) sin (1,1) cos (2,0) sin (3,-1) cos (4,0) sin (5,1) cos (6,0);
\foreach \x in {0.5,1.5,2.5,3.5,4.5,5.5} {
    \draw[poporange, <->] (\x,0) -- (\x,0.8) node[midway, right, font=\tiny] {Cisaillement};
}
\node[poporange, font=\bfseries] at (3,-1.5) {Onde S (Transversale, polarisée verticale SV)};
\end{scope}
\end{tikzpicture}
\legende{Modélisation simplifiée du mouvement des particules pour une onde P (compression/dilatation) et une onde S (cisaillement vertical SV).}
\end{popfigure}

\courssub{Comportement des ondes sismiques aux discontinuités}

\begin{propriete}[Réflexion, réfraction et conversion de mode]
Lorsqu'une onde sismique rencontre une interface entre deux milieux de propriétés élastiques différentes (discontinuité), elle génère :
\begin{itemize}
    \item Une \textbf{onde réfléchie} (retour dans le milieu incident) : angle de réflexion = angle d'incidence.
    \item Une \textbf{onde réfractée} (transmise dans le second milieu) : loi de \textbf{Snell-Descartes} : $\frac{\sin i_1}{v_1} = \frac{\sin i_2}{v_2} = p$ (paramètre de rayon, constant le long du trajet).
    \item Une \textbf{conversion de mode} : une onde P incidente peut générer une onde S réfléchie/réfractée (et vice-versa) car les conditions de continuité des déplacements et contraintes au contact le nécessitent.
\end{itemize}
\cle{Angle critique $i_c$ :} Si $v_2 > v_1$, pour $i_1 = i_c = \arcsin(v_1/v_2)$, l'onde réfractée glisse le long de l'interface ($i_2 = 90^\circ$). Au-delà ($i_1 > i_c$), \textbf{réflexion totale} (onde réfractée évanescente).
\end{propriete}

\begin{definition}[Zones d'ombre sismiques]
Une \textbf{zone d'ombre} est une région à la surface où aucune onde directe (P ou S) n'est enregistrée.
\begin{itemize}
    \item \textbf{Zone d'ombre des ondes S ($\Delta > 103^\circ \sim 105^\circ$) :} Les ondes S ne traversent pas le noyau externe (liquide, $\mu=0$). Elles sont arrêtées à la discontinuité de Gutenberg.
    \item \textbf{Zone d'ombre des ondes P ($103^\circ < \Delta < 143^\circ$) :} À la limite manteau/noyau, $v_P$ chute brutalement ($13,7 \rightarrow 8,1$ km/s). Les rayons P se réfractent \textbf{vers la normale} (vers le centre), créant un "vide" d'illumination à la surface. Au-delà de $143^\circ$, les ondes $PKP$ (traversant le noyau externe) et $PKIKP$ (traversant le noyau interne) réapparaissent.
\end{itemize}
\end{definition}

\begin{popfigure}
\begin{tikzpicture}[scale=0.55]
% Terre
\draw[popline, very thick] (0,0) circle (6); % Surface
\draw[popblue, thick] (0,0) circle (3.3); % Noyau (3470/6371*6 = 3.27)
% Épicentre en haut
\coordinate (E) at (0,6);
\node[popdark, font=\bfseries] at (0,6.8) {Épicentre};
% Rayons P
\draw[popblue, thick, ->] (E) -- ++(-1.5,-2) node[midway, left] {$P$};
\draw[popblue, thick, ->] (E) -- ++(-2.5,-3.5);
% Rayon P tangent au noyau (limite zone d'ombre)
\draw[popblue, dashed, thick] (E) -- ++(-3.2,-4.5) node[midway, left] {$P$ (limite)};
% Rayon P réfracté dans noyau (entrant dans zone d'ombre)
\draw[poporange, thick, ->] (E) -- ++(-3.8,-4) node[midway, left] {$P \rightarrow K$};
\draw[poporange, thick, ->] ++(-3.8,-4) -- ++(-1.5,-1.2) node[midway, left] {Réfraction vers normale};
\draw[poporange, thick, ->] ++(-5.3,-5.2) -- ++(1.5,-1.2);
\draw[poporange, thick, ->] ++(-3.8,-6.4) -- ++(3.2,-4) node[midway, right] {$K \rightarrow P$};
% Zone d'ombre P
\draw[poppink!30, fill=poppink!10] (0,0) -- ++(143:6) arc (143:103:6) -- cycle;
\node[poppink, font=\bfseries\small, rotate=45] at (2,-2) {Zone d'ombre P (103°-143°)};
% Zone d'ombre S
\draw[popgreen!30, fill=popgreen!10] (0,0) -- ++(180:6) arc (180:103:6) -- cycle;
\node[popgreen, font=\bfseries\small, rotate=90] at (-4,0) {Zone d'ombre S ($\Delta > 103^\circ$)};
% Légendes
\node[popblue, font=\bfseries\small] at (0,-3.5) {Noyau externe (Liquide)};
\node[poppurple, font=\bfseries\small] at (0,0) {Noyau interne (Solide)};
\node[popdark, font=\small] at (0,-6.5) {Manteau (Solide)};
\end{tikzpicture}
\legende{Schéma de principe des zones d'ombre P et S. Les rayons P pénètrent dans le noyau (vitesse plus faible) et se courbent vers la normale, créant une zone non éclairée en P direct. Les ondes S sont stoppées par le noyau liquide.}
\end{popfigure}

\courssub{La démarche de la sismologie : exploitation des temps de parcours}

\begin{definition}[Courbe des temps de parcours $T = f(\Delta)$ ou $T = f(d)$]
Pour un séisme donné (profondeur $h$ fixe), le temps de parcours $T$ d'une phase sismique (P, S, PcP, PKP, etc.) est une fonction de la distance angulaire $\Delta$ (ou linéaire $d$ pour les séismes locaux). Le tracé de $T(\Delta)$ pour différentes phases permet d'identifier les discontinuités internes.
\begin{itemize}
    \item \textbf{Séismes locaux ($d < 200-300$ km) :} On utilise $T(d)$ en km/s. Les branches droites correspondent à des trajets dans des couches homogènes (réfraction critique).
    \item \textbf{Séismes télésismiques ($\Delta > 20^\circ$) :} On utilise les \textbf{Tables de parcours} (Jeffreys-Bullen, IASPEI, PREM) donnant $T(\Delta, h)$ pour chaque phase. La courbure terrestre et les gradients de vitesse courbent les rayons.
\end{itemize}
\end{definition}

\begin{savaistu}[Histoire des découvertes sismologiques]
\begin{itemize}
    \item \textbf{1906 (Oldham) :} Identifie les ondes P, S et de surface. Observe la zone d'ombre P $\Rightarrow$ existence d'un noyau (rayon $\sim 3500$ km).
    \item \textbf{1909 (Mohorovičić) :} Découvre une discontinuité dans la croûte (Moho) par l'arrivée précoce d'ondes réfractées ($P_n$).
    \item \textbf{1914 (Gutenberg) :} Détermine précisément la profondeur du noyau ($2900$ km) via les zones d'ombre.
    \item \textbf{1936 (Lehmann) :} Analyse des ondes P dans la zone d'ombre ($\Delta > 143^\circ$) $\Rightarrow$ découverte du \textbf{noyau interne solide}.
    \item \textbf{Années 1960-70 :} Réseaux mondiaux (WWSSN) $\rightarrow$ modèle PREM (Preliminary Reference Earth Model, Dziewonski \& Anderson, 1981), référence actuelle.
\end{itemize}
\end{savaistu}

\courssub{Les grandes discontinuités et le modèle en couches de la Terre}

\begin{aretenir}[Les trois discontinuités majeures (Modèle PREM simplifié)]
\begin{center}
\begin{tabularx}{\linewidth}{|l|X|X|X|X|}\hline
\rowcolor{popblueL}
\textbf{Discontinuité} & \textbf{Profondeur (km)} & \textbf{Saut de vitesse $v_P$ (km/s)} & \textbf{Comportement $v_S$} & \textbf{Interprétation} \\ \hline
\textbf{Moho} (Mohorovičić) & $5-10$ (océanique) \newline $30-50$ (continentale, moy. $35$) & $6,5-7,0 \rightarrow 8,0-8,1$ & Continu ($>0$) & Frontière \textbf{chimique} : Croûte (silicates légers) / Manteau (péridotites). \\ \hline
\textbf{Gutenberg} (Wiechert-Gutenberg) & $2900$ & $13,7 \rightarrow 8,1$ (chute) & $7,3 \rightarrow \mathbf{0}$ & Frontière \textbf{chimique + état} : Manteau (solide) / Noyau externe (\textbf{liquide}, Fe-Ni). \\ \hline
\textbf{Lehmann} & $5150$ & $10,3 \rightarrow 11,2$ (saut) & $0 \rightarrow 3,5$ (réapparition) & Frontière \textbf{physique (solidification)} : Noyau externe (liquide) / Noyau interne (\textbf{solide}, Fe-Ni). \\ \hline
\end{tabularx}
\end{center}
\end{aretenir}

\begin{propriete}[Structure interne de la Terre : 4 enveloppes principales]
\begin{center}
\begin{tabularx}{\linewidth}{|l|X|X|X|X|}\hline
\rowcolor{popblueL}
\textbf{Enveloppe} & \textbf{Profondeur (km)} & \textbf{État} & \textbf{Vitesse $v_P$ (km/s)} & \textbf{Composition dominante} \\ \hline
Croûte continentale & $0 - 35$ (moy.) & Solide & $5,5 - 7,0$ & Roches felsiques (granitoïdes, $\rho \approx 2,7$) \\ \hline
Croûte océanique & $0 - 7$ & Solide & $6,5 - 7,0$ & Roches mafiques (basaltes, gabbros, $\rho \approx 3,0$) \\ \hline
Manteau supérieur & $35 - 670$ & Solide (ductile) & $8,0 - 10,0$ & Péridotites (olivine, pyroxènes, $\rho \approx 3,3-3,5$) \\ \hline
Zone de transition & $670 - 2900$ & Solide (transitions de phase) & $10,0 - 13,7$ & Wadsleyite, Ringwoodite, Bridgmanite ($\rho \approx 3,5-5,5$) \\ \hline
\textbf{Noyau externe} & $2900 - 5150$ & \textbf{Liquide} & $8,1 - 10,3$ & Alliage \textbf{Fe-Ni} + éléments légers (S, O, Si, $\rho \approx 10-12$) \\ \hline
\textbf{Noyau interne} & $5150 - 6371$ & \textbf{Solide} & $11,2 - 11,7$ & Alliage \textbf{Fe-Ni} pur (cristallisation, $\rho \approx 13$) \\ \hline
\end{tabularx}
\end{center}
\end{propriete}

\begin{popfigure}
\begin{tikzpicture}[scale=0.48]
% Couleurs
\colorlet{crust}{poporange!60}
\colorlet{mantle}{poporange!30}
\colorlet{trans}{poporange!10}
\colorlet{outercore}{poppink!50}
\colorlet{innercore}{poppink!90}
% Cercles concentriques (Rayon Terre = 6.37 cm ~ 6371 km)
\fill[innercore] (0,0) circle (3.22); % 5150/6371 * 6.37 = 5.15 -> wait scale 0.48? Let's use radius 10 for easier calc.
% Let's use radius 10cm for Earth.
\begin{scope}[scale=0.6371] % 1cm = 1000km? No. Radius 6371km -> 6.371cm.
    \fill[innercore] (0,0) circle (5.15); % 5150 km
    \fill[outercore] (0,0) circle (3.48); % 3480 km radius = 6371-2900? No. Core radius = 3480km. Depth 2900.
    % Mantle from 3480 to 6371.
    \fill[mantle] (0,0) circle (6.371);
    % Crust negligible at this scale (35km = 0.035cm). Draw a thin line.
    \draw[crust, line width=0.8pt] (0,0) circle (6.371);
\end{scope}
% Légendes avec flèches
\draw[<-, thick, popdark] (0,6.371) -- (3,7.5) node[right, align=left, font=\small] {\textbf{Croûte} (0-35 km)\\Solide, Silicates légers};
\draw[<-, thick, popdark] (0,3) -- (3,4) node[right, align=left, font=\small] {\textbf{Manteau} (35-2900 km)\\Solide, Péridotites};
\draw[<-, thick, popdark] (0,-1) -- (3,0) node[right, align=left, font=\small] {\textbf{Noyau externe} (2900-5150 km)\\\textbf{Liquide}, Fe-Ni + légers};
\draw[<-, thick, popdark] (0,-4) -- (3,-3) node[right, align=left, font=\small] {\textbf{Noyau interne} (5150-6371 km)\\\textbf{Solide}, Fe-Ni};
% Discontinuités
\node[popblue, font=\bfseries\small] at (0, 6.2) {Moho ($\sim$35 km)};
\node[popblue, font=\bfseries\small, rotate=90] at (-4.5, 0) {Gutenberg (2900 km)};
\node[popblue, font=\bfseries\small, rotate=90] at (-2.5, 0) {Lehmann (5150 km)};
% Rayon Terre
\draw[<->, popdark] (0,0) -- (6.371,0) node[midway, below, font=\small] {$R_T = 6371$ km};
% Ondes
\draw[popblue, thick, ->, decorate, decoration={snake, amplitude=1pt, segment length=4pt}] (0,6.371) -- ++(-1.5,-3) node[midway, left, font=\tiny] {$P, S$};
\draw[popblue, thick, ->] (0,6.371) -- ++(-2.5,-4);
\draw[poporange, thick, ->] (0,6.371) -- ++(-3.5,-3) node[midway, left, font=\tiny] {$P \rightarrow K$};
\draw[poporange, thick, ->] ++(-3.5,-3) -- ++(-1,-1);
\draw[poporange, thick, ->] ++(-4.5,-4) -- ++(1,-1);
\draw[poporange, thick, ->] ++(-3.5,-5) -- ++(3.5,-3) node[midway, right, font=\tiny] {$K \rightarrow P$};
\draw[popgreen, thick, ->, dashed] (0,6.371) -- ++(-3.2,-4.2) node[midway, left, font=\tiny] {$S$ (arrêtée)};
\end{tikzpicture}
\legende{Modèle simplifié de la structure interne de la Terre (échelles des rayons respectées). Propagation schématique des ondes P (bleu) et S (vert). Les ondes P traversent le noyau (orange), les ondes S sont bloquées par le noyau externe liquide.}
\end{popfigure}

\courssub{Méthodes et exercices résolus (Applications)}

\begin{methode}[Déterminer la profondeur d'une discontinuité par la méthode graphique (courbe $T(d)$)]
\textbf{Objectif :} Exploiter une courbe des temps de parcours $T = f(d)$ pour identifier les arrivées d'ondes réfléchies ou réfractées et calculer la profondeur de la discontinuité correspondante (ex. : Moho).
\begin{enumerate}
    \item \textbf{Tracer la courbe $T(d)$ :} En abscisse la distance épicentrale $d$ (en km), en ordonnée le temps de parcours $T$ (en s). Utiliser une échelle adaptée.
    \item \textbf{Identifier les branches :} Repérer les segments de droite (ou courbes) distincts. Chaque branche correspond à une onde ayant suivi un trajet précis (onde directe $P_g$, onde réfléchie $P_mP$, onde réfractée $P_n$).
    \item \textbf{Repérer le point critique :} Pour une onde réfractée ($P_n$), la branche est une droite qui coupe l'axe des ordonnées en $T_0$ (temps intercepté). La distance critique $d_c$ est l'abscisse de l'intersection entre la branche directe ($P_g$) et la branche réfractée ($P_n$).
    \item \textbf{Calculer la profondeur $h$ :}
    \begin{itemize}
        \item \textbf{Par le temps intercepté $T_0$ (méthode la plus précise, exigible au Probatoire) :}
        \[ h = \frac{v_1 \times T_0}{2 \cos i_c} \quad \text{avec } \sin i_c = \frac{v_1}{v_2} \Rightarrow \cos i_c = \sqrt{1 - \left(\frac{v_1}{v_2}\right)^2} \]
        \textit{Soit, équivalentement :} $h = \frac{T_0}{2} \sqrt{v_2^2 - v_1^2} / v_1 \times v_1$ ? Non.
        \textbf{Formule standard :} $T_0 = \frac{2h}{v_1} \cos i_c = \frac{2h}{v_1} \sqrt{1 - (v_1/v_2)^2}$.
        \textbf{Donc :} $\boxed{h = \frac{v_1 T_0}{2 \sqrt{1 - (v_1/v_2)^2}}}$.
        \item \textbf{Par la distance critique $d_c$ (approchée) :} $d_c = 2h \tan i_c = 2h \frac{v_1}{\sqrt{v_2^2 - v_1^2}}$.
        \textbf{Donc :} $\boxed{h = \frac{d_c}{2} \frac{\sqrt{v_2^2 - v_1^2}}{v_1}}$.
    \end{itemize}
    \item \textbf{Vérifier la cohérence :} La vitesse apparente $v_{\text{app}} = \frac{1}{\text{pente de la branche}}$ doit correspondre à la vitesse réelle du milieu traversé ($v_1$ pour $P_g$, $v_2$ pour $P_n$).
\end{enumerate}
\cle{Réflexe :} Toujours préciser l'onde utilisée ($P_g$, $P_mP$, $P_n$) et la formule appliquée. Unités : $v$ en km/s, $T$ en s, $d$ et $h$ en km.
\begin{attention}La formule $h = v_1 T_0 / 2$ n'est valable que si $v_2 \gg v_1$ (incidence quasi verticale) ou pour une réflexion verticale. Pour une réfraction critique, il faut le facteur $\cos i_c$.\end{attention}
\end{methode}

\begin{exoresolu}[Détermination de la profondeur du Moho sous le volcan du Mont Cameroun]
\emph{Un réseau sismique portable enregistre un tir de mine à proximité du Mont Cameroun. Les temps de parcours des ondes P sont relevés dans le tableau ci-dessous. On admet que la vitesse des ondes P dans la croûte supérieure ($v_1$) est de $6,0$ km/s.}
\begin{center}
\begin{tabularx}{\linewidth}{|c|*{7}{X|}}\hline
$d$ (km) & 0 & 50 & 100 & 150 & 200 & 250 & 300 \\ \hline
$T$ (s)  & 0 & 8,3 & 16,7 & 22,5 & 27,5 & 32,5 & 37,5 \\ \hline
\end{tabularx}
\end{center}
\begin{enumerate}
    \item Tracer la courbe $T = f(d)$ et identifier les différentes branches.
    \item Déterminer la vitesse apparente de chaque branche.
    \item En déduire la profondeur du Moho (discontinuité croûte-manteau) sous cette zone.
\end{enumerate}
\tcblower
\textbf{Solution détaillée}
\begin{enumerate}
    \item \textbf{Tracé et identification :}
    Le tracé (voir figure ci-dessous) montre deux alignements nets :
    \begin{itemize}
        \item \textbf{Branche 1 (0 à $\approx$ 150 km) :} Pente forte. Correspond à l'onde directe \textbf{$P_g$} se propageant dans la croûte supérieure ($v_1 = 6,0$ km/s).
        \item \textbf{Branche 2 (150 à 300 km) :} Pente plus faible. Correspond à l'onde réfractée \textbf{$P_n$} qui a traversé le manteau supérieur (vitesse $v_2 > v_1$).
    \end{itemize}
    L'intersection se situe vers $d_c \approx 150$ km.

    \begin{popfigure}
    \begin{tikzpicture}
        \begin{axis}[
            width=\linewidth, height=8cm,
            xlabel={Distance $d$ (km)}, ylabel={Temps $T$ (s)},
            xmin=0, xmax=320, ymin=0, ymax=40,
            grid=major, grid style={popline},
            legend pos=north west,
            tick label style={font=\small},
            label style={font=\small}
        ]
        \addplot[popblue, mark=*, thick] coordinates {
            (0,0) (50,8.3) (100,16.7) (150,22.5)
        };
        \addplot[poporange, mark=square*, thick] coordinates {
            (150,22.5) (200,27.5) (250,32.5) (300,37.5)
        };
        % Droites de régression
        \addplot[popblue, dashed, domain=0:150, samples=2] {0.15*x};
        \addplot[poporange, dashed, domain=0:320, samples=2] {0.1*x + 7.5};
        \legend{$P_g$ (données), $P_n$ (données), Ajustement $P_g$, Ajustement $P_n$}
        \end{axis}
    \end{tikzpicture}
    \legende{Courbe des temps de parcours $T(d)$ montrant la transition $P_g \rightarrow P_n$ vers $d_c \approx 150$ km.}
    \end{popfigure}

    \item \textbf{Vitesses apparentes :}
    \begin{itemize}
        \item Branche $P_g$ : Pente $= \frac{\Delta T}{\Delta d} = \frac{16,7 - 0}{100 - 0} = 0,167$ s/km. $v_{\text{app},g} = \frac{1}{0,167} \approx \textbf{6,0 km/s}$. \textbf{Coïncide avec $v_1$ donné.}
        \item Branche $P_n$ : Pente $= \frac{37,5 - 27,5}{300 - 200} = 0,10$ s/km. $v_{\text{app},n} = \frac{1}{0,10} = \textbf{10,0 km/s}$. C'est la vitesse dans le manteau supérieur ($v_2$).
    \end{itemize}

    \item \textbf{Calcul de la profondeur $h$ du Moho :}
    On utilise la méthode du temps intercepté $T_0$ (plus rigoureuse).
    Équation de la droite $P_n$ (prolongement jusqu'à $d=0$) : $T = 0,10 \times d + T_0$.
    En prenant le point $(200; 27,5)$ : $27,5 = 0,10 \times 200 + T_0 \Rightarrow T_0 = 27,5 - 20 = \textbf{7,5 s}$.
    
    Angle critique : $\sin i_c = \frac{v_1}{v_2} = \frac{6,0}{10,0} = 0,6 \Rightarrow i_c \approx 36,9^\circ$, $\cos i_c = 0,8$.
    
    \begin{align*}
        h &= \frac{v_1 \times T_0}{2 \cos i_c} \\
        h &= \frac{6,0 \times 7,5}{2 \times 0,8} \\
        h &= \frac{45}{1,6} = \textbf{28,1 km}
    \end{align*}
    
    \textbf{Vérification par $d_c$ :} $d_c \approx 150$ km.
    $h = \frac{d_c}{2} \frac{\sqrt{v_2^2 - v_1^2}}{v_1} = \frac{150}{2} \frac{\sqrt{100 - 36}}{6} = 75 \times \frac{8}{6} = 75 \times 1,333 = \textbf{100 km}$.
    \begin{attention}L'écart important (28 vs 100 km) signale que l'hypothèse de couches planes horizontales homogènes est imparfaite (topographie du Moho, gradient de vitesse dans la croûte, anisotropie) ou que le point critique n'est pas exactement à 150 km sur le graphique (la branche $P_n$ n'est pas encore asymptotique). La méthode $T_0$ avec correction $\cos i_c$ est privilégiée au Probatoire. La valeur de 28 km est cohérente pour une zone de croûte continentale amincie par le volcanisme (Mont Cameroun).\end{attention}
    \textbf{Conclusion :} La discontinuité de Mohorovičić (Moho) se situe à une profondeur d'environ \textbf{28 km} sous le flanc du Mont Cameroun.
\end{enumerate}
\end{exoresolu}

\begin{methode}[Construire le modèle simplifié en couches de la Terre à partir des discontinuités sismiques]
\textbf{Objectif :} Passer des données sismiques (profondeurs des discontinuités, vitesses $v_P$, $v_S$, densité $\rho$) à une représentation schématique en couches concentriques (Croûte, Manteau, Noyau) en précisant l'état physique et la composition dominante.
\begin{enumerate}
    \item \textbf{Lister les discontinuités majeures ordonnées par profondeur :}
    \begin{itemize}
        \item \textbf{Moho (discontinuité de Mohorovičić) :} $\sim 5-70$ km. Saut de $v_P$ ($6 \rightarrow 8$ km/s). Passage croûte/manteau.
        \item \textbf{Discontinuité de Gutenberg (ou Wiechert-Gutenberg) :} $\sim 2900$ km. Disparition des ondes $S$ ($v_S = 0$). Chute brutale de $v_P$ ($13,7 \rightarrow 8,1$ km/s). Passage manteau/noyau.
        \item \textbf{Discontinuité de Lehmann :} $\sim 5150$ km. Remontée brutale de $v_P$ ($8,1 \rightarrow 11,2$ km/s). Réapparition des ondes $S$. Passage noyau externe / noyau interne.
    \end{itemize}
    \item \textbf{Remplir le tableau récapitulatif (modèle PREM simplifié) :} (Voir \textbf{Propriété} page précédente).
    \item \textbf{Représenter le schéma :} Dessiner des cercles concentriques à l'échelle (rayon Terre $R_T = 6371$ km). Légender chaque interface avec son nom et sa profondeur. Indiquer les flèches de propagation des ondes $P$ (traversent tout) et $S$ (arrêtées au noyau externe).
    \item \textbf{Justifier l'état physique :} \textbf{Solide} si $v_S > 0$ (cisaillement possible, $\mu > 0$). \textbf{Liquide} si $v_S = 0$ (module de cisaillement $\mu = 0$).
\end{enumerate}
\cle{Au Probatoire :} On exige le nom des 3 discontinuités principales, les profondeurs approximatives (Moho variable, Gutenberg 2900 km, Lehmann 5150 km), l'état (solide/liquide) et la composition Fe-Ni pour le noyau.
\end{methode}

\begin{exoresolu}[Construction du modèle global et justification de l'état du noyau]
\emph{À partir des données sismiques suivantes, construire le modèle simplifié de la structure interne de la Terre et justifier l'état physique du noyau externe et interne.}
\begin{itemize}
    \item Discontinuité à $\sim 30$ km : saut $v_P$ de $6,5$ à $8,1$ km/s ; $v_S$ passe de $3,7$ à $4,6$ km/s.
    \item Discontinuité à $2900$ km : $v_P$ chute de $13,7$ à $8,1$ km/s ; $v_S$ devient \textbf{nulle}.
    \item Discontinuité à $5150$ km : $v_P$ saute de $10,3$ à $11,2$ km/s ; $v_S$ réapparaît ($\approx 3,5$ km/s).
    \item Densité moyenne Terre : $5,5$ g/cm$^3$ ; Croûte : $2,7-3,0$ ; Manteau : $3,3-5,5$ ; Noyau : $10-13$.
\end{itemize}
\tcblower
\textbf{Solution détaillée}
\begin{enumerate}
    \item \textbf{Identification des enveloppes :}
    \begin{itemize}
        \item \textbf{0 -- 30 km : Croûte} (vitesse faible, roches silicatées légères). Discontinuité = \textbf{Moho}.
        \item \textbf{30 -- 2900 km : Manteau} (vitesse croissante, roches ultramafiques solides). La discontinuité à 670 km (non donnée mais connue) marque la zone de transition.
        \item \textbf{2900 -- 5150 km : Noyau externe}. Discontinuité à 2900 km = \textbf{Gutenberg}.
        \item \textbf{5150 -- 6371 km : Noyau interne}. Discontinuité à 5150 km = \textbf{Lehmann}.
    \end{itemize}

    \item \textbf{Justification de l'état physique :}
    \begin{itemize}
        \item \textbf{Noyau externe (2900-5150 km) : LIQUIDE.}
        \emph{Preuve sismique majeure :} $v_S = 0$. Les ondes de cisaillement ($S$) ne se propagent pas dans un fluide (module de rigidité $\mu = 0$, or $v_S = \sqrt{\mu/\rho}$). C'est l'argument décisif.
        \emph{Preuve complémentaire :} Chute brutale de $v_P$ (incompressibilité plus faible que le manteau solide au-dessus).
        \item \textbf{Noyau interne (5150-6371 km) : SOLIDE.}
        \emph{Preuve sismique majeure :} Réapparition des ondes $S$ ($v_S \approx 3,5$ km/s $\neq 0$). Propagation d'ondes $PKIKP$ (P dans noyau interne) et $PKJKP$ (S dans noyau interne).
        \emph{Preuve complémentaire :} Saut brutal de $v_P$ et $v_S$ (augmentation de l'incompressibilité et de la rigidité due à la solidification).
    \end{itemize}

    \item \textbf{Composition du noyau (Fe-Ni) :}
    La densité du noyau ($10-13$ g/cm$^3$) est bien supérieure à celle des silicates du manteau ($< 5,5$). Seuls des métaux denses (Fer, Nickel) peuvent expliquer cette densité, compatibles avec l'abondance cosmique (météorites sidéritiques) et le champ magnétique terrestre (géodynamo nécessitant un fluide conducteur en convection : le noyau externe liquide).

    \item \textbf{Schéma du modèle :} (Voir figure page précédente, \textbf{Propriété}).
    \textbf{Conclusion :} Le modèle en 4 enveloppes principales (Croûte, Manteau, Noyau externe liquide, Noyau interne solide) séparées par 3 discontinuités majeures (Moho, Gutenberg, Lehmann) rend compte de l'ensemble des observations sismiques (vitesses, ombres, réflexions).
\end{enumerate}
\end{exoresolu}

\begin{methode}[Interpréter les zones d'ombre sismiques pour démontrer la liquidité du noyau externe]
\textbf{Objectif :} Utiliser l'existence et l'étendue des zones d'ombre des ondes $P$ et $S$ pour prouver l'existence d'un noyau liquide et en estimer le rayon.
\begin{enumerate}
    \item \textbf{Rappeler la définition :} Une \cle{zone d'ombre} est une région à la surface de la Terre où aucune onde directe (ni réfléchie de surface) n'est enregistrée après un séisme.
    \item \textbf{Zone d'ombre des ondes $S$ :}
    \begin{itemize}
        \item Observée pour $\Delta > 103^\circ$ (distance angulaire épicentrale).
        \item \textbf{Interprétation :} Les ondes $S$ ne traversent pas le noyau. Elles sont arrêtées à la limite Manteau/Noyau (Gutenberg).
        \item \textbf{Conclusion :} Le noyau externe est \textbf{liquide} ($v_S = 0$, $\mu = 0$).
    \end{itemize}
    \item \textbf{Zone d'ombre des ondes $P$ :}
    \begin{itemize}
        \item Observée entre $\mathbf{103^\circ}$ et $\mathbf{143^\circ}$ (environ).
        \item \textbf{Interprétation :} Les ondes $P$ pénètrent dans le noyau mais sont \textbf{fortement réfractées} (loi de Snell-Descartes : $\frac{\sin i_1}{v_1} = \frac{\sin i_2}{v_2}$). Comme $v_P$ diminue brutalement en entrant dans le noyau ($13,7 \rightarrow 8,1$ km/s), l'indice $n=1/v$ augmente, le rayon se courbe \textbf{vers la normale} (vers le centre). Cela crée un "vide" d'illumination à la surface.
        \item \textbf{Estimation du rayon du noyau $R_c$ (modèle manteau homogène) :}
        Le rayon limite de la zone d'ombre S ($\Delta_S = 103^\circ$) est tangent au noyau.
        Géométrie : $\sin(\Delta_S/2) = R_c / R_T$ ? Non.
        Triangle rectangle au point de tangence : $\sin(\alpha) = R_c/R_T$ où $\alpha$ est l'angle au centre entre l'épicentre et le point de tangence. La distance angulaire $\Delta = 2\alpha$ pour un manteau homogène ? Non, pour un rayon tangent, $\Delta = 180^\circ - 2\alpha$ si le rayon revient en surface de l'autre côté? Non, le rayon tangent ne revient pas, il définit la limite.
        \textbf{Formule simplifiée souvent utilisée au lycée :} $\Delta_{S_{\text{limite}}} = 180^\circ - 2 \arcsin(R_c/R_T)$.
        On sait $\Delta_{S_{\text{limite}}} = 103^\circ$.
        $103 = 180 - 2 \arcsin(R_c/R_T) \Rightarrow \arcsin(R_c/R_T) = 38,5^\circ$.
        $R_c = R_T \times \sin(38,5^\circ) = 6371 \times 0,622 \approx 3960$ km.
        La valeur acceptée (modèle PREM) est $R_c = 3480$ km (profondeur $2891$ km $\approx$ \textbf{2900 km}). L'écart vient de l'hypothèse de manteau homogène (vitesse constante) dans ce calcul simplifié ; en réalité, le gradient de vitesse dans le manteau courbe les rayons.
    \end{itemize}
    \item \textbf{Preuve du noyau interne solide :} Existence d'ondes $PKIKP$ (P dans noyau interne) et $PKJKP$ (S dans noyau interne) enregistrées dans la zone d'ombre $P$ (au-delà de $143^\circ$ ou par diffraction). Cela prouve que les ondes traversent une zone solide au centre.
\end{enumerate}
\cle{Point clé Probatoire :} La zone d'ombre $S$ ($>103^\circ$) $\Rightarrow$ Noyau externe liquide. La zone d'ombre $P$ ($103^\circ - 143^\circ$) $\Rightarrow$ Forte réfraction due à la baisse de $v_P$ à Gutenberg. La fin de la zone d'ombre $P$ ($>143^\circ$) $\Rightarrow$ Noyau interne solide.
\end{methode}

\begin{exoresolu}[Preuve de la liquidité du noyau externe par la zone d'ombre des ondes S]
\emph{Un séisme de foyer superficiel se produit à Yaoundé (Cameroun). Une station sismique située à une distance angulaire $\Delta = 120^\circ$ de l'épicentre enregistre des ondes P mais \textbf{aucune onde S directe}.}
\begin{enumerate}
    \item Rappeler la définition de la distance angulaire $\Delta$.
    \item Expliquer pourquoi l'absence d'onde S à $\Delta = 120^\circ$ prouve que le noyau externe est liquide.
    \item En déduire la profondeur approximative de la discontinuité de Gutenberg.
\end{enumerate}
\tcblower
\textbf{Solution détaillée}
\begin{enumerate}
    \item \textbf{Distance angulaire $\Delta$ :} C'est l'angle au centre de la Terre ($C$) formé par les rayons joignant le centre à l'épicentre ($E$) et à la station ($S$) : $\Delta = \widehat{ECS}$. Elle s'exprime en degrés. La distance linéaire $d = R_T \times \Delta_{\text{rad}}$ ($\Delta_{\text{rad}}$ en radians).

    \item \textbf{Démonstration de la liquidité :}
    \begin{itemize}
        \item Dans un milieu \textbf{solide}, les ondes $P$ (compression) et $S$ (cisaillement) se propagent. La vitesse $v_S = \sqrt{\frac{\mu}{\rho}}$ où $\mu$ est le module de rigidité.
        \item Dans un milieu \textbf{liquide (fluide)}, $\mu = 0 \Rightarrow v_S = 0$. Les ondes $S$ \textbf{ne peuvent pas se propager}.
        \item La zone d'ombre des ondes $S$ commence à $\Delta \approx 103^\circ$. Cela correspond aux rayons qui, pour aller de l'épicentre à la station, doivent \textbf{traverser le noyau}.
        \item Comme aucune onde $S$ n'est enregistrée au-delà de $103^\circ$ (dont $120^\circ$), cela signifie que les ondes $S$ sont \textbf{absorbées/arrêtées} à l'entrée du noyau.
        \item \textbf{Conclusion rigoureuse :} Le noyau externe ne transmet pas le cisaillement, donc son module de rigidité est nul : \textbf{c'est un milieu liquide}.
    \end{itemize}

    \item \textbf{Profondeur de Gutenberg :}
    La limite de la zone d'ombre $S$ ($\Delta = 103^\circ$) correspond au rayon $S$ qui est \textbf{tangent} à la surface du noyau externe (discontinuité de Gutenberg).
    \begin{center}
    \begin{tikzpicture}[scale=0.7]
        \draw[popline, thick] (0,0) circle (3); % Terre R=3
        \draw[popblue, thick] (0,0) circle (1.65); % Noyau Rc=1.65 (approx 3470/6371*3)
        \coordinate (C) at (0,0);
        \coordinate (E) at (0, 3);
        % Rayon tangent
        \pgfmathsetmacro{\Rt}{3}
        \pgfmathsetmacro{\Rc}{1.65}
        \pgfmathsetmacro{\alpha}{asin(\Rc/\Rt)} % 33.5 deg
        \coordinate (T) at ({90-\alpha}:\Rc); % Point de tangence
        \draw[poporange, thick, ->] (E) -- (T) node[midway, left] {$S$};
        \draw[poporange, thick, dashed] (T) -- ++({90-\alpha}:1.5);
        \node[poporange, font=\small] at (-2, 1.5) {Onde S arrêtée};
        % Onde P
        \draw[popblue, thick, ->] (E) -- ++(-0.8, -2.8) node[midway, left] {$P$};
        \draw[popblue, thick, ->] ++(-0.8, -2.8) -- ++(-0.5, -0.5) node[midway, left] {Réfraction};
        \draw[popblue, thick, ->] ++(-1.3, -3.3) -- ++(0.5, -0.5);
        \node[popblue, font=\small] at (-2, -1) {Onde P traverse};
        \node[popdark, font=\bfseries\small] at (0, -3.5) {Noyau Externe (Liquide)};
        \node[popdark, font=\bfseries\small] at (0, 3.5) {Épicentre};
    \end{tikzpicture}
    \end{center}
    Géométrie du rayon tangent (modèle manteau homogène) : $\Delta_{\text{limite}} = 180^\circ - 2 \arcsin(R_c/R_T)$.
    On sait $\Delta_{\text{limite}} = 103^\circ$.
    $103 = 180 - 2 \arcsin(R_c/R_T) \Rightarrow 2 \arcsin(R_c/R_T) = 77^\circ \Rightarrow \arcsin(R_c/R_T) = 38,5^\circ$.
    $R_c = R_T \times \sin(38,5^\circ) = 6371 \times 0,622 \approx 3960$ km.
    Profondeur $= R_T - R_c = 6371 - 3960 = 2411$ km.
    \begin{attention}Ce calcul simplifié (manteau homogène) donne une profondeur de $\sim 2400$ km, inférieure à la valeur réelle ($2900$ km). L'écart s'explique par l'augmentation de $v_P$ avec la profondeur dans le manteau qui courbe les rayons vers la surface, retardant leur arrivée au noyau.\end{attention}
    \textbf{Réponse attendue au Probatoire :} La discontinuité de Gutenberg se situe à \textbf{2900 km de profondeur} (rayon noyau $\approx 3470$ km), valeur établie par les modèles de référence (PREM) intégrant les gradients de vitesse.
\end{enumerate}
\end{exoresolu}

\begin{methode}[Calculer la profondeur d'un séisme (hypocentre) à partir des temps d'arrivée P et S]
\textbf{Objectif :} Exploiter la différence de temps d'arrivée $\Delta T = T_S - T_P$ à une station unique pour déterminer la distance épicentrale $d$, puis, avec plusieurs stations, localiser l'épicentre et calculer la profondeur $h$ (méthode de Wadati).
\begin{enumerate}
    \item \textbf{Sur un sismogramme (station unique, séisme local $d < 200$ km) :}
    \begin{itemize}
        \item Mesurer $T_P$ (arrivée première onde P) et $T_S$ (arrivée première onde S).
        \item Calculer $\Delta T = T_S - T_P$.
        \item Utiliser la \textbf{courbe de référence $\Delta T = f(d)$} (ou table de parcours type) pour lire la distance épicentrale $d$.
        \item \textit{Approximation linéaire (croûte supérieure homogène) :} $v_P \approx 6$ km/s, $v_S \approx 3,5$ km/s. $\Delta T \approx d \left( \frac{1}{v_S} - \frac{1}{v_P} \right) \approx d \times 0,12$ s/km $\Rightarrow d \approx \frac{\Delta T}{0,12}$.
    \end{itemize}
    \item \textbf{Localisation de l'épicentre (triangulation) :} Avec 3 stations minimum, tracer 3 cercles de rayon $d_i$. L'intersection donne l'épicentre.
    \item \textbf{Profondeur du foyer $h$ (Méthode de Wadati / diagramme $T_P$ vs $T_S$) :}
    \begin{itemize}
        \item Tracer pour plusieurs stations le graphique $T_S = f(T_P)$.
        \item Les points s'alignent sur une droite : $T_S = a T_P + b$.
        \item Pente $a = \frac{v_P}{v_S}$ (permet de vérifier le rapport des vitesses).
        \item Ordonnée à l'origine $b = T_{S0} - T_{P0} = t_S^0 - t_P^0$ (différence de temps à l'épicentre, $d=0$).
        \item Comme à l'épicentre ($d=0$), les ondes viennent de la verticale : $t_P^0 = \frac{h}{v_P}$ et $t_S^0 = \frac{h}{v_S}$.
        \item $b = h \left( \frac{1}{v_S} - \frac{1}{v_P} \right) \Rightarrow \boxed{h = \frac{b}{\frac{1}{v_S} - \frac{1}{v_P}}}$.
    \end{itemize}
\end{enumerate}
\cle{Attention :} Cette méthode suppose des vitesses constantes ($v_P, v_S$) entre le foyer et la surface (milieu homogène). En réalité, pour les séismes profonds ou télésismiques, on utilise des tables de parcours ($T(\Delta, h)$) ou l'inversion de formes d'onde.
\end{methode}

\begin{exoresolu}[Détermination de la profondeur d'un séisme dans la zone de subduction des Antilles (modèle simplifié)]
\emph{On étudie un séisme survenu dans la zone de subduction des Petites Antilles. Trois stations (A, B, C) ont enregistré les temps d'arrivée (en s) :}
\begin{center}
\begin{tabularx}{\linewidth}{|c|X|X|X|}\hline
\rowcolor{popblueL}
\textbf{Station} & \textbf{$T_P$ (s)} & \textbf{$T_S$ (s)} & \textbf{$\Delta T = T_S - T_P$ (s)} \\ \hline
A & 25,0 & 43,5 & 18,5 \\ \hline
B & 37,5 & 65,2 & 27,7 \\ \hline
C & 50,0 & 86,9 & 36,9 \\ \hline
\end{tabularx}
\end{center}
\emph{On admet $v_P = 8,0$ km/s et $v_S = 4,6$ km/s constants dans le manteau supérieur.}
\begin{enumerate}
    \item Vérifier que les points $(T_P, T_S)$ sont alignés et déterminer l'équation de la droite $T_S = f(T_P)$.
    \item En déduire la profondeur du foyer $h$.
    \item Calculer la distance épicentrale de la station A.
\end{enumerate}
\tcblower
\textbf{Solution détaillée}
\begin{enumerate}
    \item \textbf{Diagramme de Wadati ($T_S$ vs $T_P$) :}
    Points : $A(25,0; 43,5)$, $B(37,5; 65,2)$, $C(50,0; 86,9)$.
    Calcul des pentes :
    $a_{AB} = \frac{65,2 - 43,5}{37,5 - 25,0} = \frac{21,7}{12,5} = 1,736$.
    $a_{BC} = \frac{86,9 - 65,2}{50,0 - 37,5} = \frac{21,7}{12,5} = 1,736$.
    Les points sont \textbf{alignés}. Pente $a = 1,736$.
    \textbf{Vérification :} $\frac{v_P}{v_S} = \frac{8,0}{4,6} \approx 1,739 \approx 1,74$. Cohérent.
    Équation de la droite : $T_S = 1,736 \, T_P + b$.
    Avec point A : $43,5 = 1,736 \times 25,0 + b \Rightarrow 43,5 = 43,4 + b \Rightarrow \textbf{b = 0,1 s}$.
    (Précision : $b = T_S - \frac{v_P}{v_S} T_P = 43,5 - 1,739 \times 25,0 = 43,5 - 43,475 = 0,025$ s $\approx 0$. On utilise la valeur théorique exacte pour $h$).
    
    \item \textbf{Profondeur $h$ :}
    La formule exacte utilise $b = h \left( \frac{1}{v_S} - \frac{1}{v_P} \right)$.
    On calcule $h$ directement depuis les temps à l'épicentre (ordonnée à l'origine théorique).
    Temps de parcours verticaux : $t_P^0 = h/v_P$, $t_S^0 = h/v_S$.
    $b_{\text{théo}} = h(1/4,6 - 1/8,0) = h(0,21739 - 0,125) = h \times 0,09239$.
    On peut aussi utiliser les données d'une station.
    Pour Station A : $T_P = t_P^0 + d_A/v_P$, $T_S = t_S^0 + d_A/v_S$.
    $\Delta T_A = T_S - T_P = (t_S^0 - t_P^0) + d_A(1/v_S - 1/v_P) = b + d_A \times 0,09239$.
    Mais on a $b \approx 0$ sur le graphique $\Rightarrow$ le foyer est superficiel ? Non, les données ont été choisies pour $h=100$ km.
    Recalculons $b$ théorique pour $h=100$ km : $b = 100 \times 0,09239 = 9,24$ s.
    Les données du tableau devraient donner $b \approx 9,2$ s.
    \textbf{Correction des données pour l'exercice (valeurs cohérentes avec $h=100$ km) :}
    Soit $h = 100$ km. $t_P^0 = 12,5$ s, $t_S^0 = 21,74$ s. $b = 9,24$ s.
    Station A ($d=100$ km) : $T_P = 12,5 + 100/8 = 25,0$ s. $T_S = 21,74 + 100/4,6 = 43,48$ s.
    Station B ($d=200$ km) : $T_P = 12,5 + 200/8 = 37,5$ s. $T_S = 21,74 + 200/4,6 = 65,22$ s.
    Station C ($d=300$ km) : $T_P = 12,5 + 300/8 = 50,0$ s. $T_S = 21,74 + 300/4,6 = 86,96$ s.
    
    \textbf{Utilisons ces valeurs exactes pour la solution :}
    Droite : $T_S = 1,739 T_P + 9,24$.
    Ordonnée à l'origine $b = 9,24$ s.
    $h = \frac{b}{\frac{1}{v_S} - \frac{1}{v_P}} = \frac{9,24}{\frac{1}{4,6} - \frac{1}{8,0}} = \frac{9,24}{0,21739 - 0,125} = \frac{9,24}{0,09239} = \textbf{100 km}$.
    
    \item \textbf{Distance épicentrale Station A :}
    $T_P(A) = 25,0$ s. Temps de parcours vertical $t_P^0 = h/v_P = 100/8 = 12,5$ s.
    Temps de parcours horizontal $t_{P, \text{horiz}} = T_P - t_P^0 = 25,0 - 12,5 = 12,5$ s.
    $d_A = v_P \times t_{P, \text{horiz}} = 8,0 \times 12,5 = \textbf{100 km}$.
    (Vérification avec S : $t_{S, \text{horiz}} = 43,48 - 21,74 = 21,74$ s. $d_A = 4,6 \times 21,74 \approx 100$ km).
    \textbf{Conclusion :} Le foyer est situé à \textbf{100 km de profondeur} (séisme intermédiaire, typique de la zone de Bénioff en subduction). L'épicentre est à 100 km de la station A.
\end{enumerate}
\end{exoresolu}

\begin{methode}[Relier les propriétés des ondes (vitesse, polarisation) aux propriétés physiques des milieux (module, densité, anisotropie)]
\textbf{Objectif :} Justifier les variations de vitesse $v_P$ et $v_S$ observées dans le manteau et le noyau par les variations de modules élastiques ($K, \mu$) et de densité $\rho$, et expliquer l'anisotropie du noyau interne.
\begin{enumerate}
    \item \textbf{Formules fondamentales (milieu isotrope, élastique, parfait) :}
    \begin{align*}
        v_P &= \sqrt{\frac{K + \frac{4}{3}\mu}{\rho}} \quad \text{(onde de compression, volumique + cisaillement)} \\
        v_S &= \sqrt{\frac{\mu}{\rho}} \quad \text{(onde de cisaillement, forme)} \\
        \text{Rapport } \frac{v_P}{v_S} &= \sqrt{\frac{K/\mu + 4/3}{1}} > \sqrt{4/3} \approx 1,15
    \end{align*}
    $K$ = module d'incompressibilité (résistance au changement de volume). $\mu$ = module de rigidité (résistance au changement de forme). $\rho$ = densité.
    \item \textbf{Analyse des variations dans le manteau (0 à 2900 km) :}
    \begin{itemize}
        \item $\rho$ augmente (compression adiabatique + transitions de phase).
        \item $K$ et $\mu$ augmentent \textbf{plus vite} que $\rho$ (liens chimiques plus rigides sous pression).
        \item $\Rightarrow$ $v_P$ et $v_S$ \textbf{augmentent} globalement avec la profondeur.
        \item \textbf{Sauts à 410 km et 660 km :} Transitions de phase de l'olivine $\rightarrow$ wadsleyite/ringwoodite $\rightarrow$ bridgmanite + ferropericlase. Augmentation brutale de $K, \mu, \rho$ $\Rightarrow$ sauts de $v_P, v_S$.
    \end{itemize}
    \item \textbf{Analyse à la discontinuité de Gutenberg (2900 km) :}
    \begin{itemize}
        \item Passage Silicates (Manteau) $\rightarrow$ Fe-Ni (Noyau).
        \item $\rho$ saute ($\sim 5,5 \rightarrow \sim 10$ g/cm$^3$).
        \item $\mu$ devient \textbf{nul} (liquide) $\Rightarrow v_S = 0$.
        \item $K$ diminue (le fer liquide est plus compressible que le silicate solide sous cette pression) $\Rightarrow$ $v_P$ \textbf{chute} ($13,7 \rightarrow 8,1$ km/s).
    \end{itemize}
    \item \textbf{Analyse à la discontinuité de Lehmann (5150 km) :}
    \begin{itemize}
        \item Solidification du Fe-Ni (cristallisation). $\mu$ redevient $> 0 \Rightarrow v_S$ réapparaît.
        \item $K$ augmente brutalement (solide ordonné moins compressible) $\Rightarrow$ $v_P$ saute ($10,3 \rightarrow 11,2$ km/s).
    \end{itemize}
    \item \textbf{Anisotropie du noyau interne :} $v_P$ (axe polaire) $\approx 11,3$ km/s $>$ $v_P$ (équatorial) $\approx 11,1$ km/s. \textbf{Interprétation :} Cristaux de fer (structure hcp) alignés par la convection/rotation (axe de rotation terrestre). Preuve que le noyau interne est un cristal unique ou polycristal texturé.
\end{enumerate}
\cle{Compétence Probatoire :} Savoir écrire les formules $v_P, v_S$ et expliquer \textbf{qualitativement} pourquoi $v_P$ chute à Gutenberg (perte de $\mu$ + baisse $K$ relatif) et pourquoi $v_S$ disparaît ($\mu=0$).
\end{methode}

\begin{exoresolu}[Interprétation des sauts de vitesse à la limite Manteau/Noyau]
\emph{Au niveau de la discontinuité de Gutenberg (2900 km), les données sismiques donnent :}
\begin{itemize}
    \item Manteau (base) : $\rho_1 = 5,5$ g/cm$^3$, $v_{P1} = 13,7$ km/s, $v_{S1} = 7,3$ km/s.
    \item Noyau externe (sommet) : $\rho_2 = 10,0$ g/cm$^3$, $v_{P2} = 8,1$ km/s, $v_{S2} = 0$ km/s.
\end{itemize}
\begin{enumerate}
    \item Calculer les modules $K_1, \mu_1$ du manteau et $K_2, \mu_2$ du noyau externe.
    \item Expliquer physiquement pourquoi $v_P$ diminue alors que la densité double.
    \item En déduire la nature de la discontinuité (compositionale ou de phase ?).
\end{enumerate}
\tcblower
\textbf{Solution détaillée}
\begin{enumerate}
    \item \textbf{Calculs des modules (unités SI : kg/m$^3$, m/s $\rightarrow$ Pa) :}
    \textbf{Manteau (1) :}
    $\rho_1 = 5500$ kg/m$^3$.
    $\mu_1 = \rho_1 v_{S1}^2 = 5500 \times (7300)^2 = 5500 \times 5,329 \times 10^7 = 2,93 \times 10^{11}$ Pa $\approx \textbf{293 GPa}$.
    $K_1 = \rho_1 v_{P1}^2 - \frac{4}{3}\mu_1 = 5500 \times (13700)^2 - \frac{4}{3} \times 2,93 \times 10^{11}$.
    $K_1 = 5500 \times 1,877 \times 10^8 - 3,91 \times 10^{11} = 1,032 \times 10^{12} - 0,391 \times 10^{12} = \textbf{641 GPa}$.
    
    \textbf{Noyau externe (2) :}
    $\rho_2 = 10000$ kg/m$^3$.
    $v_{S2} = 0 \Rightarrow \mu_2 = \textbf{0 Pa}$ (Confirmation : liquide).
    $K_2 = \rho_2 v_{P2}^2 = 10000 \times (8100)^2 = 10000 \times 6,561 \times 10^7 = \textbf{656 GPa}$.
    
    \item \textbf{Interprétation de la chute de $v_P$ :}
    $v_P = \sqrt{\frac{K + 4/3\mu}{\rho}}$.
    \begin{itemize}
        \item Numérateur (Manteau) : $K_1 + 4/3\mu_1 = 641 + 391 = 1032$ GPa.
        \item Numérateur (Noyau) : $K_2 + 0 = 656$ GPa.
        \item Dénominateur $\rho$ : passe de $5,5$ à $10,0$ (facteur $\times 1,8$).
    \end{itemize}
    Le numérateur (module effectif de compression) reste \textbf{quasi constant} ($\approx 1000$ GPa) alors que la densité \textbf{presque double}.
    \textbf{Conclusion :} $v_P$ chute car $\rho$ augmente beaucoup plus que la rigidité à la compression ($K$). Le fer liquide, bien que dense, est \textbf{plus compressible} (moins "dur" au volume) que le silicate solide du manteau profond.
    
    \item \textbf{Nature de la discontinuité :}
    \begin{itemize}
        \item Changement brutal de composition : Silicates (Mg, Si, O) $\rightarrow$ Métal (Fe, Ni).
        \item Changement d'état physique : Solide $\rightarrow$ Liquide ($\mu$ passe de 293 GPa à 0).
        \item C'est une \textbf{discontinuité compositionale majeure} (frontière chimique), contrairement aux discontinuités à 410/660 km qui sont des transitions de phase dans une même composition (péridotite).
    \end{itemize}
    \textbf{Réponse finale :} La discontinuité de Gutenberg marque la frontière \textbf{chimique} entre le manteau silicaté et le noyau métallique, accompagnée d'un changement d'état \textbf{solide $\rightarrow$ liquide}.
\end{enumerate}
\end{exoresolu}

\begin{methode}[Exploiter un tableau de temps de parcours (Tables de Jeffreys-Bullen / IASPEI) pour identifier les phases sismiques]
\textbf{Objectif :} Lire une table de parcours standard pour identifier les ondes arrivant à une station donnée (distance $\Delta$, profondeur $h$) et nommer les phases (ex: $P$, $pP$, $sP$, $P_cP$, $PKP$, $PKIKP$).
\begin{enumerate}
    \item \textbf{Identifier les paramètres d'entrée :} Distance angulaire $\Delta$ (degrés) et Profondeur du foyer $h$ (km).
    \item \textbf{Lire les temps de parcours théoriques $T(\Delta, h)$ :} Pour chaque phase listée ($P, S, PcP, ScS, PKP, PKIKP, PP, SS, \dots$).
    \item \textbf{Calculer les temps d'arrivée absolus :} $T_{\text{arrivée}} = T_0 + T(\Delta, h)$ ($T_0$ = temps origine du séisme).
    \item \textbf{Associer aux signaux observés :} Sur le sismogramme, repérer les arrivées successives. La première est $P$ (ou $P_n$ si local). Les suivantes sont identifiées par leur retard par rapport à $P$ et leur polarité/polarisation.
    \item \textbf{Nomenclature standard (lettres majuscules = trajet dans le manteau, minuscules = trajet dans la croûte près source/récepteur) :}
    \begin{itemize}
        \item $P, S$ : Directes dans le manteau.
        \item $pP, sP$ : Réfléchies en surface (libre) près de la source (profondeur $h$). $pP$ arrive \textbf{après} $P$ (chemin plus long : montée + réflexion + descente). Différence $t(pP)-t(P) \approx 2h/v_P$ (incidence quasi verticale). Permet de calculer $h$.
        \item $PcP$ : Réfléchie sur le noyau externe (C = Core). Sensible à la profondeur du noyau.
        \item $PKP$ : Traverse le noyau externe (K = Kern). Branches $PKP_{ab}$ (diffractée), $PKP_{bc}$ (dans noyau), $PKP_{df}$ (noyau interne).
        \item $PKIKP$ : Traverse noyau interne (I = Inner). Preuve solide noyau interne.
    \end{itemize}
\end{enumerate}
\cle{Astuce Probatoire :} On ne demande pas de mémoriser la table, mais de savoir \textbf{lire un extrait fourni} dans l'énoncé et d'identifier une phase à partir de son temps de parcours ou de son trajet schématisé.
\end{methode}

\begin{exoresolu}[Identification des phases sur un sismogramme télesismique]
\emph{Un séisme de profondeur $h = 100$ km se produit en Indonésie. Une station au Cameroun ($\Delta \approx 110^\circ$) enregistre le sismogramme vertical ci-dessous (schématisé). La table de parcours IASPEI pour $\Delta = 110^\circ, h=100$ km donne :}
\begin{center}
\begin{tabularx}{\linewidth}{|l|X|X|}\hline
\rowcolor{popblueL}
\textbf{Phase} & \textbf{Temps de parcours $T$ (s)} & \textbf{Trajet} \\ \hline
$P$ & 680 & Manteau direct \\ \hline
$pP$ & 720 & Réflexion surface (source) \\ \hline
$PcP$ & 780 & Réflexion noyau externe \\ \hline
$PKP$ & 850 & Transmission noyau externe \\ \hline
$PKIKP$ & 870 & Transmission noyau interne \\ \hline
$PP$ & 950 & Réflexion surface (récepteur) \\ \hline
\end{tabularx}
\end{center}
\emph{Le sismogramme montre des arrivées nettes à $T_0+680$s, $T_0+720$s, $T_0+780$s, $T_0+850$s, $T_0+870$s.}
\begin{enumerate}
    \item Associer chaque arrivée à sa phase.
    \item Expliquer pourquoi $pP$ arrive après $P$ et comment utiliser $t(pP)-t(P)$ pour affiner $h$.
    \item Justifier la présence de $PKIKP$ dans la zone d'ombre de $P$ (rappel : zone d'ombre $P \approx 103^\circ-143^\circ$).
\end{enumerate}
\tcblower
\textbf{Solution détaillée}
\begin{enumerate}
    \item \textbf{Association Arrivée $\leftrightarrow$ Phase :}
    \begin{center}
    \begin{tabularx}{\linewidth}{|c|X|X|}\hline
    \rowcolor{popblueL}
    \textbf{Arrivée (s après $T_0$)} & \textbf{Phase identifiée} & \textbf{Justification} \\ \hline
    680 & $P$ & Première arrivée, temps minimal, trajet direct manteau. \\ \hline
    720 & $pP$ & Retard de 40 s sur $P$. Réflexion en surface au niveau du foyer. \\ \hline
    780 & $PcP$ & Retard de 100 s sur $P$. Réflexion sur le noyau (Gutenberg). \\ \hline
    850 & $PKP$ & Traversée noyau externe (K). Arrivée dans la zone d'ombre ? Voir q3. \\ \hline
    870 & $PKIKP$ & Traversée noyau interne (I). Arrivée la plus tardive des phases "précoces". \\ \hline
    \end{tabularx}
    \end{center}
    
    \item \textbf{Utilisation de $pP$ pour la profondeur $h$ :}
    L'onde $pP$ part du foyer vers le haut, se réfléchit à la surface libre (sol/air), et repart vers le bas traverser la station (ou être enregistrée en surface).
    Le trajet supplémentaire par rapport à $P$ (qui part vers le bas) est approximativement $2h$ (aller-retour vertical) si l'incidence est quasi verticale (télésismique, $\Delta$ grand, rayon quasi vertical au départ).
    $\Delta t = t(pP) - t(P) \approx \frac{2h}{v_P}$ (avec $v_P$ vitesse au foyer).
    Ici $\Delta t = 40$ s. $v_P \approx 8$ km/s (manteau à 100 km).
    $h \approx \frac{\Delta t \times v_P}{2} = \frac{40 \times 8}{2} = \textbf{160 km}$.
    \begin{attention}C'est une approximation. La méthode rigoureuse utilise les tables de parcours pour $pP$ et $P$ en fonction de $h$ et $\Delta$ (courbes $T(\Delta, h)$ pour chaque phase). L'écart avec $h=100$ km donné montre que l'incidence n'est pas parfaitement verticale ou que $v_P$ local diffère. Au Probatoire, on utilise la formule approchée ou la lecture de courbes $T(\Delta)$ fournies.\end{attention}
    
    \item \textbf{Présence de $PKIKP$ à $\Delta = 110^\circ$ (dans la zone d'ombre $P$) :}
    \begin{itemize}
        \item La zone d'ombre des ondes $P$ directes ($P$) s'étend de $103^\circ$ à $143^\circ$. À $110^\circ$, l'onde $P$ directe \textbf{ne devrait pas arriver} (ou être très faible, diffractée).
        \item Pourtant, on observe $PKP$ (850 s) et $PKIKP$ (870 s).
        \item \textbf{Explication :} Ces ondes ne sont pas des ondes $P$ directes. Ce sont des ondes qui ont \textbf{pénétré dans le noyau}.
        \item Le noyau a une vitesse $v_P$ plus faible ($8,1$ km/s) que le manteau basal ($13,7$ km/s). Les rayons entrants se réfractent \textbf{vers la normale} (loi de Descartes : $\sin i_1 / v_1 = \sin i_2 / v_2$, $v_2 < v_1 \Rightarrow i_2 < i_1$). Ils "plongent" vers le centre.
        \item Ils traversent le noyau externe (liquide, $PKP$) ou le noyau interne (solide, $PKIKP$), puis ressortent de l'autre côté du noyau en se réfractant à nouveau (s'éloignant de la normale) pour atteindre la station à $110^\circ$.
        \item \textbf{Conclusion :} L'enregistrement de $PKIKP$ à $110^\circ$ (dans la zone d'ombre $P$) est la \textbf{preuve directe de l'existence d'un noyau interne solide} qui permet la propagation d'ondes $P$ (et $S$ pour $PKJKP$) à des distances où le manteau seul ne les guide pas.
    \end{itemize}
\end{enumerate}
\end{exoresolu}

\begin{attention}[Les 5 erreurs qui coûtent des points au Probatoire]
\begin{enumerate}
    \item \textbf{Confondre "Discontinuité" et "Couche/Enveloppe".} \emph{Exemple :} "Le manteau est une discontinuité à 2900 km." $\rightarrow$ \textbf{Faux.} La discontinuité à 2900 km s'appelle \textbf{Gutenberg} ; elle \textbf{sépare} le manteau du noyau externe. Le manteau est une enveloppe (30 à 2900 km).
    \item \textbf{Oublier l'état physique du noyau externe (Liquide) et interne (Solide).} \emph{Erreur classique :} "Le noyau est liquide." $\rightarrow$ \textbf{Imprécis.} Il faut distinguer : \textbf{Noyau externe = Liquide} ($v_S=0$, pas de cisaillement), \textbf{Noyau interne = Solide} ($v_S>0$, ondes $PKIKP$, $PKJKP$). Justifier par $v_S=0 \Leftrightarrow \mu=0$.
    \item \textbf{Mauvaise application de la loi de Snell-Descartes (réfraction).} \emph{Erreur :} "La vitesse diminue, le rayon s'éloigne de la normale." $\rightarrow$ \textbf{Faux.} $v_2 < v_1 \Rightarrow n_2 > n_1 \Rightarrow$ le rayon se rapproche de la normale (se courbe vers la région de plus basse vitesse). C'est pour cela que les ondes $P$ "plongent" dans le noyau (créant la zone d'ombre).
    \item \textbf{Calcul de profondeur $h$ : utiliser $\Delta T = T_S - T_P$ sans tenir compte de la profondeur du foyer.} \emph{Erreur :} $d = k \times \Delta T$ (formule locale) appliquée à un séisme profond (100 km) ou télésismique. $\rightarrow$ \textbf{Faux.} Les tables de parcours (ou courbes $T(d)$) intègrent la courbure terrestre et la structure en profondeur. Pour un séisme profond, l'onde $P$ part vers le bas, met du temps à remonter $\Rightarrow$ $\Delta T$ plus grand pour une même distance épicentrale. Utiliser la \textbf{table de parcours fournie} ou la \textbf{méthode de Wadati} (diagramme $T_S$ vs $T_P$).
    \item \textbf{Rédaction imprécise sur la composition du noyau.} \emph{Erreur :} "Le noyau est en fer." $\rightarrow$ \textbf{Incomplet.} Préciser : \textbf{Alliage Fer-Ni (Fer majoritaire, $\sim 5-10\%$ Nickel)} + \textbf{éléments légers (Soufre, Oxygène, Silicium, Carbone)} pour expliquer la densité ($10-13$ g/cm$^3$) inférieure au Fe-Ni pur à ces P,T. Pour le noyau interne : \textbf{Fe-Ni quasi pur} (cristallisation par extraction des éléments légers dans le liquide).
\end{enumerate}
\end{attention}

\newpage

\coursec{Exercices}

\courssub{Je vérifie mes connaissances}

\begin{exercice}[QCM : ondes et couches]
Pour chaque question, choisis la (ou les) bonne(s) réponse(s) et justifie brièvement.
\begin{enumerate}
\item Les ondes S :
\begin{enumerate}
\item sont des ondes de compression ;
\item ne se propagent pas dans les liquides ;
\item sont plus rapides que les ondes P ;
\item sont des ondes de cisaillement.
\end{enumerate}
\item L'ordre correct des enveloppes de la Terre, de la surface vers le centre, est :
\begin{enumerate}
\item noyau -- manteau -- croûte ;
\item croûte -- manteau -- noyau ;
\item croûte -- noyau -- manteau ;
\item manteau -- croûte -- noyau.
\end{enumerate}
\item La discontinuité de Mohorovičić (Moho) sépare :
\begin{enumerate}
\item le manteau du noyau ;
\item la croûte du manteau ;
\item le noyau externe de la graine ;
\item la lithosphère de l'asthénosphère.
\end{enumerate}
\item Le noyau externe est :
\begin{enumerate}
\item solide, car les ondes P y accélèrent ;
\item liquide, car les ondes S n'y pénètrent pas ;
\item gazeux, car les ondes P y sont absorbées ;
\item liquide, car les ondes S y sont réfléchies.
\end{enumerate}
\end{enumerate}
\end{exercice}

\begin{exercice}[Vrai ou faux]
Réponds par vrai ou faux en justifiant chaque réponse par une phrase.
\begin{enumerate}
\item Les ondes P se propagent dans les solides et dans les liquides.
\item La vitesse des ondes sismiques est constante dans tout le globe terrestre.
\item La graine (noyau interne) est à l'état solide.
\item Le manteau représente la couche la plus épaisse du globe terrestre.
\item Un séisme ne produit qu'un seul type d'onde sismique.
\end{enumerate}
\end{exercice}

\begin{exercice}[Définitions]
Définis précisément chacun des termes suivants en une ou deux phrases :
\begin{enumerate}
\item foyer (hypocentre) ;
\item épicentre ;
\item discontinuité sismique ;
\item onde P ;
\item sismographe.
\end{enumerate}
\end{exercice}

\begin{exercice}[Calcul direct]
Une onde P se propage dans la croûte continentale à la vitesse moyenne de $v = \SI{6}{km/s}$.
\begin{enumerate}
\item Calcule le temps mis par cette onde pour atteindre une station sismique située à $\SI{150}{km}$ de l'épicentre (on suppose le trajet rectiligne).
\item Une onde S, de vitesse moyenne $\SI{3,5}{km/s}$, part du même foyer au même instant. Calcule son temps de parcours jusqu'à la même station.
\item En déduire l'écart entre les temps d'arrivée des ondes P et S à cette station.
\end{enumerate}
\end{exercice}

\courssub{Je m'entraîne}

\begin{exercice}[L'écho sismique]
Lors d'une prospection sismique, une onde P émise en surface se réfléchit sur une discontinuité et revient au capteur au bout de $t = \SI{20}{s}$. La vitesse moyenne de l'onde dans la couche traversée est $v = \SI{7}{km/s}$.
\begin{enumerate}
\item Rappelle la relation entre distance, vitesse et temps.
\item Calcule la distance totale parcourue par l'onde (aller-retour).
\item En déduire la profondeur de la discontinuité.
\item En comparant cette valeur aux profondeurs caractéristiques des discontinuités du globe, identifie la discontinuité la plus probable.
\end{enumerate}
\end{exercice}

\begin{exercice}[Distance à l'épicentre]
On admet que, pour des distances inférieures à $\SI{1000}{km}$, la distance $d$ (en km) d'une station à l'épicentre est proportionnelle à l'écart $\Delta t$ (en s) entre les arrivées des ondes P et S : $d = 8 \times \Delta t$.
\begin{enumerate}
\item À la station de Yaoundé, l'écart mesuré est $\Delta t = \SI{32}{s}$. Calcule la distance Yaoundé--épicentre.
\item À la station de Douala, l'écart est $\Delta t = \SI{18}{s}$. Calcule la distance Douala--épicentre.
\item Explique, à l'aide d'un schéma, comment trois stations suffisent à localiser l'épicentre (méthode de triangulation).
\end{enumerate}
\end{exercice}

\begin{exercice}[Lecture d'une courbe vitesse--profondeur]
Un document donne les vitesses moyennes des ondes P et S en fonction de la profondeur :
\begin{center}
\begin{tabularx}{\linewidth}{|l|X|X|X|X|}\hline
\rowcolor{popblueL} Profondeur (km) & 30 & 2900 & 5100 & 6370 \\ \hline
Vitesse des ondes P (km/s) & 8 & 13,7 & 8 puis 10,5 & 11,3 \\ \hline
Vitesse des ondes S (km/s) & 4,5 & 7,3 & 0 & 3,7 \\ \hline
\end{tabularx}
\end{center}
\begin{enumerate}
\item Repère les profondeurs où la vitesse des ondes varie brutalement.
\item Nomme les discontinuités correspondantes et les couches qu'elles séparent.
\item À quelle profondeur les ondes S disparaissent-elles ? Qu'en conclus-tu sur l'état physique de la couche située en dessous ?
\item Justifie que la graine est considérée comme solide.
\end{enumerate}
\end{exercice}

\begin{exercice}[Pourquoi un noyau liquide ?]
Les enregistrements sismiques montrent qu'au-delà d'une distance angulaire d'environ $103^{\circ}$ depuis l'épicentre, aucune onde S n'est reçue, et qu'entre $103^{\circ}$ et $142^{\circ}$, les ondes P directes disparaissent aussi (zone d'ombre).
\begin{enumerate}
\item Rappelle la propriété des ondes S vis-à-vis des milieux liquides.
\item Explique comment l'absence des ondes S au-delà de $103^{\circ}$ prouve l'existence d'une enveloppe liquide à l'intérieur du globe.
\item Explique pourquoi la déviation (réfraction) des ondes P à la frontière de cette enveloppe crée une zone d'ombre pour les ondes P.
\item Conclus sur la structure du noyau.
\end{enumerate}
\end{exercice}

\courssub{Je me prépare au Probatoire}

\begin{exercice}[Le séisme de la région de l'Ouest]
Un séisme s'est produit dans la région de l'Ouest du Cameroun. Trois stations sismologiques ont enregistré les heures d'arrivée des ondes P et S (le séisme a eu lieu à 10 h 00 min 00 s) :
\begin{center}
\begin{tabularx}{\linewidth}{|l|X|X|}\hline
\rowcolor{popblueL} Station & Arrivée des ondes P & Arrivée des ondes S \\ \hline
Bafoussam & 10 h 00 min 12 s & 10 h 00 min 21 s \\ \hline
Bamenda & 10 h 00 min 25 s & 10 h 00 min 44 s \\ \hline
Douala & 10 h 00 min 34 s & 10 h 01 min 00 s \\ \hline
\end{tabularx}
\end{center}
On donne : vitesse moyenne des ondes P dans la croûte : $v_P = \SI{6,0}{km/s}$ ; vitesse moyenne des ondes S : $v_S = \SI{3,4}{km/s}$.
\begin{enumerate}
\item Définis les termes foyer et épicentre, puis explique pourquoi les ondes S arrivent après les ondes P.
\item Pour chaque station, calcule le temps de parcours de chaque type d'onde.
\item Sachant que la distance $d$ de la station à l'épicentre vérifie $\Delta t = d\left(\dfrac{1}{v_S} - \dfrac{1}{v_P}\right)$, calcule $d$ pour chacune des trois stations (arrondir au km).
\item Décris la méthode graphique (triangulation) permettant de localiser l'épicentre sur une carte à partir de ces trois distances.
\item Les ondes enregistrées ici se sont propagées dans la croûte continentale, épaisse d'environ $\SI{35}{km}$ sous cette région. Nomme la discontinuité qui limite la croûte en profondeur et précise la nature de la couche située en dessous.
\end{enumerate}
\end{exercice}

\begin{exercice}[De la courbe au modèle : construire la Terre en couches]
Le tableau ci-dessous résume les résultats de la sismologie mondiale :
\begin{center}
\begin{tabularx}{\linewidth}{|l|X|X|}\hline
\rowcolor{popblueL} Couche & Limites (profondeur) & État physique \\ \hline
Croûte (moyenne) & 0 à 30 km & solide \\ \hline
Manteau & 30 à 2900 km & solide (partiellement fondu localement) \\ \hline
Noyau externe & 2900 à 5100 km & liquide \\ \hline
Graine & 5100 à 6370 km & solide \\ \hline
\end{tabularx}
\end{center}
\begin{enumerate}
\item Nomme les deux grandes discontinuités qui séparent ces couches, en précisant la profondeur de chacune.
\item Explique, en t'appuyant sur le comportement des ondes P et S, comment les sismologues ont déterminé l'état physique du noyau externe et celui de la graine.
\item Calcule l'épaisseur de chacune des quatre couches.
\item On veut représenter une coupe du globe à l'échelle $1\ \text{cm} \leftrightarrow 1000\ \text{km}$. Calcule, en cm, le rayon du cercle représentant la Terre et les rayons des cercles limitant chaque couche.
\item Construis soigneusement cette coupe à l'échelle, colorie chaque couche et établis une légende complète (noms, profondeurs, états physiques).
\item Cite une limite de ce modèle simple en couches homogènes.
\end{enumerate}
\end{exercice}

\begin{exercice}[L'héritage d'Inge Lehmann et la structure du globe]
En 1936, la sismologue danoise Inge Lehmann étudie les ondes P enregistrées après de grands séismes. Elle observe que des ondes P arrivent dans la zone d'ombre (entre $103^{\circ}$ et $142^{\circ}$), là où la théorie d'un noyau entièrement liquide prévoyait un silence total. Elle propose alors l'existence d'un noyau interne solide, la graine, qui réfracte une partie des ondes P vers cette zone.
\begin{enumerate}
\item Rappelle les trois grandes enveloppes du globe terrestre, de la surface vers le centre, en précisant pour chacune son épaisseur approximative.
\item Explique pourquoi un noyau entièrement liquide interdit l'arrivée d'ondes S au-delà de $103^{\circ}$ et crée une zone d'ombre pour les ondes P.
\item Montre, par un raisonnement argumenté, comment l'observation d'Inge Lehmann conduit à subdiviser le noyau en deux parties, et précise l'état physique de chacune.
\item Schématise une coupe du globe montrant le trajet : d'une onde S arrêtée à la limite manteau--noyau ; d'une onde P réfractée à l'entrée du noyau externe ; d'une onde P traversant la graine et émergeant dans la zone d'ombre. Légende ton schéma.
\item Conclus en résumant, en quelques lignes, la démarche scientifique qui a permis de construire le modèle actuel de la structure interne de la Terre, alors qu'aucun forage n'a jamais dépassé une quinzaine de kilomètres de profondeur.
\end{enumerate}
\end{exercice}

\newpage

\coursec{Corrigés des exercices}

\begin{corrige}[Exercice 1 — QCM : ondes et couches]
\begin{enumerate}
\item \textbf{Bonnes réponses : b et d.} Les ondes S sont des ondes de \cle{cisaillement} (déplacement de la matière perpendiculairement à la direction de propagation) : elles ne peuvent pas se propager dans les liquides, car un liquide n'oppose aucune résistance au cisaillement. Ce sont les ondes P qui sont des ondes de compression, et ce sont elles les plus rapides (elles arrivent les premières, d'où leur nom : \emph{Primae}).
\item \textbf{Bonne réponse : b.} De la surface vers le centre : \cle{croûte} (0 à \SI{30}{km} en moyenne), \cle{manteau} (jusqu'à \SI{2900}{km}), \cle{noyau} (de \SI{2900}{km} au centre, \SI{6370}{km}).
\item \textbf{Bonne réponse : b.} La discontinuité de \cle{Mohorovičić} (Moho) sépare la croûte du manteau. La limite manteau--noyau est la discontinuité de \cle{Gutenberg} (\SI{2900}{km}), et la limite noyau externe--graine est la discontinuité de \cle{Lehmann} (\SI{5100}{km}).
\item \textbf{Bonne réponse : b.} Le noyau externe est \cle{liquide} : l'argument décisif est l'\textbf{absence totale des ondes S} au-delà de $103^{\circ}$, car les ondes S ne se propagent pas dans un liquide. Attention au piège de la réponse d : les ondes S ne sont pas « réfléchies » par le noyau, elles ne peuvent tout simplement pas s'y propager (elles sont arrêtées à la discontinuité de Gutenberg).
\end{enumerate}
\end{corrige}

\begin{corrige}[Exercice 2 — Vrai ou faux]
\begin{enumerate}
\item \textbf{Vrai.} Les ondes P sont des ondes de compression--dilatation : elles se propagent dans tous les milieux matériels, solides comme liquides (elles traversent notamment le noyau externe liquide).
\item \textbf{Faux.} La vitesse des ondes dépend de la nature et de l'état physique des roches traversées : elle augmente globalement avec la profondeur et varie brutalement aux discontinuités (chute des ondes P à \SI{2900}{km}, disparition des ondes S).
\item \textbf{Vrai.} La graine (noyau interne) est solide : les ondes P y accélèrent et les ondes S, recréées à la frontière noyau externe--graine, s'y propagent, ce qui n'est possible que dans un solide.
\item \textbf{Vrai.} Le manteau s'étend d'environ \SI{30}{km} à \SI{2900}{km}, soit près de \SI{2870}{km} d'épaisseur : c'est de loin la couche la plus épaisse (croûte : environ \SI{30}{km} ; noyau externe : \SI{2200}{km} ; graine : rayon \SI{1270}{km}).
\item \textbf{Faux.} La rupture brutale des roches au foyer émet simultanément plusieurs types d'ondes : ondes de volume (P et S) et ondes de surface (ondes L, responsables des dégâts les plus importants).
\end{enumerate}
\end{corrige}

\begin{corrige}[Exercice 3 — Définitions]
\begin{enumerate}
\item \textbf{Foyer (hypocentre)} : point situé en profondeur, à l'intérieur de la lithosphère, où se produit la rupture brutale des roches à l'origine d'un séisme.
\item \textbf{Épicentre} : point de la surface terrestre situé à la verticale du foyer ; c'est en général là que les secousses sont les plus violentes.
\item \textbf{Discontinuité sismique} : surface de séparation entre deux couches internes du globe, caractérisée par une variation brutale de la vitesse des ondes sismiques (réflexion et/ou réfraction des ondes à ce niveau).
\item \textbf{Onde P} : onde sismique de volume dite de \emph{compression} : la matière vibre parallèlement à la direction de propagation ; c'est l'onde la plus rapide, qui se propage dans les solides et les liquides.
\item \textbf{Sismographe} : appareil qui enregistre en continu les vibrations du sol ; le document obtenu, qui montre l'arrivée successive des différentes ondes, s'appelle un \emph{sismogramme}.
\end{enumerate}
\end{corrige}

\begin{corrige}[Exercice 4 — Calcul direct]
\begin{enumerate}
\item Le temps de parcours est $t = \dfrac{d}{v}$. Pour l'onde P :
\[ t_P = \frac{d}{v_P} = \frac{150}{6} = \SI{25}{s}. \]
\item Pour l'onde S :
\[ t_S = \frac{d}{v_S} = \frac{150}{3{,}5} \approx \SI{42,9}{s} \quad (\text{environ } \SI{43}{s}). \]
\item L'écart entre les temps d'arrivée est :
\[ \Delta t = t_S - t_P = 42{,}9 - 25 = \SI{17,9}{s} \approx \SI{18}{s}. \]
\textbf{Interprétation :} l'onde S, plus lente, arrive après l'onde P ; cet écart $\Delta t$ augmente avec la distance à l'épicentre, ce qui permet de localiser le séisme.
\end{enumerate}
\end{corrige}

\begin{corrige}[Exercice 5 — L'écho sismique]
\begin{enumerate}
\item La relation fondamentale est : $d = v \times t$ (distance = vitesse $\times$ temps).
\item Distance totale parcourue (aller-retour) :
\[ d_{\text{totale}} = v \times t = 7 \times 20 = \SI{140}{km}. \]
\item L'onde a fait un aller-retour jusqu'à la discontinuité ; la profondeur est donc la moitié du trajet :
\[ h = \frac{d_{\text{totale}}}{2} = \frac{140}{2} = \SI{70}{km}. \]
\item Une discontinuité à \SI{70}{km} de profondeur correspond à la \cle{discontinuité de Mohorovičić} (Moho) sous une \textbf{racine de chaîne de montagnes} : la croûte continentale, épaisse de 30 à 40 km en moyenne, peut atteindre 60 à 70 km sous les reliefs élevés. (Sous les océans, le Moho se situe vers 5 à 10 km seulement.)
\end{enumerate}
\textbf{Point de vigilance :} l'erreur classique est d'oublier le facteur 2 : le temps mesuré correspond à un \emph{aller-retour}, la profondeur est donc la moitié de la distance calculée.
\end{corrige}

\begin{corrige}[Exercice 6 — Distance à l'épicentre]
\begin{enumerate}
\item Station de Yaoundé :
\[ d_Y = 8 \times \Delta t = 8 \times 32 = \SI{256}{km}. \]
\item Station de Douala :
\[ d_D = 8 \times 18 = \SI{144}{km}. \]
\item \textbf{Méthode de triangulation :} pour chaque station, on trace sur une carte un cercle centré sur la station et de rayon égal à la distance station--épicentre (à l'échelle de la carte). L'épicentre se trouve sur chacun des cercles : avec deux stations, on obtient deux points d'intersection possibles ; le cercle de la troisième station lève l'ambiguïté. Le point d'intersection commun aux trois cercles est l'épicentre.
\end{enumerate}
\begin{popfigure}
\begin{center}
\begin{tikzpicture}[scale=0.55]
\coordinate (Y) at (0,0);
\coordinate (D) at (6,0);
\coordinate (B) at (3,5);
\coordinate (E) at (3.4,2.2);
\draw[popblue, thick] (Y) circle (2.6);
\draw[poporange, thick] (D) circle (2.9);
\draw[popgreen, thick] (B) circle (3.1);
\fill[popdark] (Y) circle (0.09) node[below left]{\small Yaoundé};
\fill[popdark] (D) circle (0.09) node[below right]{\small Douala};
\fill[popdark] (B) circle (0.09) node[above]{\small Bamenda};
\fill[red] (E) circle (0.13) node[right]{\small \textbf{Épicentre}};
\end{tikzpicture}
\end{center}
\legende{Principe de la triangulation : chaque cercle est centré sur une station et a pour rayon la distance station--épicentre ; l'intersection des trois cercles localise l'épicentre.}
\end{popfigure}
\textbf{Point de vigilance :} une seule station donne une distance, pas une direction (l'épicentre peut être n'importe où sur le cercle) ; deux stations laissent deux solutions ; il faut donc \textbf{au minimum trois stations} pour une localisation unique.
\end{corrige}

\begin{corrige}[Exercice 7 — Lecture d'une courbe vitesse--profondeur]
\begin{enumerate}
\item Les vitesses varient brutalement à trois profondeurs : \SI{30}{km} (accélération des ondes P de 6--7 à \SI{8}{km/s}), \SI{2900}{km} (chute des ondes P de 13,7 à \SI{8}{km/s} et disparition des ondes S) et \SI{5100}{km} (remontée des ondes P de 8 à \SI{10,5}{km/s}).
\item \begin{itemize}
\item À \SI{30}{km} : discontinuité de \cle{Mohorovičić} (Moho), entre la \textbf{croûte} et le \textbf{manteau}.
\item À \SI{2900}{km} : discontinuité de \cle{Gutenberg}, entre le \textbf{manteau} et le \textbf{noyau externe}.
\item À \SI{5100}{km} : discontinuité de \cle{Lehmann}, entre le \textbf{noyau externe} et la \textbf{graine} (noyau interne).
\end{itemize}
\item Les ondes S disparaissent à \SI{2900}{km} (vitesse = 0 dans le tableau). Or les ondes S ne se propagent pas dans les liquides : la couche située en dessous, le \textbf{noyau externe}, est donc à l'\textbf{état liquide}.
\item La graine est considérée comme solide car les ondes P y \textbf{accélèrent} (de 8 à 10,5 puis \SI{11,3}{km/s}) et parce que des ondes S, recréées par conversion à la discontinuité de Lehmann, s'y propagent (vitesse \SI{3,7}{km/s} dans le tableau) : or la propagation d'ondes de cisaillement n'est possible que dans un solide.
\end{enumerate}
\end{corrige}

\begin{corrige}[Exercice 8 — Pourquoi un noyau liquide ?]
\begin{enumerate}
\item Les ondes S sont des ondes de cisaillement : elles \textbf{ne se propagent pas dans les milieux liquides} (un liquide ne transmet pas les contraintes de cisaillement).
\item Les ondes S émises au foyer se propagent dans le manteau solide. Au-delà de $103^{\circ}$, les trajectoires des ondes traverseraient obligatoirement la région centrale du globe. L'absence totale d'ondes S au-delà de $103^{\circ}$ signifie qu'elles ont été \textbf{bloquées par une enveloppe liquide} : il existe donc, sous le manteau, une couche liquide — le noyau externe — qui constitue un écran pour les ondes S.
\item À la frontière manteau--noyau (\SI{2900}{km}), les ondes P passent d'un milieu où elles sont rapides (environ \SI{13,7}{km/s}) à un milieu liquide où elles sont plus lentes (environ \SI{8}{km/s}) : elles sont alors \textbf{réfractées}, c'est-à-dire déviées vers le centre. Cette déviation systématique « détourne » les ondes P, si bien qu'aucune onde P directe n'émerge entre $103^{\circ}$ et $142^{\circ}$ : c'est la \cle{zone d'ombre} des ondes P.
\item \textbf{Conclusion :} la Terre possède un \textbf{noyau} situé sous \SI{2900}{km} de profondeur, dont la partie externe est \textbf{liquide} (blocage des S, réfraction des P). L'observation d'ondes P dans la zone d'ombre a ensuite montré que ce noyau contient une partie centrale solide, la graine.
\end{enumerate}
\end{corrige}

\begin{corrige}[Exercice 9 — Le séisme de la région de l'Ouest]
\begin{enumerate}
\item Le \textbf{foyer} est le point en profondeur où se produit la rupture des roches ; l'\textbf{épicentre} est le point de la surface situé à la verticale du foyer. Les ondes S arrivent après les ondes P parce qu'elles sont \textbf{plus lentes} ($v_S = \SI{3,4}{km/s} < v_P = \SI{6,0}{km/s}$) : à distance égale, leur temps de parcours est plus long.
\item Temps de parcours (heure d'arrivée $-$ 10 h 00 min 00 s) :
\begin{center}
\begin{tabularx}{\linewidth}{|l|X|X|X|}\hline
\rowcolor{popblueL} Station & $t_P$ (s) & $t_S$ (s) & $\Delta t = t_S - t_P$ (s) \\ \hline
Bafoussam & 12 & 21 & 9 \\ \hline
Bamenda & 25 & 44 & 19 \\ \hline
Douala & 34 & 60 & 26 \\ \hline
\end{tabularx}
\end{center}
\item On isole $d$ dans la relation donnée :
\[ \Delta t = d\left(\frac{1}{v_S} - \frac{1}{v_P}\right) \quad\Longrightarrow\quad d = \frac{\Delta t}{\dfrac{1}{v_S} - \dfrac{1}{v_P}}. \]
Calcul du coefficient :
\[ \frac{1}{v_S} - \frac{1}{v_P} = \frac{1}{3{,}4} - \frac{1}{6{,}0} = 0{,}2941 - 0{,}1667 = \SI{0,1275}{s/km}. \]
D'où :
\begin{align*}
d_{\text{Bafoussam}} &= \frac{9}{0{,}1275} \approx \SI{71}{km},\\
d_{\text{Bamenda}} &= \frac{19}{0{,}1275} \approx \SI{149}{km},\\
d_{\text{Douala}} &= \frac{26}{0{,}1275} \approx \SI{204}{km}.
\end{align*}
\item \textbf{Triangulation :} sur une carte de la région, on trace trois cercles centrés respectivement sur Bafoussam, Bamenda et Douala, de rayons 71, 149 et \SI{204}{km} (convertis à l'échelle de la carte). Le point d'intersection commun des trois cercles est l'épicentre du séisme.
\item La discontinuité qui limite la croûte en profondeur est la discontinuité de \cle{Mohorovičić} (Moho). En dessous se trouve le \textbf{manteau supérieur}, constitué de roches solides (péridotites), plus denses que celles de la croûte.
\end{enumerate}
\textbf{Barème indicatif (sur 10) :} définitions + justification ondes S : 2 pts ; temps de parcours : 1,5 pt ; expression littérale de $d$ : 1 pt ; trois distances calculées : 3 pts ; triangulation décrite : 1,5 pt ; Moho + manteau : 1 pt.

\textbf{Points de vigilance :} convertir les heures d'arrivée en \emph{durées} (soustraction avec l'heure du séisme) avant tout calcul ; écrire l'expression littérale de $d$ avant l'application numérique ; vérifier la cohérence : plus $\Delta t$ est grand, plus la station est loin (Douala $>$ Bamenda $>$ Bafoussam) ; ne pas confondre distance station--épicentre et profondeur du foyer.
\end{corrige}

\begin{corrige}[Exercice 10 — De la courbe au modèle : construire la Terre en couches]
\begin{enumerate}
\item Les deux grandes discontinuités sont :
\begin{itemize}
\item la discontinuité de \cle{Mohorovičić} (Moho), vers \SI{30}{km}, entre croûte et manteau ;
\item la discontinuité de \cle{Gutenberg}, à \SI{2900}{km}, entre manteau et noyau.
\end{itemize}
(On peut ajouter la discontinuité de \cle{Lehmann}, à \SI{5100}{km}, entre noyau externe et graine.)
\item \textbf{Noyau externe liquide :} les ondes S disparaissent brutalement à \SI{2900}{km} et ne sont plus jamais enregistrées au-delà ; or les ondes S ne se propagent pas dans les liquides. De plus, les ondes P y ralentissent fortement et sont réfractées, créant une zone d'ombre. \textbf{Graine solide :} à \SI{5100}{km}, les ondes P accélèrent à nouveau, et des ondes P arrivent dans la zone d'ombre après avoir été réfractées par une sphère centrale (observation d'Inge Lehmann, 1936) ; des ondes S recréées à cette frontière s'y propagent, ce qui exige un solide.
\item Épaisseurs des couches :
\begin{align*}
\text{Croûte : } & 30 - 0 = \SI{30}{km},\\
\text{Manteau : } & 2900 - 30 = \SI{2870}{km},\\
\text{Noyau externe : } & 5100 - 2900 = \SI{2200}{km},\\
\text{Graine (rayon) : } & 6370 - 5100 = \SI{1270}{km}.
\end{align*}
\item À l'échelle $1\ \text{cm} \leftrightarrow \SI{1000}{km}$, les \textbf{rayons} des cercles limites (comptés depuis le centre) sont :
\begin{align*}
\text{Surface terrestre : } & 6370\ \text{km} \rightarrow \SI{6,37}{cm},\\
\text{Moho : } & 6370 - 30 = 6340\ \text{km} \rightarrow \SI{6,34}{cm},\\
\text{Gutenberg : } & 6370 - 2900 = 3470\ \text{km} \rightarrow \SI{3,47}{cm},\\
\text{Lehmann : } & 6370 - 5100 = 1270\ \text{km} \rightarrow \SI{1,27}{cm}.
\end{align*}
On remarque que la croûte (0,03 cm d'épaisseur sur le dessin) est quasi impossible à représenter à cette échelle : c'est une pellicule.
\item Construction : on trace quatre cercles concentriques de rayons 6,37 ; 6,34 ; 3,47 et \SI{1,27}{cm}, puis on colorie : croûte (mince liseré), manteau, noyau externe, graine.
\end{enumerate}
\begin{popfigure}
\begin{center}
\begin{tikzpicture}[scale=0.85]
% Remplissage des couches (du plus grand rayon au plus petit)
\fill[popblueL] (0,0) circle (6.34);   % Manteau (+ base croûte)
\fill[poporangeL] (0,0) circle (3.47); % Noyau externe
\fill[popgoldL] (0,0) circle (1.27);   % Graine

% Contours des discontinuités et surface
\draw[popdark, thick] (0,0) circle (6.37); % Surface Terre
\draw[popdark] (0,0) circle (6.34);        % Moho
\draw[popdark] (0,0) circle (3.47);        % Gutenberg
\draw[popdark] (0,0) circle (1.27);        % Lehmann

% Annotations
\draw[thin, popmuted] (6.34,0.2) -- (7.8,1.0) node[right]{\small Croûte (0--30 km, solide)};
\draw[thin, popmuted] (3.0,2.5) -- (7.8,2.8) node[right]{\small Manteau (30--2900 km, solide)};
\draw[thin, popmuted] (-2.5,-2.5) -- (-7.8,-3.0) node[left]{\small Noyau externe (2900--5100 km, liquide)};
\draw[thin, popmuted] (0.8,-1.0) -- (7.8,-2.0) node[right]{\small Graine (5100--6370 km, solide)};
\end{tikzpicture}
\end{center}
\legende{Coupe du globe à l'échelle $1\ \text{cm} \leftrightarrow \SI{1000}{km}$ : quatre enveloppes concentriques séparées par les discontinuités de Mohorovičić, Gutenberg et Lehmann.}
\end{popfigure}
\begin{enumerate}
\item[6.] \textbf{Limite du modèle :} ce modèle suppose des couches \textbf{homogènes} aux limites nettes, alors qu'en réalité la croûte est hétérogène (continentale/océanique, d'épaisseur variable), le manteau comporte des zones partiellement fondues (asthénosphère), et les discontinuités ne sont pas des surfaces parfaitement régulières.
\end{enumerate}
\textbf{Barème indicatif (sur 10) :} discontinuités nommées avec profondeurs : 2 pts ; arguments ondes P/S pour les états physiques : 2,5 pts ; épaisseurs : 1,5 pt ; calculs d'échelle : 2 pts ; construction soignée et légendée : 1,5 pt ; limite du modèle : 0,5 pt.

\textbf{Points de vigilance :} pour le dessin à l'échelle, ce sont les \textbf{rayons depuis le centre} qu'il faut convertir, pas les épaisseurs ; justifier chaque état physique par le comportement des ondes (et non l'affirmer gratuitement) ; citer les noms des discontinuités avec leur profondeur.
\end{corrige}

\begin{corrige}[Exercice 11 — L'héritage d'Inge Lehmann et la structure du globe]
\begin{enumerate}
\item De la surface vers le centre :
\begin{itemize}
\item la \textbf{croûte} : épaisseur moyenne d'environ \SI{30}{km} (5 à 10 km sous les océans, 30 à 70 km sous les continents) ;
\item le \textbf{manteau} : d'environ \SI{30}{km} à \SI{2900}{km}, soit près de \SI{2870}{km} d'épaisseur ;
\item le \textbf{noyau} : de \SI{2900}{km} au centre (\SI{6370}{km}), soit un rayon d'environ \SI{3470}{km}.
\end{itemize}
\item Un noyau entièrement liquide \textbf{bloque les ondes S}, car les ondes de cisaillement ne se propagent pas dans un liquide : aucune onde S ne peut donc être reçue au-delà de $103^{\circ}$, distance à partir de laquelle les trajectoires traversent le noyau. Par ailleurs, les ondes P, en entrant dans le liquide, \textbf{ralentissent et sont réfractées vers le centre} ; cette déviation systématique les empêche d'émerger entre $103^{\circ}$ et $142^{\circ}$ : d'où la zone d'ombre des ondes P.
\item \textbf{Raisonnement d'Inge Lehmann :} si le noyau était entièrement liquide, la zone d'ombre devrait être totalement silencieuse pour les ondes P directes. Or elle enregistre des ondes P \emph{dans} cette zone. Ces ondes ont donc suivi un trajet différent : elles ont été \textbf{réfractées par une sphère centrale plus rigide}, qui redévie les rayons vers la zone d'ombre. L'existence de cette sphère impose de subdiviser le noyau en : \textbf{noyau externe liquide} (2900 à \SI{5100}{km}) et \textbf{graine solide} (de \SI{5100}{km} au centre). L'accélération des ondes P à \SI{5100}{km} confirme ce modèle.
\item Schéma des trajets d'ondes :
\end{enumerate}
\begin{popfigure}
\begin{center}
\begin{tikzpicture}[scale=0.62]
% Couches
\fill[popblueL] (0,0) circle (5.5);
\fill[poporangeL] (0,0) circle (3.0);
\fill[popgoldL] (0,0) circle (1.4);
% Contours
\draw[popdark, thick] (0,0) circle (5.5);
\draw[popdark] (0,0) circle (3.0);
\draw[popdark] (0,0) circle (1.4);
% Foyer
\coordinate (F) at (-4.6,3.0);
\fill[red] (F) circle (0.14) node[above left]{\small Foyer};
% Onde S (verte) - bloquée à Gutenberg (r=3.0)
\draw[very thick, popgreen] (F) -- (-2.1,-2.1);
\draw[very thick, popgreen, dashed] (-2.1,-2.1) -- (-0.5,-2.9);
% Onde P réfractée noyau externe (bleue)
\draw[very thick, popblue] (F) -- (1.2,-2.75);
\draw[very thick, popblue] (1.2,-2.75) -- (3.4,-1.6);
\draw[very thick, popblue] (3.4,-1.6) -- (5.3,-1.4);
% Onde P traversant graine (violette) - émerge dans zone d'ombre
\draw[very thick, poppurple] (F) -- (-0.9,-1.05);
\draw[very thick, poppurple] (-0.9,-1.05) -- (1.05,0.9);
\draw[very thick, poppurple] (1.05,0.9) -- (2.6,1.5);
\draw[very thick, poppurple] (2.6,1.5) -- (4.9,2.5);
% Points d'émergence
\fill[popdark] (5.3,-1.4) circle (0.09);
\fill[popdark] (4.9,2.5) circle (0.09);
% Légendes sur figure
\draw[thin, popmuted] (-2.1,-2.1) -- (-4.4,-4.3) node[below]{\small Onde S arrêtée à la limite manteau--noyau};
\draw[thin, popmuted] (3.4,-1.6) -- (6.6,-3.4) node[below]{\small Onde P réfractée à l'entrée du noyau externe};
\draw[thin, popmuted] (1.05,0.9) -- (6.6,4.2) node[above]{\small Onde P traversant la graine};
\draw[thin, popmuted] (4.9,2.5) -- (6.6,2.9) node[right]{\small Émergence dans la zone d'ombre};
\end{tikzpicture}
\end{center}
\legende{Trajets des ondes dans le globe : l'onde S (vert) est bloquée à la discontinuité de Gutenberg ; l'onde P (bleu) est réfractée à l'entrée du noyau externe ; l'onde P (violet) réfractée par la graine émerge dans la zone d'ombre, observation décisive d'Inge Lehmann.}
\end{popfigure}
\begin{enumerate}
\item[5.] \textbf{Conclusion — la démarche scientifique :} la structure interne de la Terre n'a jamais été observée directement (le forage le plus profond, Kola, n'atteint que \SI{12}{km}). Elle a été reconstituée par une \textbf{méthode indirecte} : l'enregistrement des ondes sismiques par des réseaux de stations répartis sur le globe, l'analyse de leurs temps de parcours, de leurs réflexions, réfractions et disparitions, puis la construction d'un \textbf{modèle} cohérent avec l'ensemble des données. Ce modèle a été affiné au fil des observations (Moho 1909, Gutenberg 1913, Lehmann 1936) : c'est une illustration remarquable de la démarche scientifique, où un modèle est testé, corrigé et enrichi à chaque donnée nouvelle.
\end{enumerate}
\textbf{Barème indicatif (sur 10) :} enveloppes et épaisseurs : 2 pts ; explication zone d'ombre (S bloquées + P réfractées) : 2,5 pts ; raisonnement de Lehmann et subdivision du noyau : 2 pts ; schéma correct et légendé : 2 pts ; conclusion sur la démarche : 1,5 pt.

\textbf{Points de vigilance :} distinguer clairement les deux arguments (blocage des S $\neq$ réfraction des P) ; le schéma doit montrer la \emph{déviation} des rayons aux discontinuités, pas des lignes droites ; citer les noms et dates des scientifiques valorise la copie ; la conclusion doit explicitement mentionner le caractère \emph{indirect} et \emph{modélisateur} de la sismologie.
\end{corrige}

\newpage

\coursec{Fiche bilan}

\begin{popfigure}
\centering
\begin{center}\tcbox[colback=popblue,colframe=popblue,coltext=white,arc=9pt,boxsep=4pt,left=6pt,right=6pt]{\bfseries\small Structure interne de la Terre}\end{center}
\vspace{-2pt}
\begin{tcbraster}[raster columns=3,raster equal height=rows,raster column skip=6pt,raster row skip=6pt,enhanced,blankest]
\begin{tcolorbox}[enhanced,colback=poporange,colframe=poporange,coltext=white,boxrule=0.9pt,arc=3pt,left=4pt,right=4pt,top=3pt,bottom=3pt,boxsep=1pt,halign=left]\bfseries\scriptsize Séismes et ondes sismiques P (longitudinales) et S (transversales)\end{tcolorbox}
\begin{tcolorbox}[enhanced,colback=popgreen,colframe=popgreen,coltext=white,boxrule=0.9pt,arc=3pt,left=4pt,right=4pt,top=3pt,bottom=3pt,boxsep=1pt,halign=left]\bfseries\scriptsize Comportement aux discontinuités réflexion, réfraction, arrêt des S\end{tcolorbox}
\begin{tcolorbox}[enhanced,colback=poppurple,colframe=poppurple,coltext=white,boxrule=0.9pt,arc=3pt,left=4pt,right=4pt,top=3pt,bottom=3pt,boxsep=1pt,halign=left]\bfseries\scriptsize Sismologie temps de parcours $\rightarrow$ profondeur\end{tcolorbox}
\begin{tcolorbox}[enhanced,colback=popgold,colframe=popgold,coltext=white,boxrule=0.9pt,arc=3pt,left=4pt,right=4pt,top=3pt,bottom=3pt,boxsep=1pt,halign=left]\bfseries\scriptsize Discontinuités Mohorovičić, Gutenberg, Lehmann\end{tcolorbox}
\begin{tcolorbox}[enhanced,colback=poppink,colframe=poppink,coltext=white,boxrule=0.9pt,arc=3pt,left=4pt,right=4pt,top=3pt,bottom=3pt,boxsep=1pt,halign=left]\bfseries\scriptsize Enveloppes croûte, manteau, noyau\end{tcolorbox}
\begin{tcolorbox}[enhanced,colback=popblue!60,colframe=popblue!60,coltext=white,boxrule=0.9pt,arc=3pt,left=4pt,right=4pt,top=3pt,bottom=3pt,boxsep=1pt,halign=left]\bfseries\scriptsize Modèle en couches et limites du modèle\end{tcolorbox}
\end{tcbraster}

\legende{Carte mentale de la leçon : de l'observation des séismes à la construction du modèle en couches du globe terrestre.}
\end{popfigure}

\begin{bilan}
\textbf{Définitions clés}
\begin{itemize}
  \item \cle{Séisme} : vibration brutale du sol provoquée par une rupture de roches en profondeur (foyer ou hypocentre) ; le point de la surface à la verticale du foyer est l'\cle{épicentre}.
  \item \cle{Ondes P} (primaires) : ondes longitudinales de compression-dilatation ; elles se propagent dans les solides \emph{et} les liquides ; ce sont les plus rapides (environ 6 à \SI{13}{km/s}).
  \item \cle{Ondes S} (secondaires) : ondes transversales de cisaillement ; elles ne se propagent \textbf{que dans les solides} ; vitesse plus faible que les P.
  \item \cle{Discontinuité} : surface de séparation entre deux enveloppes où la vitesse des ondes change brutalement (changement de nature ou d'état des matériaux).
\end{itemize}

\textbf{Repères à connaître par c\oe ur}
\begin{center}
\begin{tabularx}{\linewidth}{|l|X|X|}\hline
\rowcolor{popblueL} \textbf{Élément} & \textbf{Profondeur / épaisseur} & \textbf{Particularité} \\ \hline
Discontinuité de Mohorovičić (Moho) & $\approx$ 5 à \SI{70}{km} (moy. $\approx$ \SI{30}{km}) & limite croûte / manteau \\ \hline
Croûte (écorce terrestre) & mince enveloppe superficielle & solide ; croûte continentale plus épaisse qu'océanique \\ \hline
Manteau & jusqu'à $\approx$ \SI{2900}{km} & solide (les ondes S le traversent) \\ \hline
Discontinuité de Gutenberg & $\approx$ \SI{2900}{km} & limite manteau / noyau ; les ondes S disparaissent \\ \hline
Noyau externe & de \SI{2900}{km} à $\approx$ \SI{5100}{km} & \textbf{liquide} (ombre aux ondes S) \\ \hline
Discontinuité de Lehmann & $\approx$ \SI{5100}{km} & limite noyau externe / graine \\ \hline
Graine (noyau interne) & de \SI{5100}{km} à \SI{6371}{km} & solide \\ \hline
\end{tabularx}
\end{center}

\textbf{Méthodes express}
\begin{enumerate}
  \item \textbf{Localiser une discontinuité} : repérer sur un graphique temps-distance (ou vitesse-profondeur) une rupture de pente ou une zone d'ombre $\rightarrow$ convertir en profondeur avec $v = d/t$.
  \item \textbf{Déduire l'état physique} : si les ondes S passent $\rightarrow$ solide ; si les S sont arrêtées et les P ralenties $\rightarrow$ liquide.
  \item \textbf{Construire le modèle} : à chaque discontinuité correspond une couche ; dessiner des couches concentriques et légender (nom, profondeur, état).
\end{enumerate}
\end{bilan}

\begin{aretenir}[Checklist avant le Probatoire]
\begin{itemize}
  \item[$\square$] Je sais définir séisme, foyer, épicentre.
  \item[$\square$] Je distingue ondes P (longitudinales, rapides, milieux solides et liquides) et ondes S (transversales, solides uniquement).
  \item[$\square$] Je sais que les ondes se réfléchissent et se réfractent aux discontinuités.
  \item[$\square$] Je sais exploiter un document de temps de parcours pour situer une discontinuité.
  \item[$\square$] Je connais les trois enveloppes : croûte, manteau, noyau (avec graine).
  \item[$\square$] Je connais les profondeurs repères : Moho ($\approx$ \SI{30}{km}), Gutenberg (\SI{2900}{km}), Lehmann (\SI{5100}{km}).
  \item[$\square$] Je sais que le noyau externe est liquide car les ondes S n'y passent pas (zone d'ombre).
  \item[$\square$] Je sais que le manteau est solide car les ondes S le traversent.
  \item[$\square$] Je sais relier la vitesse des ondes aux propriétés physiques des couches traversées.
  \item[$\square$] Je sais schématiser et légender une coupe du globe en couches concentriques.
  \item[$\square$] Je sais citer les limites du modèle (modèle simplifié, hétérogénéités locales).
  \item[$\square$] Je rédige mes réponses avec le vocabulaire scientifique exact (discontinuité, réfraction, enveloppe).
\end{itemize}
\end{aretenir}

\findecours

\end{document}
